\input hofstadter.opm

Douglas Richard Hofstadter (born February 15, 1945 -) is an American scholar of cognitive science, physics, and comparative literature whose research includes concepts such as the sense of self in relation to the external world, consciousness, analogy-making, artistic creation, literary translation, and discovery in mathematics and physics. His 1979 book Gödel, Escher, Bach: An Eternal Golden Braid won both the Pulitzer Prize for general nonfiction and a National Book Award (at that time called The American Book Award) for Science. His 2007 book I Am a Strange Loop won the Los Angeles Times Book Prize for Science and Technology. Hofstadter was born in New York City to Jewish parents: Nobel Prize- winning physicist Robert Hofstadter and Nancy Givan Hofstadter. He grew up on the campus of Stanford University, where his father was a professor, and attended the International School of Geneva in 1958–59. He graduated with Distinction in mathematics from Stanford University in 1965, and received his Ph.D. in physics from the University of Oregon in 1975, where his study of the energy levels of Bloch electrons in a magnetic field led to his discovery of the fractal known as Hofstadter's butterfly. Contents Overview                                                      viii List of Illustrations                                         xiv Words of Thanks                                               xix 

Part I: GEB 
\chapter {IV}








Contents                                                             II Part II EGB 


\chapter {XIV}












Contents                                                         III Overview Part I: GEB Introduction: A Musico-Logical Offering. The book opens with the story of Bach's Musical Offering. Bach made an impromptu visit to King Frederick the Great of Prussia, and was requested to improvise upon a theme presented by the King. His improvisations formed the basis of that great work. The Musical Offering and its story form a theme upon which I "improvise" throughout the book, thus making a sort of "Metamusical Offering". Self-reference and the interplay between different levels in Bach are discussed: this leads to a discussion of parallel ideas in Escher's drawings and then Gödel’s Theorem. A brief presentation of the history of logic and paradoxes is given as background for Gödel’s Theorem. This leads to mechanical reasoning and computers, and the debate about whether Artificial Intelligence is possible. I close with an explanation of the origins of the book-particularly the why and wherefore of the Dialogues. 

Three-Part Invention. Bach wrote fifteen three-part inventions. In this three-part Dialogue, the Tortoise and Achilles-the main fictional protagonists in the Dialogues-are "invented" by Zeno (as in fact they were, to illustrate Zeno's paradoxes of motion). Very short, it simply gives the flavor of the Dialogues to come. 


\chapter {I}
is urged to work out a puzzle to gain familiarity with formal systems in general. A number of fundamental notions are introduced: string, theorem, axiom, rule of inference, derivation, formal system, decision procedure, working inside/outside the system. 

Two-Part Invention. Bach also wrote fifteen two-part inventions. This two-part Dialogue was written not by me, but by Lewis Carroll in 1895. Carroll borrowed Achilles and the Tortoise from Zeno, and I in turn borrowed them from Carroll. The topic is the relation between reasoning, reasoning about reasoning, reasoning about reasoning about reasoning, and so on. It parallels, in a way, Zeno's paradoxes about the impossibility of motion, seeming to show, by using infinite regress, that reasoning is impossible. It is a beautiful paradox, and is referred to several times later in the book. 


\chapter {II}
presented, even simpler than the MIU-system of Chapter I. Apparently meaningless at first, its symbols are suddenly revealed to possess meaning by virtue of the form of the theorems they appear in. This revelation is the first important insight into meaning: its deep connection to isomorphism. Various issues related to meaning are then discussed, such as truth, proof, symbol manipulation, and the elusive concept, "form". 

Sonata for Unaccompanied Achilles. A Dialogue which imitates the Bach Sonatas for unaccompanied violin. In particular, Achilles is the only speaker, since it is a transcript of one end of a telephone call, at the far end of which is the Tortoise. Their conversation concerns the concepts of "figure" and "ground" in various 







Overview                                                                                          IV contexts- e.g., Escher's art. The Dialogue itself forms an example of the distinction, since Achilles' lines form a "figure", and the Tortoise's lines-implicit in Achilles' lines-form a "ground". 


\chapter {III}
the distinction between theorems and nontheorems in formal systems. The question "Does a figure necessarily contain the same information as its ground%" leads to the distinction between recursively enumerable sets and recursive sets. 

Contracrostipunctus. This Dialogue is central to the book, for it contains a set of paraphrases of Gödel’s self-referential construction and of his Incompleteness Theorem. One of the paraphrases of the Theorem says, "For each record player there is a record which it cannot play." The Dialogue's title is a cross between the word "acrostic" and the word "contrapunctus", a Latin word which Bach used to denote the many fugues and canons making up his Art of the Fugue. Some explicit references to the Art of the Fugue are made. The Dialogue itself conceals some acrostic tricks. 


\chapter {IV}
the extent it is possible at this stage. This leads back to the question of how and when symbols in a formal system acquire meaning. The history of Euclidean and non-Euclidean geometry is given, as an illustration of the elusive notion of "undefined terms". This leads to ideas about the consistency of different and possibly "rival" geometries. Through this discussion the notion of undefined terms is clarified, and the relation of undefined terms to perception and thought processes is considered. 

Little Harmonic Labyrinth. This is based on the Bach organ piece by the same name. It is a playful introduction to the notion of recursive-i.e., nested structures. It contains stories within stories. The frame story, instead of finishing as expected, is left open, so the reader is left dangling without resolution. One nested story concerns modulation in music-particularly an organ piece which ends in the wrong key, leaving the listener dangling without resolution. 


\chapter {V}
different contexts: musical patterns, linguistic patterns, geometric structures, mathematical functions, physical theories, computer programs, and others. 

Canon by Intervallic Augmentation. Achilles and the Tortoise try to resolve the question, "Which contains more information-a record, or the phonograph which plays it This odd question arises when the Tortoise describes a single record which, when played on a set of different phonographs, produces two quite different melodies: B-A-C-H and C-A-G-E. It turns out, however, that these melodies are "the same", in a peculiar sense. 


\chapter {VI}
message, decoder, and receiver. Examples presented include strands of DNA, undeciphered inscriptions on ancient tablets, and phonograph records sailing out in space. The relationship of intelligence to "absolute" meaning is postulated. 

Chromatic Fantasy, And Feud. A short Dialogue bearing hardly any resemblance, except in title, to Bach's Chromatic Fantasy and Fugue. It concerns the proper way to manipulate sentences so as to preserve truth-and in particular the question 







Overview                                                                                              V of whether there exist rules for the usage of the word "arid". This Dialogue has much in common with the Dialogue by Lewis Carroll. 


\chapter {VII}
governed by formal rules. Once again, the ideas of isomorphism and automatic acquisition of meaning by symbols in such a system are brought up. All the examples in this Chapter, incidentally, are "Zentences"-sentences taken from Zen koans. This is purposefully done, somewhat tongue-in-cheek, since Zen koans are deliberately illogical stories. 

Crab Canon. A Dialogue based on a piece by the same name from the Musical Offering. Both are so named because crabs (supposedly) walk backwards. The Crab makes his first appearance in this Dialogue. It is perhaps the densest Dialogue in the book in terms of formal trickery and level- play. Gödel, Escher, and Bach are deeply intertwined in this very short Dialogue. 


\chapter {VIII}
"TNT" is presented. In TNT, number-theoretical reasoning can be done by rigid symbol manipulation. Differences between formal reasoning and human thought are considered. 

A Mu Offering. This Dialogue foreshadows several new topics in the book. Ostensibly concerned with Zen Buddhism and koans, it is actually a thinly veiled discussion of theoremhood and nontheoremhood, truth and falsity, of strings in number theory. There are fleeting references to molecular biology-particular) the Genetic Code. There is no close affinity to the Musical Offering, other than in the title and the playing of self-referential games. 


\chapter {IX}
Buddhism. The Zen monk Mumon, who gave well known commentaries on many koans, is a central figure. In a way, Zen ideas bear a metaphorical resemblance to some contemporary ideas in the philosophy of mathematics. After this "Zennery", Gödel’s fundamental idea of Gödel- numbering is introduced, and a first pass through Gödel’s Theorem is made. 

Part II: EGB Prelude ... This Dialogue attaches to the next one. They are based on preludes and fugues from Bach's Well-Tempered Clavier. Achilles and the Tortoise bring a present to the Crab, who has a guest: the Anteater. The present turns out to be a recording of the W.T.C.; it is immediately put on. As they listen to a prelude, they discuss the structure of preludes and fugues, which leads Achilles to ask how to hear a fugue: as a whole, or as a sum of parts? This is the debate between holism and reductionism, which is soon taken up in the Ant Fugue. 


\chapter {X}
chessboards, and computer systems are discussed. The last of these is then examined in detail. This involves describing machine languages, assembly languages, compiler languages, operating systems, and so forth. Then the discussion turns to composite systems of other types, such as sports teams, nuclei, atoms, the weather, and so forth. The question arises as to how man intermediate levels exist-or indeed whether any exist. 







Overview                                                                                      VI …Ant Fugue. An imitation of a musical fugue: each voice enters with the same statement. The theme-holism versus reductionism-is introduced in a recursive picture composed of words composed of smaller words. etc. The words which appear on the four levels of this strange picture are "HOLISM", "REDLCTIONIsM", and "ML". The discussion veers off to a friend of the Anteater's Aunt Hillary, a conscious ant colony. The various levels of her thought processes are the topic of discussion. Many fugal tricks are ensconced in the Dialogue. As a hint to the reader, references are made to parallel tricks occurring in the fugue on the record to which the foursome is listening. At the end of the Ant Fugue, themes from the Prelude return. transformed considerably. 


\chapter {XI}
the topic of the Chapter. An overview of the large scale and small-scale structure of the brain is first given. Then the relation between concepts and neural activity is speculatively discussed in some detail. 

English French German Suite. An interlude consisting of Lewis Carroll's nonsense poem "Jabberwocky`' together with two translations: one into French and one into German, both done last century. 


\chapter {XII}
of whether languages, or indeed minds, can be "mapped" onto each other. How is communication possible between two separate physical brains: What do all human brains have in common? A geographical analogy is used to suggest an answer. The question arises, "Can a brain be understood, in some objective sense, by an outsider?" 

Aria with Diverse Variations. A Dialogue whose form is based on Bach's Goldberg Variations, and whose content is related to number-theoretical problems such as the Goldbach conjecture. This hybrid has as its main purpose to show how number theory's subtlety stems from the fact that there are many diverse variations on the theme of searching through an infinite space. Some of them lead to infinite searches, some of them lead to finite searches, while some others hover in between. 


\chapter {XIII}
BlooP programs can carry out only predictably finite searches, while FlooP programs can carry out unpredictable or even infinite searches. The purpose of this Chapter is to give an intuition for the notions of primitive recursive and general recursive functions in number theory, for they are essential in Gödel’s proof. 

Air on G's String. A Dialogue in which Gödel’s self-referential construction is mirrored in words. The idea is due to W. V. O. Quine. This Dialogue serves as a prototype for the next Chapter. 


\chapter {XIV}
Chapter's title is an adaptation of the title of Gödel’s 1931 article, in which his Incompleteness Theorem was first published. The two major parts of Gödel’s proof are gone through carefully. It is shown how the assumption of consistency of TNT forces one to conclude that TNT (or any similar system) is incomplete. Relations to Euclidean and non-Euclidean geometry are discussed. Implications for the philosophy of mathematics are gone into with some care. 







Overview                                                                                       VII Birthday Cantatatata ... In which Achilles cannot convince the wily and skeptical Tortoise that today is his (Achilles') birthday. His repeated but unsuccessful tries to do so foreshadow the repeatability of the Gödel argument. 


\chapter {XV}
the implication that TNT is not only incomplete, but "essentially incomplete The fairly notorious argument by J. R. Lucas, to the effect that Gödel’s Theorem demonstrates that human thought cannot in any sense be "mechanical", is analyzed and found to be wanting. 

Edifying Thoughts of a Tobacco Smoker. A Dialogue treating of many topics, with the thrust being problems connected with self-replication and self-reference. Television cameras filming television screens, and viruses and other subcellular entities which assemble themselves, are among the examples used. The title comes from a poem by J. S. Bach himself, which enters in a peculiar way. 


\chapter {XVI}
in its various guises, and self-reproducing entities e.g., computer programs or DNA molecules). The relations between a self-reproducing entity and the mechanisms external to it which aid it in reproducing itself (e.g., a computer or proteins) are discussed-particularly the fuzziness of the distinction. How information travels between various levels of such systems is the central topic of this Chapter. 

The Magnificrab, Indeed. The title is a pun on Bach's Magnifacat in D. The tale is about the Crab, who gives the appearance of having a magical power of distinguishing between true and false statements of number theory by reading them as musical pieces, playing them on his flute, and determining whether they are "beautiful" or not. 


\chapter {XVII}
is replaced by various real people with amazing mathematical abilities. The Church-Turing Thesis, which relates mental activity to computation, is presented in several versions of differing strengths. All are analyzed, particularly in terms of their implications for simulating human thought mechanically, or programming into a machine an ability to sense or create beauty. The connection between brain activity and computation brings up some other topics: the halting problem of Turing, and Tarski's Truth Theorem. 

SHRDLU, Toy of Man's Designing. This Dialogue is lifted out of an article by Terry Winograd on his program SHRDLU: only a few names have been changed. In it. a program communicates with a person about the so-called "blocks world" in rather impressive English. The computer program appears to exhibit some real understanding-in its limited world. The Dialogue's title is based on Jesu, joy of Mans Desiring, one movement of Bach's Cantata 147. 


\chapter {XVIII}
famous "Turing test"-a proposal by the computer pioneer Alan Turing for a way to detect the presence or absence of "thought" in a machine. From there, we go on to an abridged history of Artificial Intelligence. This covers programs that can-to some degree-play games, prove theorems, solve problems, compose music, do mathematics, and use "natural language" (e.g., English). 







Overview                                                                                     VIII Contrafactus. About how we unconsciously organize our thoughts so that we can imagine hypothetical variants on the real world all the time. Also about aberrant variants of this ability- such as possessed by the new character, the Sloth, an avid lover of French fries, and rabid hater of counterfactuals. 


\chapter {XIX}
how knowledge is represented in layers of contexts. This leads to the modern Al idea of "frames". A frame-like way of handling a set of visual pattern puzzles is presented, for the purpose of concreteness. Then the deep issue of the interaction of concepts in general is discussed, which leads into some speculations on creativity. The Chapter concludes with a set of personal "Questions and Speculations" on Al and minds in general. 

Sloth Canon. A canon which imitates a Bach canon in which one voice plays the same melody as another, only upside down and twice as slowly, while a third voice is free. Here, the Sloth utters the same lines as the Tortoise does, only negated (in a liberal sense of the term) and twice as slowly, while Achilles is free. 


\chapter {XX}
about hierarchical systems and self-reference. It is concerned with the snarls which arise when systems turn back on themselves-for example, science probing science, government investigating governmental wrongdoing, art violating the rules of art, and finally, humans thinking about their own brains and minds. Does Gödel’s Theorem have anything to say about this last "snarl"? Are free will and the sensation of consciousness connected to Gödel’s Theorem? The Chapter ends by tying Gödel, Escher, and Bach together once again. 

Six-Part Ricercar. This Dialogue is an exuberant game played with many of the ideas which have permeated the book. It is a reenactment of the story of the Musical Offering, which began the book; it is simultaneously a "translation" into words of the most complex piece in the Musical Offering: the Six-Part Ricercar. This duality imbues the Dialogue with more levels of meaning than any other in the book. Frederick the Great is replaced by the Crab, pianos by computers, and so on. Many surprises arise. The Dialogue's content concerns problems of mind, consciousness, free will, Artificial Intelligence, the Turing test, and so forth, which have been introduced earlier. It concludes with an implicit reference to the beginning of the book, thus making the book into one big self-referential loop, symbolizing at once Bach's music, Escher's drawings, and Gödel’s Theorem. 







Overview                                                                                           IX FIGURE 1. Johann Sebastian Bach, in 1748. From a painting by Elias Gottlieb Hanssmann. 







Introduction: 

A Musico-Logical Offering Author: FREDERICK THE GREAT, King of Prussia, came to power in 1740. Although he is remembered in history books mostly for his military astuteness, he was also devoted to the life of the mind and the spirit. His court in Potsdam was one of the great centers of intellectual activity in Europe in the eighteenth century. The celebrated mathematician Leonhard Euler spent twenty-five years there. Many other mathematicians and scientists came, as well as philosophers-including Voltaire and La Mettrie, who wrote some of their most influential works while there. But music was Frederick's real love. He was an avid flutist and composer. Some of his compositions are occasionally performed even to this day. Frederick was one of the first patrons of the arts to recognize the virtues of the newly developed "piano-forte" ("soft- loud"). The piano had been developed in the first half of the eighteenth century as a modification of the harpsichord. The problem with the harpsichord was that pieces could only be played at a rather uniform loudness-there was no way to strike one note more loudly than its neighbors. The "soft-loud", as its name implies, provided a remedy to this problem. From Italy, where Bartolommeo Cristofori had made the first one, the soft-loud idea had spread widely. Gottfried Silbermann, the foremost German organ builder of the day, was endeavoring to make a "perfect" piano-forte. Undoubtedly King Frederick was the greatest supporter of his efforts-it is said that the King owned as many as fifteen Silbermann pianos! 

Bach 

Frederick was an admirer not only of pianos, but also of an organist and composer by the name of J. S. Bach. This Bach's compositions were somewhat notorious. Some called them "turgid and confused", while others claimed they were incomparable masterpieces. But no one disputed Bach's ability to improvise on the organ. In those days, being an organist not only meant being able to play, but also to extemporize, and Bach was known far and wide for his remarkable extemporizations. (For some delightful anecdotes about Bach's extemporization, see The Bach Reader, by H. T. David and A. Mendel.) In 1747, Bach was sixty-two, and his fame, as well as one of his sons, had reached Potsdam: in fact, Carl Philipp Emanuel Bach was the Capellmeister (choirmaster) at the court of King Frederick. For years the King had let it be known, through gentle hints to Philipp Emanuel, how 







pleased he would be to have the elder Bach come and pay him a visit; but this wish had never been realized. Frederick was particularly eager for Bach to try out his new Silbermann pianos, which lie (Frederick) correctly foresaw as the great new wave in music. It was Frederick's custom to have evening concerts of chamber music in his court. Often he himself would be the soloist in a concerto for flute Here we have reproduced a painting of such an evening by the German painter Adolph von Menzel, who, in the 1800's, made a series of paintings illustrating the life of Frederick the Great. At the cembalo is C. P. E. Bach, and the figure furthest to the right is Joachim Quantz, the King's flute master-and the only person allowed to find fault with the King's flute playing. One May evening in 1747, an unexpected guest showed up. Johann Nikolaus Forkel, one of Bach's earliest biographers, tells the story as follows: One evening, just as lie was getting his flute ready, and his musicians were ssembled, an officer brought him a list of the strangers who had arrived. With his flute in his hand he ran ever the list, but immediately turned to the assembled musicians, and said, with a kind of agitation, "Gentlemen, old Bach is come." The Hute was now laid aside, and old Bach, who had alighted at his son's lodgings, was immediately summoned to the Palace. Wilhelm Friedemann, who accompanied his father, told me this story, and I must say that 1 still think with pleasure on the manner in which lie related it. At that time it was the fashion to make rather prolix compliments. The first appearance of J. S. Bach before se great a King, who did not even give him time to change his traveling dress for a black chanter's gown, must necessarily be attended with many apologies. I will net here dwell en these apologies, but merely observe, that in Wilhelm Friedemann's mouth they made a formal Dialogue between the King and the Apologist. But what is mere important than this is that the King gave up his Concert for this evening, and invited Bach, then already called the Old Bach, to try his fortepianos, made by Silbermann, which steed in several rooms of the palace. [Forkel here inserts this footnote: "The pianofortes manufactured by Silbermann, of Frevberg, pleased the King se much, that he resolved to buy them all up. He collected fifteen. I hear that they all now stand unfit for use in various corners of the Royal Palace."] The musicians went with him from room to room, and Bach was invited everywhere to try them and to play unpremeditated compositions. After he had gene en for some time, he asked the King to give him a subject for a Fugue, in order to execute it immediately without any preparation. The King admired the learned manner in which his subject was thus executed extempore: and, probably to see hew far such art t could be carried, expressed a wish to hear a Fugue with six Obligato parts. But as it is not every subject that is fit for such full harmony, Bach chose one himself, and immediately executed it to the astonishment of all present in the same magnificent and learned manner as he had done that of the King. His Majesty desired also to hear his performance en the organ. The next day therefore Bach was taken to all the organs in Potsdam, as lie had before been to Silbermann's fortepianos. After his return to Leipzig, he composed the subject, which he had received from the King, in three and six parts. added several artificial passages in strict canon to it, and had it engraved, under the title of "Musikalisches Opfer" [Musical Offering], and dedicated it to the Inventor.' 







Introduction: A Musico-Logical Offering   13 FIGURE 3. The Royal Theme. 

When Bach sent a copy of his Musical Offering to the King, he included a dedicatory letter, which is of interest for its prose style if nothing else rather submissive and flattersome. From a modern perspective it seems comical. Also, it probably gives something of the flavor of Bach's apology for his appearance.2 

MOST GRACIOUS KING! 

In deepest humility I dedicate herewith to Your Majesty a musical offering, the noblest part of which derives from Your Majesty's own august hand. With awesome pleasure I still remember the very special Royal grace when, some time ago, during my visit in Potsdam, Your Majesty's Self deigned to play to me a theme for a fugue upon the clavier, and at the same time charged me most graciously to carry it out in Your Majesty's most august presence. To obey Your Majesty's command was my most humble dim. I noticed very soon, however, that, for lack of necessary preparation, the execution of the task did not fare as well as such an excellent theme demanded. I resoled therefore and promptly pledged myself to work out this right Royal theme more fully, and then make it known to the world. This resolve has now been carried out as well as possible, and it has none other than this irreproachable intent, to glorify, if only in a small point, the fame of a monarch whose greatness and power, as in all the sciences of war and peace, so especially in music, everyone must admire and revere. I make bold to add this most humble request: may Your Majesty deign to dignify the present modest labor with a gracious acceptance, and continue to grant Your Majesty's most august Royal grace to 

Your Majesty's most humble and obedient servant, THE AUTHOR Leipzig, July 7 1747 

Some twenty-seven years later, when Bach had been dead for twentyfour years, a Baron named Gottfried van Swieten-to whom, incidentally, Forkel dedicated his biography of Bach, and Beethoven dedicated his First Symphony-had a conversation with King Frederick, which he reported as follows: 

He [Frederick] spoke to me, among other things, of music, and of a great organist named Bach, who has been for a while in Berlin. This artist [Wilhelm Friedemann Bach] is endowed with a talent superior, in depth of harmonic knowledge and power of execution, to any 1 have heard or can imagine, while those who knew his father claim that he, in turn, was even greater. The King 





is of this opinion, and to prove it to me he sang aloud a chromatic fugue subject which he had given this old Bach, who on the spot had made of it a fugue in four parts, then in five parts, and finally in eight parts.' 

Of course there is no way of knowing whether it was King Frederick or Baron van Swieten who magnified the story into larger-than-life proportions. But it shows how powerful Bach's legend had become by that time. To give an idea of how extraordinary a six-part fugue is, in the entire Well-Tempered Clavier by Bach, containing forty-eight Preludes and Fugues, only two have as many as five parts, and nowhere is there a six-part fugue! One could probably liken the task of improvising a six-part fugue to the playing of sixty simultaneous blindfold games of chess, and winning them all. To improvise an eight-part fugue is really beyond human capability. In the copy which Bach sent to King Frederick, on the page preceding the first sheet of music, was the following inscription: 







FIG URE 4. 


\intro ("At the King's Command, the Song and the Remainder Resolved with Canonic Art.") \endintro
Here Bach is punning on the word "canonic", since it means not only "with canons" but also "in the best possible way". The initials of this inscription are 

RICERCAR 

-an Italian word, meaning "to seek". And certainly there is a great deal to seek in the Musical Offering. It consists of one three-part fugue, one six-part fugue, ten canons, and a trio sonata. Musical scholars have concluded that the three-part fugue must be, in essence, identical with the one which Bach improvised for King Frederick. The six-part fugue is one of Bach's most complex creations, and its theme is, of course, the Royal Theme. That theme, shown in Figure 3, is a very complex one, rhythmically irregular and highly chromatic (that is, filled with tones which do not belong to the key it is in). To write a decent fugue of even two voices based on it would not be easy for the average musician! Both of the fugues are inscribed "Ricercar", rather than "Fuga". This is another meaning of the word; "ricercar" was, in fact, the original name for the musical form now known as "fugue". By Bach's time, the word "fugue" (or fuga, in Latin and Italian) had become standard, but the term "ricercar" had survived, and now designated an erudite kind of fugue, perhaps too austerely intellectual for the common ear. A similar usage survives in English today: the word "recherche" means, literally, "sought out", but carries the same kind of implication, namely of esoteric or highbrow cleverness. The trio sonata forms a delightful relief from the austerity of the fugues and canons, because it is very melodious and sweet, almost dance- 





able. Nevertheless, it too is based largely on the King's theme, chromatic and austere as it is. It is rather miraculous that Bach could use such a theme to make so pleasing an interlude. The ten canons in the Musical Offering are among the most sophisticated canons Bach ever wrote. However, curiously enough, Bach himself never wrote them out in full. This was deliberate. They were posed as puzzles to King Frederick. It was a familiar musical game of the day to give a single theme, together with some more or less tricky hints, and to let the canon based on that theme be "discovered" by someone else. In order to know how this is possible, you must understand a few facts about canons. 

Canons and Fugues 

The idea of a canon is that one single theme is played against itself. This is done by having "copies" of the theme played by the various participating voices. But there are means' ways to do this. The most straightforward of all canons is the round, such as "Three Blind Mice", "Row, Row, Row Your Boat", or " Frere Jacques". Here, the theme enters in the first voice and, after a fixed time-delay, a "copy" of it enters, in precisely the same key. After the same fixed time-delay in the second voice, the third voice enters carrying the theme, and so on. Most themes will not harmonize with themselves in this way. In order for a theme to work as a canon theme, each of its notes must be able to serve in a dual (or triple, or quadruple) role: it must firstly be part of a melody, and secondly it must be part of a harmonization of the same melody. When there are three canonical voices, for instance, each note of the theme must act in two distinct harmonic ways, as well as melodically. Thus, each note in a canon has more than one musical meaning; the listener's ear and brain automatically figure out the appropriate meaning, by referring to context. There are more complicated sorts of canons, of course. The first escalation in complexity comes when the "copies" of the theme are staggered not only in time, but also in pitch; thus, the first voice might sing the theme starting on C, and the second voice, overlapping with the first voice, might sing the identical theme starting five notes higher, on G. A third voice, starting on the D yet five notes higher, might overlap with the first two, and so on. The next escalation in complexity comes when the speeds of' the different voices are not equal; thus, the second voice might sing twice as quickly, or twice as slowly, as the first voice. The former is called diminution, the latter augmentation (since the theme seems to shrink or to expand). We are not yet done! The next stage of complexity in canon construction is to invert the theme, which means to make a melody which jumps down wherever the original theme jumps up, and by exactly the same number of semitones. This is a rather weird melodic transformation, but when one has heard many themes inverted, it begins to seem quite natural. Bach was especially fond of inversions, and they show up often in his work-and the Musical Offering is no exception. (For a simple example of 







inversion, try the tune "Good King Wenceslas". When the original and its inversion are sung together, starting an octave apart and staggered with a time-delay of two beats, a pleasing canon results.) Finally, the most esoteric of "copies" is the retrograde copy- where the theme is played backwards in time. A canon which uses this trick is affectionately known as a crab canon, because of the peculiarities of crab locomotion. Bach included a crab canon in the Musical Offering, needless to say. Notice that every type of "copy" preserves all the information in the original theme, in the sense that the theme is fully recoverable from any of the copies. Such an information preserving transformation is often called an isomorphism, and we will have much traffic with isomorphisms in this book. Sometimes it is desirable to relax the tightness of the canon form. One way is to allow slight departures from perfect copying, for the sake of more fluid harmony. Also, some canons have "free" voices-voices which do not employ the canon's theme, but which simply harmonize agreeably with the voices that are in canon with each other. Each of the canons in the Musical Offering has for its theme a different variant of the King's Theme, and all the devices described above for making canons intricate are exploited to the hilt; in fact, they are occasionally combined. Thus, one three-voice canon is labeled "Canon per Augmentationem, contrario Motu"; its middle voice is free (in fact, it sings the Royal Theme), while the other two dance canonically above and below it, using the devices of augmentation and inversion. Another bears simply the cryptic label "Quaerendo invenietis" ("By seeking, you will discover"). All of the canon puzzles have been solved. The canonical solutions were given by one of Bach's pupils, Johann Philipp Kirnberger. But one might still wonder whether there are more solutions to seek! I should also explain briefly what a fugue is. A fugue is like a canon, in that it is usually based on one theme which gets played in different voices and different keys, and occasionally at different speeds or upside down or backwards. However, the notion of fugue is much less rigid than that of canon, and consequently it allows for more emotional and artistic expression. The telltale sign of a fugue is the way it begins: with a single voice singing its theme. When it is done, then a second voice enters, either five scale-notes up, or four down. Meanwhile the first voice goes on, singing the "countersubject": a secondary theme, chosen to provide rhythmic, harmonic, and melodic contrasts to the subject. Each of the voices enters in turn, singing the theme, often to the accompaniment of the countersubject in some other voice, with the remaining voices doing whatever fanciful things entered the composer's mind. When all the voices have "arrived", then there are no rules. There are, to be sure, standard kinds of things to do-but not so standard that one can merely compose a fugue by formula. The two fugues in the Musical Offering are outstanding examples of fugues that could never have been "composed by formula". There is something much deeper in them than mere fugality. All in all, the Musical Offering represents one of Bach's supreme accomplishments in counterpoint. It is itself one large intellectual fugue, in 







which many ideas and forms have been woven together, and in which playful double meanings and subtle allusions are commonplace. And it is a very beautiful creation of the human intellect which we can appreciate forever. (The entire work is wonderfully described in the book f. S. Bach's Musical Offering, by H. T. David.) 

An Endlessly Rising Canon 

There is one canon in the Musical Offering which is particularly unusual. Labeled simply "Canon per Tonos", it has three voices. The uppermost voice sings a variant of the Royal Theme, while underneath it, two voices provide a canonic harmonization based on a second theme. The lower of this pair sings its theme in C minor (which is the key of the canon as a whole), and the upper of the pair sings the same theme displaced upwards in pitch by an interval of a fifth. What makes this canon different from any other, however, is that when it concludes-or, rather, seems to conclude-it is no longer in the key of C minor, but now is in D minor. Somehow Bach has contrived to modulate (change keys) right under the listener's nose. And it is so constructed that this "ending" ties smoothly onto the beginning again; thus one can repeat the process and return in the key of E, only to join again to the beginning. These successive modulations lead the ear to increasingly remote provinces of tonality, so that after several of them, one would expect to be hopelessly far away from the starting key. And yet magically, after exactly six such modulations, the original key of C minor has been restored! All the voices are exactly one octave higher than they were at the beginning, and here the piece may be broken off in a musically agreeable way. Such, one imagines, was Bach's intention; but Bach indubitably also relished the implication that this process could go on ad infinitum, which is perhaps why he wrote in the margin "As the modulation rises, so may the King's Glory." To emphasize its potentially infinite aspect, I like to call this the "Endlessly Rising Canon". In this canon, Bach has given us our first example of the notion of Strange Loops. The "Strange Loop" phenomenon occurs whenever, by moving upwards (or downwards) through the levels of some hierarchical system, we unexpectedly find ourselves right back where we started. (Here, the system is that of musical keys.) Sometimes I use the term Tangled Hierarchy to describe a system in which a Strange Loop occurs. As we go on, the theme of Strange Loops will recur again and again. Sometimes it will be hidden, other times it will be out in the open; sometimes it will be right side up, other times it will be upside down, or backwards. "Quaerendo invenietis" is my advice to the reader. 



Escher 

To my mind, the most beautiful and powerful visual realizations of this notion of Strange Loops exist in the work of the Dutch graphic artist M. C. Escher, who lived from 1902 to 1972. Escher was the creator of some of the 







FIGURE 5. Waterfall, by M. C. Escher (lithograph, 1961). 

most intellectually stimulating drawings of all time. Many of them have their origin in paradox, illusion, or double-meaning. Mathematicians were among the first admirers of Escher's drawings, and this is understandable because they often are based on mathematical principles of symmetry or pattern ... But there is much more to a typical Escher drawing than just symmetry or pattern; there is often an underlying idea, realized in artistic form. And in particular, the Strange Loop is one of the most recurrent themes in Escher's work. Look, for example, at the lithograph Waterfall (Fig. 5), and compare its six-step endlessly falling loop with the six-step endlessly rising loop of the "Canon per Tonos". The similarity of vision is 







FIGURE 6. Ascending and Descending, by M. C. Escher (lithograph, 1960). 







remarkable. Bach and Escher are playing one single theme in two different "keys": music and art. Escher realized Strange Loops in several different ways, and they can be arranged according to the tightness of the loop. The lithograph Ascending and Descending (Fig. 6), in which monks trudge forever in loops, is the loosest version, since it involves so many steps before the starting point is regained. A tighter loop is contained in Waterfall, which, as we already observed, involves only six discrete steps. You may be thinking that there is some ambiguity in the notion of a single "step"-for instance, couldn't Ascending and Descending be seen just as easily as having four levels (staircases) as forty-five levels (stairs)% It is indeed true that there is an inherent 







FIGURE 7. Hand with Reflecting Globe. Self-portrait In, M. C. Escher (lithograph, 1935). 







Introduction: A Musico-Logical Offering   22 haziness in level-counting, not only in Escher pictures, but in hierarchical, many-level systems. We will sharpen our understanding of this haziness later on. But let us not get too distracted now' As we tighten our loop, we come to the remarkable Drawing Hands (Fig. 135), in which each of two hands draws the other: a two-step Strange Loop. And finally, the tightest of all Strange Loops is realized in Print Gallery (Fig. 142): a picture of a picture which contains itself. Or is it a picture of a gallery which contains itself? Or of a town which contains itself? Or a young man who contains himself'? (Incidentally, the illusion underlying Ascending and Descending and Waterfall was not invented by Escher, but by Roger Penrose, a British mathematician, in 1958. However, the theme of the Strange Loop was already present in Escher's work in 1948, the year he drew Drawing Hands. Print Gallery dates from 1956.) Implicit in the concept of Strange Loops is the concept of infinity, since what else is a loop but a way of representing an endless process in a finite way? And infinity plays a large role n many of Escher's drawings. Copies of one single theme often fit into each' other, forming visual analogues to the canons of Bach. Several such patterns can be seen in Escher's famous print Metamorphosis (Fig. 8). It is a little like the "Endlessly Rising Canon": wandering further and further from its starting point, it suddenly is back. In the tiled planes of Metamorphosis and other pictures, there are already suggestions of infinity. But wilder visions of infinity appear in other drawings by Escher. In some of his drawings, one single theme can appear on different levels of reality. For instance, one level in a drawing might clearly be recognizable as representing fantasy or imagination; another level would be recognizable as reality. These two levels might be the only explicitly portrayed levels. But the mere presence of these two levels invites the viewer to look upon himself as part of yet another level; and by taking that step, the viewer cannot help getting caught up in Escher's implied chain of levels, in which, for any one level, there is always another level above it of greater "reality", and likewise, there is always a level below, "more imaginary" than it is. This can be mind-boggling in itself. However, what happens if the chain of levels is not linear, but forms a loop? What is real, then, and what is fantasy? The genius of Escher was that he could not only concoct, but actually portray, dozens of half-real, half-mythical worlds, worlds filled with Strange Loops, which he seems to be inviting his viewers to enter. 

Gödel 

In the examples we have seen of Strange Loops by Bach and Escher, there is a conflict between the finite and the infinite, and hence a strong sense of paradox. Intuition senses that there is something mathematical involved here. And indeed in our own century a mathematical counterpart was discovered, with the most enormous repercussions. And, just as the Bach and Escher loops appeal to very simple and ancient intuitions-a musical scale, a staircase-so this discovery, by K. Gödel, of a Strange Loop in 







FIGURE 9. Kurt Godel. 







mathematical systems has its origins in simple and ancient intuitions. In its absolutely barest form, Godel's discovery involves the translation of an ancient paradox in philosophy into mathematical terms. That paradox is the so-called Epimenides paradox, or liar paradox. Epimenides was a Cretan who made one immortal statement: "All Cretans are liars." A sharper version of the statement is simply "I am lying"; or, "This statement is false". It is that last version which I will usually mean when I speak of the Epimenides paradox. It is a statement which rudely violates the usually assumed dichotomy of statements into true and false, because if you tentatively think it is true, then it immediately backfires on you and makes you think it is false. But once you've decided it is false, a similar backfiring returns you to the idea that it must be true. Try it! The Epimenides paradox is a one-step Strange Loop, like Escher's Print Gallery. But how does it have to do with mathematics? That is what Godel discovered. His idea was to use mathematical reasoning in exploring mathematical reasoning itself. This notion of making mathematics "introspective" proved to be enormously powerful, and perhaps its richest implication was the one Godel found: Godel's Incompleteness Theorem. What the Theorem states and how it is proved are two different things. We shall discuss both in quite some detail in this book. The Theorem can De likened to a pearl, and the method of proof to an oyster. The pearl is prized for its luster and simplicity; the oyster is a complex living beast whose innards give rise to this mysteriously simple gem. Godel's Theorem appears as Proposition VI in his 1931 paper "On Formally Undecidable Propositions in Principia Mathematica and Related Systems I." It states: 

To every w-consistent recursive class K of formulae there correspond recursive class-signs r, such that neither v Gen r nor Neg (v Gen r) belongs to Fig (K) (where v is the free variable of r). 

Actually, it was in German, and perhaps you feel that it might as well be in German anyway. So here is a paraphrase in more normal English: 

All consistent axiomatic formulations of number theory include undecidable propositions. 

This is the pearl. In this pearl it is hard to see a Strange Loop. That is because the Strange Loop is buried in the oyster-the proof. The proof of Godel's Incompleteness Theorem hinges upon the writing of a self-referential mathematical statement, in the same way as the Epimenides paradox is a self-referential statement of language. But whereas it is very simple to talk about language in language, it is not at all easy to see how a statement about numbers can talk about itself. In fact, it took genius merely to connect the idea of self-referential statements with number theory. Once Godel had the intuition that such a statement could be created, he was over the major hurdle. The actual creation of the statement was the working out of this one beautiful spark of intuition. 







We shall examine the Godel construction quite carefully in Chapters to come, but so that you are not left completely in the dark, I will sketch here, in a few strokes, the core of the idea, hoping that what you see will trigger ideas in your mind. First of all, the difficulty should be made absolutely clear. Mathematical statements-let us concentrate on number- theoretical ones-are about properties of whole numbers. Whole numbers are not statements, nor are their properties. A statement of number theory is not about a. statement of number theory; it just is a statement of number theory. This is the problem; but Godel realized that there was more here than meets the eye. Godel had the insight that a statement of number theory could be about a statement of number theory (possibly even itself), if only numbers could somehow stand for statements. The idea of a code, in other words, is at the heart of his construction. In the Godel Code, usually called "Godel-numbering", numbers are made to stand for symbols and sequences of symbols. That way, each statement of number theory, being a sequence of specialized symbols, acquires a Godel number, something like a telephone number or a license plate, by which it can be referred to. And this coding trick enables statements of number theory to be understood on two different levels: as statements of number theory, and also as statements about statements of number theory. Once Godel had invented this coding scheme, he had to work out in detail a way of transporting the Epimenides paradox into a numbertheoretical formalism. His final transplant of Epimenides did not say, "This statement of number theory is false", but rather, "This statement of number theory does not have any proof". A great deal of confusion can be caused by this, because people generally understand the notion of "proof" rather vaguely. In fact, Godel's work was just part of a long attempt by mathematicians to explicate for themselves what proofs are. The important thing to keep in mind is that proofs are demonstrations within fixed systems of propositions. In the case of Godel's work, the fixed system of numbertheoretical reasoning to which the word "proof" refers is that of Principia Mathematica (P.M.), a giant opus by Bertrand Russell and Alfred North Whitehead, published between 1910 and 1913. Therefore, the Godel sentence G should more properly be written in English as: 

This statement of number theory does not have any proof in the system of Principia Mathematica. 

Incidentally, this Godel sentence G is not Godel's Theorem-no more than the Epimenides sentence is the observation that "The Epimenides sentence is a paradox." We can now state what the effect of discovering G is. Whereas the Epimenides statement creates a paradox since it is neither true nor false, the Godel sentence G is unprovable (inside P.M.) but true. The grand conclusion% That the system of Principia Mathematica is "incomplete"-there are true statements of number theory which its methods of proof are too weak to demonstrate. 







But if Principia Mathematica was the first victim of this stroke, it was certainly not the last! The phrase "and Related Systems" in the title of Godel's article is a telling one: for if Godel's result had merely pointed out a defect in the work of Russell and Whitehead, then others could have been inspired to improve upon P.M. and to outwit Godel's Theorem. But this was not possible: Godel's proof pertained to any axiomatic system which purported to achieve the aims which Whitehead and Russell had set for themselves. And for each different system, one basic method did the trick. In short, Godel showed that provability is a weaker notion than truth, no matter what axiomatic system is involved. Therefore Godel's Theorem had an electrifying effect upon logicians, mathematicians, and philosophers interested in the foundations of mathematics, for it showed that no fixed system, no matter how complicated, could represent the complexity of the whole numbers: 0, 1, 2, 3, ... Modern readers may not be as nonplussed by this as readers of 1931 were, since in the interim our culture has absorbed Godel's Theorem, along with the conceptual revolutions of relativity and quantum mechanics, and their philosophically disorienting messages have reached the public, even if cushioned by several layers of translation (and usually obfuscation). There is a general mood of expectation, these days, of "limitative" results-but back in 1931, this came as a bolt from the blue. 

Mathematical Logic: A Synopsis 

A proper appreciation of Godel's Theorem requires a setting of context. Therefore, I will now attempt to summarize in a short space the history of mathematical logic prior to 1931-an impossible task. (See DeLong, Kneebone, or Nagel and Newman, for good presentations of history.) It all began with the attempts to mechanize the thought processes of reasoning. Now our ability to reason has often been claimed to be what distinguishes us from other species; so it seems somewhat paradoxical, on first thought, to mechanize that which is most human. Yet even the ancient Greeks knew that reasoning is a patterned process, and is at least partially governed by statable laws. Aristotle codified syllogisms, and Euclid codified geometry; but thereafter, many centuries had to pass before progress in the study of axiomatic reasoning would take place again. One of the significant discoveries of nineteenth-century mathematics was that there are different, and equally valid, geometries-where by "a geometry" is meant a theory of properties of abstract points and lines. It had long been assumed that geometry was what Euclid had codified, and that, although there might be small flaws in Euclid's presentation, they were unimportant and any real progress in geometry would be achieved by extending Euclid. This idea was shattered by the roughly simultaneous discovery of non-Euclidean geometry by several people-a discovery that shocked the mathematics community, because it deeply challenged the idea that mathematics studies the real world. How could there be many differ 







ent kinds of "points" and "lines" in one single reality? Today, the solution to the dilemma may be apparent, even to some nonmathematicians-but at the time, the dilemma created havoc in mathematical circles. Later in the nineteenth century, the English logicians George Boole and Augustus De Morgan went considerably further than Aristotle in codifying strictly deductive reasoning patterns. Boole even called his book "The Laws of Thought"-surely an exaggeration, but it was an important contribution. Lewis Carroll was fascinated by these mechanized reasoning methods, and invented many puzzles which could be solved with them. Gottlob Frege in Jena and Giuseppe Peano in Turin worked on combining formal reasoning with the study of sets and numbers. David Hilbert in Gottingen worked on stricter formalizations of geometry than Euclid's. All of these efforts were directed towards clarifying what one means by "proof". In the meantime, interesting developments were taking place in classical mathematics. A theory of different types of infinities, known as the theory of sets, was developed by Georg Cantor in the 1880's. The theory was powerful and beautiful, but intuition-defying. Before long, a variety of set-theoretical paradoxes had been unearthed. The situation was very disturbing, because just as mathematics seemed to be recovering from one set of paradoxes-those related to the theory of limits, in the calculusalong came a whole new set, which looked worse! The most famous is Russell's paradox. Most sets, it would seem, are not members of themselves-for example, the set of walruses is not a walrus, the set containing only Joan of Arc is not Joan of Arc (a set is not a person)-and so on. In this respect, most sets are rather "run-of-the-mill". However, some "self-swallowing" sets do contain themselves as members, such as the set of all sets, or the set of all things except Joan of Arc, and so on. Clearly, every set is either run-of-the-mill or self-swallowing, and no set can be both. Now nothing prevents us from inventing R: the set of all run-o,-the-mill sets. At first, R might seem a rather run-of-the-mill invention-but that opinion must be revised when you ask yourself, "Is R itself "a run-of-the-mill set or a self-swallowing set?" You will find that the answer is: "R is neither run-of-the-mill nor self-swallowing, for either choice leads to paradox." Try it! But if R is neither run-of-the-mill nor self-swallowing, then what is it? At the very least, pathological. But no one was satisfied with evasive answers of that sort. And so people began to dig more deeply into the foundations of set theory. The crucial questions seemed to be: "What is wrong with our intuitive concept of 'set'? Can we make a rigorous theory of sets which corresponds closely with our intuitions, but which skirts the paradoxes?" Here, as in number theory and geometry, the problem is in trying to line up intuition with formalized, or axiomatized, reasoning systems. A startling variant of Russell's paradox, called "Grelling's paradox", can be made using adjectives instead of sets. Divide the adjectives in English into two categories: those which are self-descriptive, such as "pentasyllabic", "awkwardnessful", and "recherche", and those which are not, such 







as "edible", "incomplete", and "bisyllabic". Now if we admit "non-selfdescriptive" as an adjective, to which class does it belong? If it seems questionable to include hyphenated words, we can use two terms invented specially for this paradox: autological (= "self- descriptive"), and heterological (= "non-self-descriptive"). The question then becomes: "Is 'heterological' heterological?" Try it! There seems to he one common culprit in these paradoxes, namely self-reference, or "Strange Loopiness". So if the goal is to ban all paradoxes, why not try banning self- reference and anything that allows it to arise? This is not so easy as it might seem, because it can be hard to figure out just where self-reference is occurring. It may be spread out over a whole Strange Loop with several steps, as in this "expanded" version of Epimenides, reminiscent of Drawing Hands: 

The following sentence is false. The preceding sentence is true. 

Taken together, these sentences have the same effect as the original Epimenides paradox: yet separately, they are harmless and even potentially useful sentences. The "blame" for this Strange Loop can't he pinned on either sentence-only on the way they "point" at each other. In the same way, each local region of Ascending and Descending is quite legitimate; it is only the way they are globally put together that creates an impossibility. Since there are indirect as well as direct ways of achieving self-reference, one must figure out how to ban both types at once-if one sees selfreference as the root of all evil. Banishing Strange Loops Russell and Whitehead did subscribe to this view, and accordingly, Principia Mathematica was a mammoth exercise in exorcising Strange Loops from logic, set theory, and number theory. The idea of their system was basically this. A set of the lowest "type" could contain only "objects" as membersnot sets. A set of the next type up could only contain objects, or sets of the lowest type. In general, a set of a given type could only contain sets of lower type, or objects. Every set would belong to a specific type. Clearly, no set could contain itself because it would have to belong to a type higher than its own type. Only "run-of'-the-mill" sets exist in such a system; furthermore, old R- the set of all run-of-the-mill sets-no longer is considered a set at all, because it does not belong to any finite type. To all appearances, then, this theory of types, which we might also call the "theory of the abolition of Strange Loops", successfully rids set theory of its paradoxes, but only at the cost of introducing an artificial-seeming hierarchy, and of disallowing the formation of certain kinds of sets-such as the set of all run-of-the-mill sets. Intuitively, this is not the way we imagine sets. The theory of types handled Russell's paradox, but it did nothing about the Epimenides paradox or Grelling's paradox. For people whose 







interest went no further than set theory, this was quite adequate-but for people interested in the elimination of paradoxes generally, some similar "hierarchization" seemed necessary, to forbid looping back inside language. At the bottom of such a hierarchy would be an object language. Here, reference could be made only to a specific domain- not to aspects of the object language itself (such as its grammatical rules, or specific sentences in it). For that purpose there would be a metalanguage. This experience of two linguistic levels is familiar to all learners of foreign languages. Then there would be a metametalanguage for discussing the metalanguage, and so on. It would be required that every sentence should belong to some precise level of the hierarchy. Therefore, if one could find no level in which a given utterance fit, then the utterance would be deemed meaningless, and forgotten. An analysis can be attempted on the two-step Epimenides loop given above. The first sentence, since it speaks of the second, must be on a higher level than the second. But by the same token, the second sentence must be on a higher level than the first. Since this is impossible, the two sentences are "meaningless". More precisely, such sentences simply cannot be formulated at all in a system based on a strict hierarchy of languages. This prevents all versions of the Epimenides paradox as well as Grelling's paradox. (To what language level could "heterological" belong?) Now in set theory, which deals with abstractions that we don't use all the time, a stratification like the theory of types seems acceptable, even if a little strange-but when it comes to language, an all-pervading part of life, such stratification appears absurd. We don't think of ourselves as jumping up and down a hierarchy of languages when we speak about various things. A rather matter-of-fact sentence such as, "In this book, I criticize the theory of types" would be doubly forbidden in the system we are discussing. Firstly, it mentions "this book", which should only be mentionable in a metabook"-and secondly, it mentions me-a person whom I should not be allowed to speak of at all! This example points out how silly the theory of types seems, when you import it into a familiar context. The remedy it adopts for paradoxes-total banishment of self-reference in any form-is a real case of overkill, branding many perfectly good constructions as meaningless. The adjective "meaningless", by the way, would have to apply to all discussions of the theory of linguistic types (such as that of this very paragraph) for they clearly could not occur on any of the levels-neither object language, nor metalanguage, nor metametalanguage, etc. So the very act of discussing the theory would be the most blatant possible violation of it! Now one could defend such theories by saying that they were only intended to deal with formal languages-not with ordinary, informal language. This may be so, but then it shows that such theories are extremely academic and have little to say about paradoxes except when they crop up in special tailor-made systems. Besides, the drive to eliminate paradoxes at any cost, especially when it requires the creation of highly artificial formalisms, puts too much stress on bland consistency, and too little on the 







quirky and bizarre, which make life and mathematics interesting. It is of course important to try to maintain consistency, but when this effort forces you into a stupendously ugly theory, you know something is wrong. These types of issues in the foundations of mathematics were responsible for the high interest in codifying human reasoning methods which was present in the early part of this century. Mathematicians and philosophers had begun to have serious doubts about whether even the most concrete of theories, such as the study of whole numbers (number theory), were built on solid foundations. If paradoxes could pop up so easily in set theory-a theory whose basic concept, that of a set, is surely very intuitively appealing- then might they not also exist in other branches of mathematics? Another related worry was that the paradoxes of logic, such as the Epimenides paradox, might turn out to be internal to mathematics, and thereby cast in doubt all of mathematics. This was especially worrisome to those-and there were a good number-who firmly believed that mathematics is simply a branch of logic (or conversely, that logic is simply a branch of mathematics). In fact, this very question-"Are mathematics and logic distinct, or separate%"-was the source of much controversy. This study of mathematics itself became known as metamathematics-or occasionally, metalogic, since mathematics and logic are so intertwined. The most urgent priority of metamathematicians was to determine the true nature of mathematical reasoning. What is a legal method of procedure, and what is an illegal one? Since mathematical reasoning had always been done in "natural language" (e.g., French or Latin or some language for normal communication), there was always a lot of possible ambiguity. Words had different meanings to different people, conjured up different images, and so forth. It seemed reasonable and even important to establish a single uniform notation in which all mathematical work could be done, and with the aid of which any two mathematicians could resolve disputes over whether a suggested proof was valid or not. This would require a complete codification of the universally acceptable modes of human reasoning, at least as far as they applied to mathematics. 

Consistency, Completeness, Hilbert's Program 

This was the goal of Principia Mathematica, which purported to derive all of mathematics from logic, and, to be sure, without contradictions! It was widely admired, but no one was sure if (1) all of mathematics really was contained in the methods delineated by Russell and Whitehead, or (2) the methods given were even self-consistent. Was it absolutely clear that contradictory results could never be derived, by any mathematicians whatsoever, following the methods of Russell and Whitehead? This question particularly bothered the distinguished German mathematician (and metamathematician) David Hilbert, who set before the world community of mathematicians (and metamathematicians) this chal 







lenge: to demonstrate rigorously-perhaps following the very methods outlined by Russell and Whitehead-that the system defined in Principia Mathematica was both consistent (contradiction-free), and complete (i.e., that every true statement of, number theory could be derived within the framework drawn up in P.M.). This was a tall order, and one could criticize it on the grounds that it was somewhat circular: how can you justify your methods of reasoning on the basis of those same methods of reasoning? It is like lifting yourself up by your own bootstraps. (We just don't seem to be able to get away from these Strange Loops!) Hilbert was fully aware of this dilemma, of course, and therefore expressed the hope that a demonstration of consistency or completeness could be found which depended only on "finitistic" modes of reasoning. "these were a small set of reasoning methods usually accepted by mathematicians. In this way, Hilbert hoped that mathematicians could partially lift themselves by their own bootstraps: the sum total of mathematical methods might be proved sound, by invoking only a smaller set of methods. This goal may sound rather esoteric, but it occupied the minds of many of the greatest mathematicians in the world during the first thirty years of this century. In the thirty-first year, however, Godel published his paper, which in some ways utterly demolished Hilbert's program. This paper revealed not only that there were irreparable "holes" in the axiomatic system proposed by Russell and Whitehead, but more generally, that no axiomatic system whatsoever could produce all number-theoretical truths, unless it were an inconsistent system! And finally, the hope of proving the consistency of a system such as that presented in P.M. was shown to be vain: if such a proof could be found using only methods inside P.M., then-and this is one of the most mystifying consequences of Godel's work-P.M. itself would be inconsistent! The final irony of it all is that the proof of Gi del's Incompleteness Theorem involved importing the Epimenides paradox right into the heart ofPrincipia Mathematica, a bastion supposedly invulnerable to the attacks of Strange Loops! Although Godel's Strange Loop did not destroy Principia Mathematica, it made it far less interesting to mathematicians, for it showed that Russell and Whitehead's original aims were illusory. 

Babbage, Computers, Artificial Intelligence ... 

When Godel's paper came out, the world was on the brink of developing electronic digital computers. Now the idea of mechanical calculating engines had been around for a while. In the seventeenth century, Pascal and Leibniz designed machines to perform fixed operations (addition and multiplication). These machines had no memory, however, and were not, in modern parlance, programmable. The first human to conceive of the immense computing potential of machinery was the Londoner Charles Babbage (1792-1871). A character who could almost have stepped out of the pages of the Pickwick Papers, 







Babbage was most famous during his lifetime for his vigorous campaign to rid London of "street nuisances"-organ grinders above all. These pests, loving to get his goat, would come and serenade him at any time of day or night, and he would furiously chase them down the street. Today, we recognize in Babbage a man a hundred years ahead of his time: not only inventor of the basic principles of modern computers, he was also one of the first to battle noise pollution. His first machine, the "Difference Engine", could generate mathematical tables of many kinds by the "method of differences". But before any model of the "D.E." had been built, Babbage became obsessed with a much more revolutionary idea: his "Analytical Engine". Rather immodestly, he wrote, "The course through which I arrived at it was the most entangled and perplexed which probably ever occupied the human mind."' Unlike any previously designed machine, the A.E. was to possess both a "store" (memory) and a "mill" (calculating and decision-making unit). These units were to be built of thousands of intricate geared cylinders interlocked in incredibly complex ways. Babbage had a vision of numbers swirling in and out of the mill tinder control of a program contained in punched cards-an idea inspired by the jacquard loom, a card-controlled loom that wove amazingly complex patterns. Babbage's brilliant but ill-fated Countess friend, Lady Ada Lovelace (daughter of Lord Byron), poetically commented that "the Analytical Engine weaves algebraic patterns just as the Jacquard-loom weaves flowers and leaves." Unfortunately, her use of the present tense was misleading, for no A.E. was ever built, and Babbage died a bitterly disappointed man. Lady Lovelace, no less than Babbage, was profoundly aware that with the invention of the Analytical Engine, mankind was flirting with mechanized intelligence-particularly if the Engine were capable of "eating its own tail" (the way Babbage described the Strange Loop created when a machine reaches in and alters its own stored program). In an 1842 memoir,5 she wrote that the A.E. "might act upon other things besides number". While Babbage dreamt of creating_ a chess or tic-tac-toe automaton, she suggested that his Engine, with pitches and harmonies coded into its spinning cylinders, "might compose elaborate and scientific pieces of music of any degree of complexity or extent." In nearly the same breath, however, she cautions that "The Analytical Engine has no pretensions whatever to originate anything. It can do whatever we know how to order it to perform." Though she well understood the power of artificial computation, Lady Lovelace was skeptical about the artificial creation of intelligence. However, could her keen insight allow her to dream of the potential that would be opened up with the taming of electricity? In our century the time was ripe for computers-computers beyond the wildest dreams of Pascal, Leibniz, Babbage, or Lady Lovelace. In the 1930's and 1940's, the first "giant electronic brains" were designed and built. They catalyzed the convergence of three previously disparate areas: the theory of axiomatic reasoning, the study of mechanical computation, and the psychology of intelligence. These same years saw the theory of computers develop by leaps and 







bounds. This theory was tightly linked to metamathematics. In fact, Godel's Theorem has a counterpart in the theory of computation, discovered by Alan Turing, which reveals the existence of inelucPable "holes" in even the most powerful computer imaginable. Ironically, just as these somewhat eerie limits were being mapped out, real computers were being built whose powers seemed to grow and grow beyond their makers' power of prophecy. Babbage, who once declared he would gladly give up the rest of his life if he could come back in five hundred years and have a three-day guided scientific tour of the new age, would probably have been thrilled speechless a mere century after his death- both by the new machines, and by their unexpected limitations. By the early 1950's, mechanized intelligence seemed a mere stone's throw away; and yet, for each barrier crossed, there always cropped up some new barrier to the actual creation of a genuine thinking machine. Was there some deep reason for this goal's mysterious recession? No one knows where the borderline between non-intelligent behavior and intelligent behavior lies; in fact, to suggest that a sharp borderline exists is probably silly. But essential abilities for intelligence are certainly: 

to respond to situations very flexibly; to take advantage of fortuitous circumstances; to make sense out of ambiguous or contradictory messages; to recognize the relative importance of different elements of a situation; to find similarities between situations despite differences which may separate them; to draw distinctions between situations despite similarities may link them; to synthesize new concepts by taking old them together in new ways; to come up with ideas which are novel. 

Here one runs up against a seeming paradox. Computers by their very nature are the most inflexible, desireless, rule-following of beasts. Fast though they may be, they are nonetheless the epitome of unconsciousness. How, then, can intelligent behavior be programmed? Isn't this the most blatant of contradictions in terms? One of the major theses of this book is that it is not a contradiction at all. One of the major purposes of this book is to urge each reader to confront the apparent contradiction head on, to savor it, to turn it over, to take it apart, to wallow in it, so that in the end the reader might emerge with new insights into the seemingly unbreathable gulf between the formal and the informal, the animate and the inanimate, the flexible and the inflexible. This is what Artificial Intelligence (A1) research is all about. And the strange flavor of AI work is that people try to put together long sets of rules in strict formalisms which tell inflexible machines how to be flexible. What sorts of "rules" could possibly capture all of what we think of as intelligent behavior, however? Certainly there must be rules on all sorts of 







different levels. There must be many "just plain" rules. There must be "metarules" to modify the "just plain" rules; then "metametarules" to modify the metarules, and so on. The flexibility of intelligence comes from the enormous number of different rules, and levels of rules. The reason that so many rules on so many different levels must exist is that in life, a creature is faced with millions of situations of completely different types. In some situations, there are stereotyped responses which require "just plain" rules. Some situations are mixtures of stereotyped situations-thus they require rules for deciding which of the 'just plain" rules to apply. Some situations cannot be classified-thus there must exist rules for inventing new rules ... and on and on. Without doubt, Strange Loops involving rules that change themselves, directly or indirectly, are at the core of intelligence. Sometimes the complexity of our minds seems so overwhelming that one feels that there can be no solution to the problem of understanding intelligence-that it is wrong to think that rules of any sort govern a creature's behavior, even if one takes "rule" in the multilevel sense described above. 

...and Bach 

In the year 1754, four years after the death of J. S. Bach, the Leipzig theologian Johann Michael Schmidt wrote, in a treatise on music and the soul, the following noteworthy passage: Not many years ago it was reported from France that a man had made a statue that could play various pieces on the Fleuttraversiere, placed the flute to its lips and took it down again, rolled its eyes, etc. But no one has yet invented an image that thinks, or wills, or composes, or even does anything at all similar. Let anyone who wishes to be convinced look carefully at the last fugal work of the above-praised Bach, which has appeared in copper engraving, but which was left unfinished because his blindness intervened, and let him observe the art that is contained therein; or what must strike him as even more wonderful, the Chorale which he dictated in his blindness to the pen of another: Wenn wir in hochsten Nothen seen. I am sure that he will soon need his soul if he wishes to observe all the beauties contained therein, let alone wishes to play it to himself or to form a judgment of the author. Everything that the champions of Materialism put forward must fall to the ground in view of this single example.6 Quite likely, the foremost of the "champions of Materialism" here alluded to was none other than Julien Offroy de la Mettrie-philosopher at the court of Frederick the Great, author of L'homme machine ("Man, the Machine"), and Materialist Par Excellence. It is now more than 200 years later, and the battle is still raging between those who agree with Johann Michael Schmidt, and those who agree with Julien Offroy de la Mettrie. I hope in this book to give some perspective on the battle. 

"Godel, Escher, Bach" 

The book is structured in an unusual way: as a counterpoint between Dialogues and Chapters. The purpose of this structure is to allow me to 







present new concepts twice: almost every new concept is first presented metaphorically in a Dialogue, yielding a set of concrete, visual images; then these serve, during the reading of the following`Chapter, as an intuitive background for a more serious and abstract presentation of the same concept. In many of the Dialogues I appear to be talking about one idea on the surface, but in reality I am talking about some other idea, in a thinly disguised way. Originally, the only characters in my Dialogues were Achilles and the Tortoise, who came to me from Zeno of Elea, by way of Lewis Carroll. Zeno of Elea, inventor of paradoxes, lived in the fifth century B.C. One of his paradoxes was an allegory, with Achilles and the Tortoise as protagonists. Zeno's invention of the happy pair is told in my first Dialogue, Three-Part Invention. In 1895, Lewis Carroll reincarnated Achilles and the Tortoise for the purpose of illustrating his own new paradox of infinity. Carroll's paradox, which deserves to be far better known than it is, plays a significant role in this book. Originally titled "What the Tortoise Said to Achilles", it is reprinted here as Two-Part Invention. When I began writing Dialogues, somehow I connected them up with musical forms. I don't remember the moment it happened; I just remember one day writing "Fugue" above an early Dialogue, and from then on the idea stuck. Eventually I decided to pattern each Dialogue in one way or another on a different piece by Bach. This was not so inappropriate. Old Bach himself used to remind his pupils that the separate parts in their compositions should behave like "persons who conversed together as if in a select company". I have taken that suggestion perhaps rather more literally than Bach intended it; nevertheless I hope the result is faithful to the meaning. I have been particularly inspired by aspects of Bach's compositions which have struck me over and over, and which are so well described by David and Mendel in The Bach Reader: His form in general was based on relations between separate sections. These relations ranged from complete identity of passages on the one hand to the return of a single principle of elaboration or a mere thematic allusion on the other. The resulting patterns were often symmetrical, but by no means necessarily so. Sometimes the relations between the various sections make up a maze of interwoven threads that only detailed analysis can unravel. Usually, however, a few dominant features afford proper orientation at first sight or hearing, and while in the course of study one may discover unending sub tleties, one is never at a loss to grasp the unity that holds together every single creation by Bach.' I have sought to weave an Eternal Golden Braid out of these three strands: Godel, Escher, Bach. I began, intending to write an essay at the core of which would be Godel's Theorem. I imagined it would be a mere pamphlet. But my ideas expanded like a sphere, and soon touched Bach and Escher. It took some time for me to think of making this connection explicit, instead of just letting it be a private motivating force. But finally 1 realized that to me, Godel and Escher and Bach were only shadows cast in different directions by some central solid essence. I tried to reconstruct the central object, and came up with this book. 







Three-Part Invention 

Achilles (a Greek warrior, the fleetest of foot of all mortals) and a Tortoise are standing together on a dusty runway in the hot sun. Far down the runway, on a tall flagpole, there hangs a large rectangular flag. The flag is sold red, except where a thin ring-shaped holes has been cut out of it, through which one can see the sky. 


\achilles What is that strange flag down at the other end of the track? It reminds me\par
somehow of a print by my favourite artists M.C. Escher. 
\tortoise That is Zeno’ s flag\par

\achilles Could it be that the hole in it resembles the holes in a Mobian strip Escher once\par
drew? Something is wrong about the flag, I can tell. 
\tortoise The ring which has been cut from it has the shape of the numeral for zero, which\par
is Zeno´s favourite number. 
\achilles The ring which hasn´t been invented yet! It will only be invented by a Hindu\par
mathematician some millennia hence. And thus, Mr. T, mt argument proves that such a flag is impossible. 
\tortoise Your argument is persuasive, Achilles, and I must agree that such a flag is indeed\par
impossible. But it is beautiful anyway, is it not? 
\achilles Oh, yes, there is no doubt of its beauty.\par

\tortoise I wonder if it´s beauty is related to it´s impossibility. I don´t know, I´ve never had\par
the time to analyze Beauty. It´s a Capitalized Essence, and I never seem to have time for Capitalized Essences. 
\achilles Speaking of Capitalized Essences, Mr. T, have you ever wondered about the\par
Purpose of Life? 
\tortoise Oh, heavens, no;\par

\achilles Haven’t you ever wondered why we are here, or who invented us?\par

\tortoise Oh, that is quite another matter. We are inventions of Zeno (as you will shortly\par
see) and the reason we are here is to have a footrace. 
\achilles A footrace? How outrageous! Me, the fleetest of foot of all mortals, versus you,\par
the ploddingest of the plodders! There can be no point to such a race. 
\tortoise You might give me a head start.\par

\achilles It would have to be a huge one.\par

\tortoise I don’t object.\par

\achilles But I will catch you, sooner or later – most likely sooner.\par

\tortoise Not if things go according to Zeno´s paradox, you won’t. Zeno is hoping to use\par
our footrace to show that motion is impossible, you see. It is only in the mind that motion seems possible, according to Zeno. In truth, Motion Is Inherently Impossible. He proves it quite elegantly. 







Figure 10. Mobius strip by M.C.Escher (wood-engraving printed from four blocks, 1961) 


\achilles Oh, yes, it comes back to me now: the famous Zen koan about Zen\par
Master Zeno. As you say it is very simple indeed. 
\tortoise Zen Koan? Zen Master? What do you mean?\par

\achilles It goes like this: Two monks were arguing about a flag. One said, “The\par
flag is moving.” The other said, “The wind is moving.” The sixth patriarch, Zeno, happened to be passing by. He told them, “Not the wind, not the flag, mind is moving.” 
\tortoise I am afraid you are a little befuddled, Achilles. Zeno is no Zen master, far\par
from it. He is in fact, a Greek philosopher from the town of Elea (which lies halfway between points A and B). Centuries hence, he will be celebrated for his paradoxes of motion. In one of those paradoxes, this very footrace between you and me will play a central role. 
\achilles I’m all confused. I remember vividly how I used to repeat over and over\par
the names of the six patriarchs of Zen, and I always said, “The sixth patriarch is Zeno, The sixth patriarch is Zeno…” (Suddenly a soft warm breeze picks up.) Oh, look Mr. Tortoise – look at the flag waving! How I love to watch the ripples shimmer through it’s soft fabric. And the ring cut out of it is waving, too! 






\tortoise Don´t be silly. The flag is impossible, hence it can’t be waving. The wind is\par
waving. 


\intro (At this moment, Zeno happens by.) \endintro



\zeno Hallo! Hulloo! What’s up? What’s new?\par

\achilles The flag is moving.\par

\tortoise The wind is moving.\par

\zeno O friends, Friends! Cease your argumentation! Arrest your vitriolics! Abandon your\par
discord! For I shall resolve the issue for you forthwith. Ho! And on such a fine day. 
\achilles This fellow must be playing the fool.\par

\tortoise No, wait, Achilles. Let us hear what he has to say. Oh Unknown Sir, do impart to\par
us your thoughts on this matter. 
\zeno Most willingly. Not thw ind, not the flag – neither one is moving, nor is anything moving\par
at all. For I have discovered a great Theorem, which states; “Motion Is Inherently Impossible.” And from this Theorem follows an even greater Theorem – Zeno’s Theorem: “Motion Unexists.” 
\achilles “Zeno’s Theorem”? Are you, sir, by any chance, the philosopher Zeno of Elea?\par

\zeno I am indeed, Achilles.\par

\achilles (scratching his head in puzzlement). Now how did he know my name?\par

\zeno Could I possibly persuade you two to hear me out as to why this is the case? I’ve come\par
all the way to Elea from point A this afternoon, just trying to find someone who’ll pay some attention to my closely honed argument. But they’re all hurrying hither and thither, and they don’t have time. You’ve no idea how disheartening it is to meet with refusal after refusal. Oh, I’m sorry to burden you with my troubles, I’d just like to ask you one thing: Would the two of you humour a sill old philosopher for a few moments – only a few, I promise you – in his eccentric theories. 
\achilles Oh, by all means! Please do illuminate us! I know I speak for both of us, since my\par
companion, Mr. Tortoise, was only moments ago speaking of you with great veneration – and he mentioned especially your paradoxes. 
\zeno Thank you. You see, my Master, the fifth patriarch, taught me that reality is one,\par
immutable, and unchanging, all plurality, change, and motion are mere illusions of the sense. Some have mocked his views; but I will show the absurdity of their mockery. My argument is quite simple. I will illustrate it with two characters of my own Invention: Achilles )a Greek warrior, the fleetest of foot of all mortals), and a Tortoise. In my tale, they are persuaded by a passerby to run a footrace down a runway towards a distant flag waving in the breeze. Let us assume that, since the Tortoise is a much slowerrunner, he gets a head start of, say, ten rods. Now the race begins. In a few bounds Achilles has reached the spot where the Tortoise started. 








\achilles Hah!\par

\zeno And now the Tortoise is but a single rod ahead of Achilles. Within only a moment,\par
Achilles has attained that spot. 
\achilles Ho ho!\par

\zeno Yet, in that short moment, the Tortoise has managed to advance a slight amount. In a\par
flash, Achilles covers that distance too. 
\achilles Hee hee hee!\par

\zeno But in that very short flash, the Tortoise has managed to inch ahead by ever so little, and\par
so Achilles is still behind. Now you see that in order for Achilles to catch the Tortoise, this game of “try-to-catch-me” will have to be played an INFINITE number of times – and therefore Achilles can NEVER catch up with the Tortoise. 
\tortoise Heh heh heh heh!\par

\achilles Hmm… Hmm… Hmm… Hmm… Hmm…That argument sounds wrong to me.\par
And yes, I can’t quite make out what’s wrong with it 
\zeno Isn’t it a teaser? It’s my favourite paradox.\par

\tortoise Excuse me, Zeno, but I believe your tale illustrates the wrong principle, doe sit\par
not? You have just told us what will come to known, centuries hence, as Zeno’s “Achilles paradox” , which shows (ahem!) that Achilles will never catch the Tortoise; but the proof that Motion Is Inherently Impossible (and thence that Motion Unexists) is your “dichotomy paradox”, isn’t that so? 
\zeno Oh, shame on me. Of course, you’re right. That’s the new one about how, in going from\par
A to B, one has to go halfway first – and of that stretch one also has to go halfway, and so on and so forth. But you see, both those paradoxes really have the same flavour. Frankly, I’ve only had one Great Idea – I just exploit it in different ways. 
\achilles I swear, these arguments contain a flaw. I don’t quite see where, but they cannot\par
be correct. 
\zeno You doubt the validity of my paradox? Why not just try it out|? You see that red flag\par
waving down here, at the far end of the runway? 
\achilles The impossible one, based on an Escher print?\par

\zeno Exactly. What do you say to you and Mr. Tortoise racing for it, allowing Mr. T a fair\par
head start of, well, I don’t know – 
\tortoise How about ten rods?\par

\zeno Very good – ten rods.\par

\achilles Any time.\par

\zeno Excellent! How exciting! An empirical test of my rigorously proven Theorem! Mr.\par
Tortoise, will you position yourself ten rods upwind? 


\intro (The Tortoise moves ten rods closer to the flag) \endintro


Tortoise and Achlles: Ready! 
\zeno On your mark! Get set! Go!\par









\chapter {1}




The MU-puzzle 

Formal Systems 

ONE OF THE most central notions in this book is that of a formal system. The type of formal system I use was invented by the American logician Emil Post in the 1920's, and is often called a "Post production system". This Chapter introduces you to a formal system and moreover, it is my hope that you will want to explore this formal system at least a little; so to provoke your curiosity, I have posed a little puzzle. "Can you produce MU?" is the puzzle. To begin with, you will be supplied with a string (which means a string of letters).* Not to keep you in suspense, that string will be MI. Then you will be told some rules, with which you can change one string into another. If one of those rules is applicable at some point, and you want to use it, you may, but- there is nothing that will dictate which rule you should use, in case there are several applicable rules. That is left up to you-and of course, that is where playing the game of any formal system can become something of an art. The major point, which almost doesn't need stating, is that you must not do anything which is outside the rules. We might call this restriction the "Requirement of Formality". In the present Chapter, it probably won't need to be stressed at all. Strange though it may sound, though, I predict that when you play around with some of the formal systems of Chapters to come, you will find yourself violating the Requirement of Formality over and over again, unless you have worked with formal systems before. The first thing to say about our formal system-the MIU-system-is that it utilizes only three letters of the alphabet: M, I, U. That means that the only strings of the MIU- system are strings which are composed of those three letters. Below are some strings of the MIU-system: 

MU UIM MUUMUU UIIUMIUUIMUIIUMIUUIMUIIU 

* In this book, we shall employ the following conventions when we refer to strings. When the string is in the same typeface as the text, then it will be enclosed in single or double quotes. Punctuation which belongs to the sentence and not to the string under discussion will go outside of the quotes, as logic dictates. For example, the first letter of this sentence is 'F', while the first letter of 'this ‘sentence’.is 't'. When the string is in Quadrata Roman, however, quotes will usually be left off, unless clarity demands them. For example, the first letter of Quadrata is Q. 







But although all of these are legitimate strings, they are not strings which are "in your possession". In fact, the only string in your possession so far is MI. Only by using the rules, about to be introduced, can you enlarge your private collection. Here is the first rule: 


\rule I: If you possess a string whose last letter is I, you can add on a U at the end.\par


By the way, if up to this point you had not guessed it, a fact about the meaning of "string" is that the letters are in a fixed order. For example, MI and IM are two different strings. A string of symbols is not just a "bag" of symbols, in which the order doesn't make any difference. Here is the second rule: 


\rule II: Suppose you have Mx. Then you may add Mxx to your collection.\par


What I mean by this is shown below, in a few examples. 

From MIU, you may get MIUIU. From MUM, you may get MUMUM. From MU, you may get MUU. 

So the letter `x' in the rule simply stands for any string; but once you have decided which string it stands for, you have to stick with your choice (until you use the rule again, at which point you may make a new choice). Notice the third example above. It shows how, once you possess MU, you can add another string to your collection; but you have to get MU first! I want to add one last comment about the letter `x': it is not part of the formal system in the same way as the three letters `M', `I', and `U' are. It is useful for us, though, to have some way to talk in general about strings of the system, symbolically-and that is the function of the `x': to stand for an arbitrary string. If you ever add a string containing an 'x' to your "collection", you have done something wrong, because strings of the MIU-system never contain "x" “s”! Here is the third rule: 


\rule III: If III occurs in one of the strings in your collection, you may make a new\par
string with U in place of III. 

Examples: 

From UMIIIMU, you could make UMUMU. From MII11, you could make MIU (also MUI). From IIMII, you can't get anywhere using this rule. 
\intro (The three I's have to be consecutive.) \endintro
From MIII, make MU. 

Don't, under any circumstances, think you can run this rule backwards, as in the following example: 





From MU, make MIII               <- This is wrong. 

Rules are one-way. Here is the final rule. 


\rule IV: If UU occurs inside one of your strings, you can drop it.\par


From UUU, get U. From MUUUIII, get MUIII. 

There you have it. Now you may begin trying to make MU. Don't worry you don't get it. Just try it out a bit-the main thing is for you to get the flavor of this MU-puzzle. Have fun. 

Theorems, Axioms, Rules The answer to the MU-puzzle appears later in the book. For now, what important is not finding the answer, but looking for it. You probably hay made some attempts to produce MU. In so doing, you have built up your own private collection of strings. Such strings, producible by the rules, are called theorems. The term "theorem" has, of course, a common usage mathematics which is quite different from this one. It means some statement in ordinary language which has been proven to be true by a rigorous argument, such as Zeno's Theorem about the "unexistence" of motion, c Euclid's Theorem about the infinitude of primes. But in formal system theorems need not be thought of as statements- they are merely strings c symbols. And instead of being proven, theorems are merely produced, as if F machine, according to certain typographical rules. To emphasize this important distinction in meanings for the word "theorem", I will adopt the following convention in this book: when "theorem" is capitalized, its meaning will be the everyday one-a Theorem is a statement in ordinary language which somebody once proved to be true by some sort of logic argument. When uncapitalized, "theorem" will have its technical meaning a string producible in some formal system. In these terms, the MU- puzzle asks whether MU is a theorem of the MIU-system. 

I gave you a theorem for free at the beginning, namely MI. Such "free" theorem is called an axiom-the technical meaning again being qui different from the usual meaning. A formal system may have zero, or several, or even infinitely many axioms. Examples of all these types v appear in the book. 

Every formal system has symbol-shunting rules, such as the four rules of the MIU- system. These rules are called either rules of production or rules of inference. I will use both terms. 

The last term which I wish to introduce at this point is derivation. Shown below is a derivation of the theorem MUIIU: 

(1) MI                            axiom (2) MII                           from (1) by rule II 



(3) MIII                         from (2) by rule II (4) MIIIIU                       from (3) by rule I (5) MUIU                         from (4) by rule III (6) MUIUUIU                      from (5) by rule II (7) MUIIU                        from (6) by rule IV 

A derivation of a theorem is an explicit, line-by-line demonstration of how to produce that theorem according to the rules of the formal system. The concept of derivation is modeled on that of proof, but a derivation is an austere cousin of a proof. It would sound strange to say that you had proven MUIIU, but it does not sound so strange to say you have derived MUIIU. 

Inside and Outside the System Most people go about the MU-puzzle by deriving a number of theorems, quite at random, just to see what kind of thing turns up. Pretty soon, they begin to notice some properties of the theorems they have made; that is where human intelligence enters the picture. For instance, it was probably not obvious to you that all theorems would begin with M, until you had tried a few. Then, the pattern emerged, and not only could you see the pattern, but you could understand it by looking at the rules, which have the property that they make each new theorem inherit its first letter from an earlier theorem; ultimately, then, all theorems' first letters can be traced back to the first letter of the sole axiom MI-and that is a proof that theorems of the MIU-system must all begin with M. There is something very significant about what has happened here. It shows one difference between people and machines. It would certainly be possible-in fact it would be very easy-to program a computer to generate theorem after theorem of the MIU- system; and we could include in the program a command to stop only upon generating U. You now know that a computer so programmed would never stop. And this does not amaze you. But what if you asked a friend to try to generate U? It would not surprise you if he came back after a while, complaining that he can't get rid of the initial M, and therefore it is a wild goose chase. Even if a person is not very bright, he still cannot help making some observations about what he is doing, and these observations give him good insight into the task-insight which the computer program, as we have described it, lacks. Now let me be very explicit about what I meant by saying this shows a difference between people and machines. I meant that it is possible to program a machine to do a routine task in such a way that the machine will never notice even the most obvious facts about what it is doing; but it is inherent in human consciousness to notice some facts about the things one is doing. But you knew this all along. If you punch "1" into an adding machine, and then add 1 to it, and then add 1 again, and again, and again, and continue doing so for hours and hours, the machine will never learn to anticipate you, and do it itself, although any person would pick up the 







pick up the idea, no matter how much or how well it is driven, that it i supposed to avoid other cars and obstacles on the road; and it will never learn even the most frequently traveled routes of its owner. The difference, then, is that it is possible for a machine to act unobservant; it is impossible for a human to act unobservant. Notice I am not saying that all machines are necessarily incapable of making sophisticated observations; just that some machines are. Nor am I saying that all people are always making sophisticated observations; people, in fact, are often very unobservant. But machines can be made to be totally unobservant; any people cannot. And in fact, most machines made so far are pretty close ti being totally unobservant. Probably for this reason, the property of being; unobservant seems to be the characteristic feature of machines, to most people. For example, if somebody says that some task is "mechanical", i does not mean that people are incapable of doing the task; it implies though, that only a machine could do it over and over without eve complaining, or feeling bored. 

Jumping out of the System It is an inherent property of intelligence that it can jump out of the tas which it is performing, and survey what it has done; it is always looking for and often finding, patterns. Now I said that an intelligence can jump out o its task, but that does not mean that it always will. However, a little prompting will often suffice. For example, a human being who is reading a boo may grow sleepy. Instead of continuing to read until the book is finished he is just as likely to put the book aside and turn off the light. He ha stepped "out of the system" and yet it seems the most natural thing in the world to us. Or, suppose person A is watching television when person B comes in the room, and shows evident displeasure with the situation Person A may think he understands the problem, and try to remedy it b exiting the present system (that television program), and flipping the channel knob, looking for a better show. Person B may have a more radio concept of what it is to "exit the system"-namely to turn the television oft Of course, there are cases where only a rare individual will have the vision to perceive a system which governs many peoples lives, a system which ha never before even been recognized as a system; then such people often devote their lives to convincing other people that the system really is there and that it ought to be exited from! How well have computers been taught to jump out of the system? I w cite one example which surprised some observers. In a computer chess: tournament not long ago in Canada, one program-the weakest of all the competing ones-had the unusual feature of quitting long before the game was over. It was not a very good chess player, but it at least had the redeeming quality of being able to spot a hopeless position, and to resign then and there, instead of waiting for the other program to go through the 







boring ritual of checkmating. Although it lost every game it played, it did it in style. A lot of local chess experts were impressed. Thus, if you define "the system" as "making moves in a chess game", it is clear that this program had a sophisticated, preprogrammed ability to exit from the system. On the other hand, if you think of "the system" as being "whatever the computer had been programmed to do", then there is no doubt that the computer had no ability whatsoever to exit from that system. It is very important when studying formal systems to distinguish working within the system from making statements or observations about the system. I assume that you began the MU-puzzle, as do most people, by working within the system; and that you then gradually started getting anxious, and this anxiety finally built up to the point where without any need for further consideration, you exited from the system, trying to take stock of what you had produced, and wondering why it was that you had not succeeded in producing MU. Perhaps you found a reason why you could not produce MU; that is thinking about the system. Perhaps you produced MIU somewhere along the way; that is working within the system. Now I do not want to make it sound as if the two modes are entirely incompatible; I am sure that every human being is capable to some extent of working inside a system and simultaneously thinking about what he is doing. Actually, in human affairs, it is often next to impossible to break things neatly up into "inside the system" and "outside the system"; life is composed of so many interlocking and interwoven and often inconsistent "systems" that it may seem simplistic to think of things in those terms. But it is often important to formulate simple ideas very clearly so that one can use them as models in thinking about more complex ideas. And that is why I am showing you formal systems; and it is about time we went back to discussing the MIU- system. 



M-Mode, I-Mode, U-Mode The MU-puzzle was stated in such a way that it encouraged some amount of exploration within the MIU-system-deriving theorems. But it was also stated in a way so as not to imply that staying inside the system would necessarily yield fruit. Therefore it encouraged some oscillation between the two modes of work. One way to separate these two modes would be to have two sheets of paper; on one sheet, you work "in your capacity as a machine", thus filling it with nothing but M's, I's, and U's; on the second sheet, you work "in your capacity as a thinking being", and are allowed to do whatever your intelligence suggests-which might involve using English, sketching ideas, working backwards, using shorthand (such as the letter `x'), compressing several steps into one, modifying the rules of the system to see what that gives, or whatever else you might dream up. One thing you might do is notice that the numbers 3 and 2 play an important role, since I's are gotten rid of in three's, and U's in two's-and doubling of length (except for the M) is allowed by rule II. So the second sheet might 







also have some figuring on it. We will occasionally refer back to these two modes of dealing with a formal system, and we will call them the Mechanic mode (M-mode) and the Intelligent mode (I-mode). To round out our mode with one for each letter of the MIU-system, I will also mention a fin mode-the Un-mode (U-mode), which is the Zen way of approaching thing. More about this in a few Chapters. 



Decision Procedures An observation about this puzzle is that it involves rules of two opposite tendencies-the lengthening rules and the shortening rules. Two rules (I and II) allow you to increase the size of strings (but only in very rigid, pr scribed ways, of course); and two others allow you to shrink strings somewhat (again in very rigid ways). There seems to be an endless variety to the order in which these different types of rules might be applied, and this gives hope that one way or another, MU could be produced. It might involve lengthening the string to some gigantic size, and then extracting piece after piece until only two symbols are left; or, worse yet, it might involve successive stages of lengthening and then shortening and then lengthening and then shortening, and so on. But there is no guarantee it. As a matter of fact, we already observed that U cannot be produced at all and it will make no difference if you lengthen and shorten till kingdom come. Still, the case of U and the case of MU seem quite different. It is by very superficial feature of U that we recognize the impossibility of producing it: it doesn't begin with an M (whereas all theorems must). It is very convenient to have such a simple way to detect nontheorems. However who says that that test will detect all nontheorems? There may be lots strings which begin with M but are not producible. Maybe MU is one of them. That would mean that the "first-letter test" is of limited usefulness able only to detect a portion of the nontheorems, but missing others. B there remains the possibility of some more elaborate test which discriminates perfectly between those strings which can be produced by the rules and those which cannot. Here we have to face the question, "What do mean by a test?" It may not be obvious why that question makes sense, of important, in this context. But I will give an example of a "test" which somehow seems to violate the spirit of the word. Imagine a genie who has all the time in the world, and who enjoys using it to produce theorems of the MIU-system, in a rather methodical way. Here, for instance, is a possible way the genie might go about it 

Step 1: Apply every applicable rule to the axiom MI. This yields two new theorems MIU, MII. Step 2: Apply every applicable rule to the theorems produced in step 1. This yields three new theorems: MIIU, MIUIU, MIIII. 







Step 3: Apply every applicable rule to the theorems produced in step 2. This yields five new theorems: MIIIIU, MIIUIIU, MIUIUIUIU, MIIIIIIII, MUI. 



This method produces every single theorem sooner or later, because the rules are applied in every conceivable order. (See Fig. 11.) All of the lengthening-shortening alternations which we mentioned above eventually get carried out. However, it is not clear how long to wait for a given string 







FIGURE 11. A systematically constructed "tree" of all the theorems of the MIU-system. The N th level down contains those theorems whose derivations contain exactly N steps. The encircled numbers tell which rule was employed. Is MU anywhere in this tree? 



to appear on this list, since theorems are listed according to the shortness of their derivations. This is not a very useful order, if you are interested in a specific string (such as MU), and you don't even know if it has any derivation, much less how long that derivation might be. Now we state the proposed "theoremhood-test": 

Wait until the string in question is produced; when that happens, you know it is a theorem-and if it never happens, you know that it is not a theorem. 

This seems ridiculous, because it presupposes that we don't mind waiting around literally an infinite length of time for our answer. This gets to the crux of the matter of what should count as a "test". Of prime importance is a guarantee that we will get our answer in a finite length of time. If there is a test for theoremhood, a test which does always terminate in a finite 







amount of time, then that test is called a decision procedure for the given formal system. When you have a decision procedure, then you have a very concrete characterization of the nature of all theorems in the system. Offhand, it might seem that the rules and axioms of the formal system provide no less complete a characterization of the theorems of the system than a decision procedure would. The tricky word here is "characterization". Certainly the rules of inference and the axioms of the MIU-system do characterize, implicitly, those strings that are theorems. Even more implicitly, they characterize those strings that are not theorems. But implicit characterization is not enough, for many purposes. If someone claims to have a characterization of all theorems, but it takes him infinitely long to deduce that some particular string is not a theorem, you would probably tend to say that there is something lacking in that characterization-it is not quite concrete enough. And that is why discovering that a decision procedure exists is a very important step. What the discovery means, in effect, is that you can perform a test for theoremhood of a string, and that, even if the test is complicated, it is guaranteed to terminate. In principle, the test is just as easy, just as mechanical, just as finite, just as full of certitude, as checking whether the first letter of the string is M. A decision procedure is a "litmus test" for theoremhood! Incidentally, one requirement on formal systems is that the set of axioms must be characterized by a decision procedure-there must be a litmus test for axiomhood. This ensures that there is no problem in getting off the ground at the beginning, at least. That is the difference between the set of axioms and the set of theorems: the former always has a decision procedure, but the latter may not. I am sure you will agree that when you looked at the MIU-system for the first time, you had to face this problem exactly. The lone axiom was known, the rules of inference were simple, so the theorems had been implicitly characterized-and yet it was still quite unclear what the consequences of that characterization were. In particular, it was still totally unclear whether MU is, or is not, a theorem. 







FIGURE 12. Sky Castle, by M. C.: Escher (woodcut, 1928). 







Two-Part Invention or, What the Tortoise Said to Achilles by Lewis Carroll' 

Achilles had overtaken the Tortoise, and had seated himself comfortably on its back. "So you've got to the end of our race-course?" said the Tortoise. "Even though it DOES consist of an infinite series of distances? I thought some wiseacre or other had proved that the thing couldn't be done?" "It CAN be done," said Achilles. "It HAS been done! Solvitur ambulando. You see the distances were constantly DIMINISHING; and so-" "But if they had been constantly INCREASING?" the Tortoise interrupted. "How then?" "Then I shouldn't be here," Achilles modestly replied; "and You would have got several times round the world, by this time!" "You flatter me-FLATTEN, I mean," said the Tortoise; "for you ARE a heavy weight, and NO mistake! Well now, would you like to hear of a race-course, that most people fancy they can get to the end of in two or three steps, while it REALLY consists of an infinite number of distances, each one longer than the previous one?" "Very much indeed!" said the Grecian warrior, as he drew from his helmet (few Grecian warriors possessed POCKETS in those days) an enormous note-book and pencil. "Proceed! And speak SLOWLY, please! SHORTHAND isn't invented yet!" "That beautiful First Proposition by Euclid!" the Tortoise murmured dreamily. "You admire Euclid?" "Passionately! So far, at least, as one CAN admire a treatise that won't be published for some centuries to come!" "Well, now, let's take a little bit of the argument in that First Proposition just TWO steps, and the conclusion drawn from them. Kindly enter them in your note-book. And in order to refer to them conveniently, let's call them A, B, and Z: 

(A) Things that are equal to the same are equal to each other. (B) The two sides of this Triangle are things that are equal to the same. (Z) The two sides of this Triangle are equal to each other. 

Readers of Euclid will grant, I suppose, that Z follows logically from A and B, so that any one who accepts A and B as true, MUST accept Z as true?" "Undoubtedly! The youngest child in a High School-as soon as High 







Schools are invented, which will not be till some two thousand years later-will grant THAT." "And if some reader had NOT yet accepted A and B as true, he might still accept the SEQUENCE as a VALID one, I suppose?" "No doubt such a reader might exist. He might say, `I accept as true the Hypothetical Proposition that, IF A and B be true, Z must be true; but I DON'T accept A and B as true.' Such a reader would do wisely in abandoning Euclid, and taking to football." "And might there not ALSO be some reader who would say `I accept A and B as true, but I DON'T accept the Hypothetical'?" "Certainly there might. HE, also, had better take to football." "And NEITHER of these readers," the Tortoise continued, "is AS YET under any logical necessity to accept Z as true?" "Quite so," Achilles assented. "Well, now, I want you to consider ME as a reader of the SECOND kind, and to force me, logically, to accept Z as true." "A tortoise playing football would be-" Achilles was beginning. `-an anomaly, of course," the Tortoise hastily interrupted. "Don't wander from the point. Let's have Z first, and football afterwards!" "I'm to force you to accept Z, am I?" Achilles said musingly. "And your present position is that you accept A and B, but you DON'T accept the Hypothetical-" "Let's call it C," said the Tortoise. "-but you DON'T accept 

(C) If A and B are true, Z must be true." 

"That is my present position," said the Tortoise. "Then I must ask you to accept C." "I'll do so," said the Tortoise, "as soon as you've entered it in that notebook of yours. What else have you got in it?" "Only a few memoranda," said Achilles, nervously fluttering the leaves: "a few memoranda of-of the battles in which I have distinguished myself!" "Plenty of blank leaves, I see!" the Tortoise cheerily remarked. "We shall need them ALL!" (Achilles shuddered.) "Now write as I dictate: 

(A) Things that are equal to the same are equal to each other. (B) The two sides of this Triangle are things that are equal to the same. (C) If A and B are true, Z must be true. (Z) The two sides of this Triangle are equal to each other." 

"You should call it D, not Z," said Achilles. "It comes NEXT to the other three. If you accept A and B and C, you MUST accept Z. 







“And why must I?” "Because it follows LOGICALLY from them. If A and B and C are true, Z MUST be true. You can't dispute THAT, I imagine?" "If A and B and C are true, Z MUST be true," the Tortoise thoughtfully repeated. "That's ANOTHER Hypothetical, isn't it? And, if I failed to see its truth, I might accept A and B and C, and STILL not accept Z, mightn't I?" "You might," the candid hero admitted; "though such obtuseness would certainly be phenomenal. Still, the event is POSSIBLE. So I must ask you to grant ONE more Hypothetical." "Very good, I'm quite willing to grant it, as soon as you've written it down. We will call it 

(D) If A and B and C are true, Z must be true. 

Have you entered that in your note-book?" "I HAVE!" Achilles joyfully exclaimed, as he ran the pencil into its sheath. "And at last we've got to the end of this ideal race-course! Now that you accept A and B and C and D, OF COURSE you accept Z." "Do I?" said the Tortoise innocently. "Let's make that quite clear. I accept A and B and C and D. Suppose I STILL refused to accept Z?" "Then Logic would take you by the throat, and FORCE you to do it!" Achilles triumphantly replied. "Logic would tell you, `You can't help yourself. Now that you've accepted A and B and C and D, you MUST accept Z!' So you've no choice, you see.", "Whatever LOGIC is good enough to tell me is worth WRITING DOWN," said the Tortoise. "So enter it in your book, please. We will call it 

(E) If A and B and C and D are true, Z must be true. 

Until I've granted THAT, of course I needn't grant Z. So it's quite a NECESSARY step, you see?" "I see," said Achilles; and there was a touch of sadness in his tone. Here the narrator, having pressing business at the Bank, was obliged to leave the happy pair, and did not again pass the spot until some months afterwards. When he did so, Achilles was still seated on the back of the much-enduring Tortoise, and was writing in his notebook, which appeared to be nearly full. The Tortoise was saying, "Have you got that last step written down? Unless I've lost count, that makes a thousand and one. There are several millions more to come. And WOULD you mind, as a personal favour, considering what a lot of instruction this colloquy of ours will provide for the Logicians of the Nineteenth Century-WOULD you mind adopting a pun that my cousin the Mock- Turtle will then make, and allowing yourself to be renamed TAUGHT-US?" "As you please," replied the weary warrior, in the hollow tones of despair, as he buried his face in his hands. "Provided that YOU, for YOUR part, will adopt a pun the Mock- Turtle never made, and allow yourself to be re-named A KILL-EASE!" 








\chapter {11}




Meaning and Form in Mathematics. 

THIS Two-Part Invention was the inspiration for my two characters. Just as Lewis Carroll took liberties with Zeno's Tortoise and Achilles, so have I taken liberties with Lewis Carroll's Tortoise and Achilles. In Carroll's dialogue, the same events take place over and over again, only each time on a higher and higher level; it is a wonderful analogue to Bach's Ever-Rising Canon. The Carrollian Dialogue, with its wit subtracted out, still leaves a deep philosophical problem: Do words and thoughts follow formal rules, or do they not? That problem is the problem of this book. In this Chapter and the next, we will look at several new formal systems. This will give us a much wider perspective on the concept of formal system. By the end of these two Chapters, you should have quite a good idea of the power of formal systems, and why they are of interest to mathematicians and logicians. 

The pq-System The formal system of this Chapter is called the pq-system. It is not important to mathematicians or logicians-in fact, it is just a simple invention of mine. Its importance lies only in the fact that it provides an excellent example of many ideas that play a large role in this book. There are three distinct symbols of the pq-system: 

p       q    - -the letters p, q, and the hyphen. The pq-system has an infinite number of axioms. Since we can't write them all down, we have to have some other way of describing what they are. Actually, we want more than just a description of the axioms; we want a way to tell whether some given string is an axiom or not. A mere description of axioms might characterize them fully and yet weakly-which was the problem with the way theorems in the MIU-system were characterized. We don't want to have to struggle for an indeterminate-possibly infinite length of time, just to find out if some string is an axiom or not. Therefore, we will define axioms in such a way that there is an obvious decision procedure for axiomhood of a string composed of p's, q's, and hyphens. 







DEFINITION: xp-qx is an axiom, whenever x is composed of hyphens only. 

Note that 'x' must stand for the same string of hyphens in both occurrences For example, - -p-q---is an axiom. The literal expression `xp-qx-' i,, not an axiom, of course (because `x' does not belong to the pq-system); it is more like a mold in which all axioms are cast-and it is called an axiom schema. The pq-system has only one rule of production: 

RULE: Suppose x, y, and z all stand for particular strings containing only hyphens. And suppose that x py qz is known to be a theorem. The` xpy-qz- is a theorem. 

For example, take x to be'--', y to be'---', and z to be'-'. The rule tells us: 

If --p---q- turns out to be a theorem, then so will --p----q--. 



As is typical of rules of production, the statement establishes a causal connection between the theoremhood of two strings, but without asserting theoremhood for either one on its own. A most useful exercise for you is to find a decision procedure for the theorems of the pq-system. It is not hard; if you play around for a while you will probably pick it up. Try it. 



The Decision Procedure I presume you have tried it. First of all, though it may seem too obvious to mention, I would like to point out that every theorem of the pq-system has three separate groups of hyphens, and the separating elements are one p, and one q, in that order. (This can be shown by an argument based on "heredity", just the way one could show that all MIU- system theorems had to begin with M.) This means that we can rule out, from its form alone, o string such as --p--p--p--q . Now, stressing the phrase "from its form alone" may seem silly; what else is there to a string except its form? What else could possibly play a roll in determining its properties? Clearly nothing could. But bear this in mint as the discussion of formal systems goes on; the notion of "form" will star to get rather more complicated and abstract, and we will have to think more about the meaning of the word "form". In any case, let us give the name well formed string to any string which begins with a hyphen-group, then ha one p, then has a second hyphen-group, then a q, and then a final hyphen-group. Back to the decision procedure ... The criterion for theoremhood is that the first two hyphen-groups should add up, in length, to the third 







hyphen-group. for instance, --p--q - is a theorem, since 2 plus 2 equals 4, whereas --p--q- is not, since 2 plus 2 is not 1. To see why this is the proper criterion, look first at the axiom schema. Obviously, it only manufactures axioms which satisfy the addition criterion. Second, look at the rule of production. If the first string satisfies the addition criterion, so must the second one-and conversely, if the first string does not satisfy the addition criterion, then neither does the second string. The rule makes the addition criterion into a hereditary property of theorems: any theorem passes the property on to its offspring. This shows why the addition criterion is correct. There is, incidentally, a fact about the pq-system which would enable us to say with confidence that it has a decision procedure, even before finding the addition criterion. That fact is that the pq-system is not complicated by the opposing currents of lengthening and shortening rules; it has only lengthening rules. Any formal system which tells you how to make longer theorems from shorter ones, but never the reverse, has got to have a decision procedure for its theorems. For suppose you are given a string. First check whether it's an axiom or not (I am assuming that there is a decision procedure for axiomhood-otherwise, things are hopeless). If it is an axiom, then it is by definition a theorem, and the test is over. So suppose instead that it's not an axiom. Then, to be a theorem, it must have come from a shorter string, via one of the rules. By going over the various rules one by one, you can pinpoint not only the rules that could conceivably produce that string, but also exactly which shorter strings could be its forebears on the "family tree". In this way, you "reduce" the problem to determining whether any of several new but shorter strings is a theorem. Each of them can in turn be subjected to the same test. The worst that can happen is a proliferation of more and more, but shorter and shorter, strings to test. As you continue inching your way backwards in this fashion, you must be getting closer to the source of all theorems-the axiom schemata. You just can't get shorter and shorter indefinitely; therefore, eventually either you will find that one of your short strings is an axiom, or you'll come to a point where you're stuck, in that none of your short strings is an axiom, and none of them can be further shortened by running some rule or other backwards. This points out that there really is not much deep interest in formal systems with lengthening rules only; it is the interplay of lengthening and shortening rules that gives formal systems a certain fascination.. 

Bottom-up vs. Top-down 

The method above might be called a top-down decision procedure, to be contrasted with a bottom-up decision procedure, which I give now. It is very reminiscent of the genie's systematic theorem-generating method for the MIU-system, but is complicated by the presence of an axiom schema. We are going to form a "bucket" into which we throw theorems as they are generated. Here is how it is done: 







(1a) Throw the simplest possible axiom (-p-q--) into the bucket. (I b) Apply the rule of inference to the item in the bucket, and put the result into the bucket. (2a) Throw the second-simplest axiom into the bucket. (2b) Apply the rule to each item in the bucket, and throw all results into the bucket. (3a) Throw the third-simplest axiom into the bucket. (3b) Apply the rule to each item in the bucket, and throw all results into the bucket. 

etc., etc. 

A moment's reflection will show that you can't fail to produce every theorem of the pq- system this way. Moreover, the bucket is getting filled with longer and longer theorems, as time goes on. It is again a consequence of that lack of shortening rules. So if you have a particular string, such as --p---q---- , which you want to test for theoremhood, just follow the numbered steps, checking all the while for the string in question. If it turns up- theorem! If at some point everything that goes into the bucket is longer than the string in question, forget it-it is not a theorem. This decision procedure is bottom=up because it is working its way up from the basics, which is to say the axioms. The previous decision procedure is top-down because it does precisely the reverse: it works its way back down towards the basics. 

Isomorphisms Induce Meaning Now we come to a central issue of this Chapter-indeed of the book. Perhaps you have already thought to yourself that the pq-theorems are like additions. The string --p---q--- is a theorem because 2 plus 3 equals 5. It could even occur to you that the theorem --p---q-- is a statement, written in an odd notation, whose meaning is that 2 plus 3 is 5. Is this a reasonable way to look at things? Well, I deliberately chose 'p' to remind you of 'plus', and 'q' to remind you of 'equals' . . . So, does the string --p---q---- actually mean "2 plus 3 equals 5"? What would make us feel that way? My answer would be that we have perceived an isomorphism between pq-theorems and additions. In the Introduction, the word "isomorphism" was defined as an information preserving transformation. We can now go into that notion a little more deeply, and see it from another perspective. The word "isomorphism' applies when two complex structures can be mapped onto each other, in such a way that to each part of one structure there is a corresponding part in the other structure, where "corresponding" means that the two part play similar roles in their respective structures. This usage of the word "isomorphism" is derived from a more precise notion in mathematics. 







It is cause for joy when a mathematician discovers an isomorphism between two structures which he knows. It is often a "bolt from the blue", and a source of wonderment. The perception of an isomorphism between two known structures is a significant advance in knowledge-and I claim that it is such perceptions of isomorphism which create meanings in the minds of people. A final word on the perception of isomorphisms: since they come in many shapes and sizes, figuratively speaking, it is not always totally clear when you really have found an isomorphism. Thus, "isomorphism" is a word with all the usual vagueness of words-which is a defect but an advantage as well. In this case, we have an excellent prototype for the concept of isomorphism. There is a "lower level" of our isomorphism-that is, a mapping between the parts of the two structures: 

p <= => plus q <= => equals - <= => one -- <= => two ---- <= => three etc. 

This symbol-word correspondence has a name: interpretation. Secondly, on a higher level, there is the correspondence between true statements and theorems. But-note carefully-this higher-level correspondence could not be perceived without the prior choice of an interpretation for the symbols. Thus it would be more accurate to describe it as a correspondence between true statements and interpreted theorems. In any case we have displayed a two-tiered correspondence, which is typical of all isomorphisms. When you confront a formal system you know nothing of, and if you hope to discover some hidden meaning in it, your problem is how to assign interpretations to its symbols in a meaningful way-that is, in such a way that a higher-level correspondence emerges between true statements and theorems. You may make several tentative stabs in the dark before finding a good set of words to associate with the symbols. It is very similar to attempts to crack a code, or to decipher inscriptions in an unknown language like Linear B of Crete: the only way to proceed is by trial and error, based on educated guesses. When you hit a good choice, a "meaningful" choice, all of a sudden things just feel right, and work speeds up enormously. Pretty soon everything falls into place. The excitement of such an experience is captured in The Decipherment of Linear B by John Chadwick. But it is uncommon, to say the least, for someone to be in the position of "decoding" a formal system turned up in the excavations of a ruined civilization! Mathematicians (and more recently, linguists, philosophers, and some others) are the only users of formal systems, and they invariably have an interpretation in mind for the formal systems which they use and publish. The idea of these people is to set up a formal system whose 







Theorems reflect some portion of reality isomorphically. In such a case, the choice of symbols is a highly motivated one, as is the choice of typographical rules of production. When I devised the pq-system, I was in position. You see why I chose the symbols I chose. It is no accident theorems are isomorphic to additions; it happened because I deliberately sought out a way to reflect additions typographically. 

Meaningless and Meaningful Interpretations You can choose interpretations other than the one I chose. You need make every theorem come out true. But there would be very little reason make an interpretation in which, say, all theorems came out false, certainly even less reason to make an interpretation under which there is no correlation at all, positive or negative, between theoremhood and tri Let us therefore make a distinction between two types of interpretations a formal system. First, we can have a meaningless interpretation, one un which we fail to see any isomorphic connection between theorems of system, and reality. Such interpretations abound-any random choice a will do. For instance, take this one: 

p <= => horse q <= => happy - <= => apple 



Now -p-q-- acquires a new interpretation: "apple horse apple hat apple apple"-a poetic sentiment, which might appeal to horses, and mi! even lead them to favor this mode of interpreting pq-strings! However, t interpretation has very little "meaningfulness"; under interpretative, theorems don't sound any truer, or any better, than nontheorems. A ho might enjoy "happy happy happy apple horse" (mapped onto q q q) just as much as any interpreted theorem. The other kind of interpretation will be called meaningful. Under si an interpretation, theorems and truths correspond-that is, an isomorphism exists between theorems and some portion of reality. That is why it is good to distinguish between interpretations and meanings. Any old word can be used as an interpretation for `p', but `plus' is the only meaningful choice we've come up with. In summary, the meaning of `p' seems to be 'plus’ though it can have a million different interpretations. 

Active vs. Passive Meanings Probably the most significant fact of this Chapter, if understood deeply this: the pq- system seems to force us into recognizing that symbols of a formal system, though initially without meaning, cannot avoid taking on "meaning" of sorts at least if an isomorphism is found. The difference between meaning it formal system and in a language is a very important one, however. It is this: 







in a language, when we have learned a meaning for a word, we then mar-c new statements based on the meaning of the word. In a sense the meaning becomes active, since it brings into being a new rule for creating sentences. This means that our command of language is not like a finished product: the rules for making sentences increase when we learn new meanings. On the other hand, in a formal system, the theorems are predefined, by the rules of production. We can choose "meanings" based on an isomorphism (if we can find one) between theorems and true statements. But this does not give us the license to go out and add new theorems to the established theorems. That is what the Requirement of Formality in Chapter I was warning you of. In the MIU-system, of course, there was no temptation to go beyond the four rules, because no interpretation was sought or found. But here, in our new system, one might be seduced by the newly found "meaning" of each symbol into thinking that the string 

--p--p--p--q 

is a theorem. At least, one might wish that this string were a theorem. But wishing doesn't change the fact that it isn't. And it would be a serious mistake to think that it "must" be a theorem, just because 2 plus 2 plus 2 plus 2 equals 8. It would even be misleading to attribute it any meaning at all, since it is not well-formed, and our meaningful interpretation is entirely derived from looking at well-formed strings. In a formal system, the meaning must remain passive; we can read each string according to the meanings of its constituent symbols, but we do not have the right to create new theorems purely on the basis of the meanings we've assigned the symbols. Interpreted formal systems straddle the line between systems without meaning, and systems with meaning. Their strings can be thought of as "expressing" things, but this must come only as a consequence of the formal properties of the system. 

Double-Entendre! And now, I want to destroy any illusion about having found the meanings for the symbols of the pq-system. Consider the following association: 

p <= => equals q <= => taken from - <= => one -- <= => two etc. 

Now, --p---q---- has a new interpretation: "2 equals 3 taken from 5". Of course it is a true statement. All theorems will come out true under this new interpretation. It is just as meaningful as the old one. Obviously, it is silly to ask, "But which one is the meaning of the string?" An interpreta 







tion will me meaningful to the extent that it accurately reflects some isomorphism to the real world. When different aspects of the real world a isomorphic to each other (in this case, additions and subtractions), or single formal system can be isomorphic to both, and therefore can take ( two passive meanings. This kind of double-valuedness of symbols at strings is an extremely important phenomenon. Here it seems trivial curious, annoying. But it will come back in deeper contexts and bring with it a great richness of ideas. Here is a summary of our observations about the pq-system. Und either of the two meaningful interpretations given, every well-form( string has a grammatical assertion for its counterpart-some are true, son false. The idea of well formed strings in any formal system is that they a those strings which, when interpreted symbol for symbol, yield grammatical sentences. (Of course, it depends on the interpretation, but usually, there one in mind.) Among the well-formed strings occur the theorems. The: are defined by an axiom schema, and a rule of production. My goal in inventing the pq-system was to imitate additions: I wanted every theorem] to express a true addition under interpretation; conversely, I wanted every true addition of precisely two positive integers to be translatable into a string, which would be a theorem. That goal was achieved. Notice, then fore, that all false additions, such as "2 plus 3 equals 6", are mapped into strings which are well-formed, but which are not theorems. 



Formal Systems and Reality This is our first example of 'a case where a formal system is based upon portion of reality, and seems to mimic it perfectly, in that its theorems a] isomorphic to truths about that part of reality. However, reality and tt formal system are independent. Nobody need be aware that there is a isomorphism between the two. Each side stands by itself-one plus or equals two, whether or not we know that -p-q-- is a theorem; and -p-q-- is still a theorem whether or not we connect it with addition. You might wonder whether making this formal system, or any form system, sheds new light on truths in the domain of its interpretation. Hat we learned any new additions by producing pq-theorems? Certainly not but we have learned something about the nature of addition as process-namely, that it is easily mimicked by a typographical rule governing meaningless symbols. This still should not be a big surprise sing addition is such a simple concept. It is a commonplace that addition can I captured in the spinning gears of a device like a cash register. But it is clear that we have hardly scratched the surface, as far formal systems go; it is natural to wonder about what portion of reality co be imitated in its behavior by a set of meaningless symbols governed I formal rules. Can all of reality be turned into a formal system? In a very broad sense, the answer might appear to be yes. One could suggest, for instance, that reality is itself nothing but one very complicated formal 







system. Its symbols do not move around on paper, but rather in a three-dimensional vacuum (space); they are the elementary particles of which everything is composed. (Tacit assumption: that there is an end to the descending chain of matter, so that the expression "elementary particles" makes sense.) The "typographical rules" are the laws of physics, which tell how, given the positions and velocities of all particles at a given instant, to modify them, resulting in a new set of positions and velocities belonging to the "next" instant. So the theorems of this grand formal system are the possible configurations of particles at different times in the history of the universe. The sole axiom is (or perhaps, was) the original configuration of all the particles at the "beginning of time". This is so grandiose a conception, however, that it has only the most theoretical interest; and besides, quantum mechanics (and other parts of physics) casts at least some doubt on even the theoretical worth of this idea. Basically, we are asking if the universe operates deterministically, which is an open question. 

Mathematics and Symbol Manipulation Instead of dealing with such a big picture, let's limit ourselves to mathematics as our "real world". Here, a serious question arises: How can we be sure, if we've tried to model a formal system on some part of mathematics, that we've done the job accurately- especially if we're not one hundred per cent familiar with that portion of mathematics already? Suppose the goal of the formal system is to bring us new knowledge in that discipline. How will we know that the interpretation of every theorem is true, unless we've proven that the isomorphism is perfect? And how will we prove that the isomorphism is perfect, if we don't already know all about the truths in the discipline to begin with? Suppose that in an excavation somewhere, we actually did discover some mysterious formal system. We would try out various interpretations and perhaps eventually hit upon one which seemed to make every theorem come out true, and every nontheorem come out false. But this is something which we could only check directly in a finite number of cases. The set of theorems is most likely infinite. How will we know that all theorems express truths under this interpretation, unless we know everything there is to know about both the formal system and the corresponding domain of interpretation? It is in somewhat this odd position that we will find ourselves when we attempt to match the reality of natural numbers (i.e., the nonnegative integers: 0, 1, 2, ...) with the typographical symbols of a formal system. We will try to understand the relationship between what we call "truth" in number theory and what we can get at by symbol manipulation. So let us briefly look at the basis for calling some statements of number theory true, and others false. How much is 12 times 12? Everyone knows it is 144. But how many of the people who give that answer have actually at 







any time in their lives drawn a 12 by 12 rectangle, and then counted the little squares in it? Most people would regard the drawing and counting unnecessary. They would instead offer as proof a few marks on paper, such as are shown below: 

X 12 ------ ------ 



And that would be the "proof". Nearly everyone believes that if you counted the squares, you would get 144 of them; few people feel that outcome is in doubt. The conflict between the two points of view comes into sharper focus when you consider the problem of determining the value 987654321 x 123456789. First of all, it is virtually impossible to construct the appropriate rectangle; and what is worse, even if it were constructed and huge armies of people spent centuries counting the little squares, o a very gullible person would be willing to believe their final answer. It is just too likely that somewhere, somehow, somebody bobbled just a little bit. So is it ever possible to know what the answer is? If you trust the symbolic process which involves manipulating digits according to certain simple rules, yes. That process is presented to children as a device which gets right answer; lost in the shuffle, for many children, are the rhyme reason of that process. The digit-shunting laws for multiplication are based mostly on a few properties of addition and multiplication which are assumed to hold for all numbers. 

The Basic Laws of Arithmetic The kind of assumption I mean is illustrated below. Suppose that you down a few sticks: 

/ // // // / / 

Now you count them. At the same time, somebody else counts them, starting from the other end. Is it clear that the two of you will get the s: answer? The result of a counting process is independent of the way in which it is done. This is really an assumption about what counting i would be senseless to try to prove it, because it is so basic; either you s or you don't-but in the latter case, a proof won't help you a bit. From this kind of assumption, one can get to the commutativity and associativity of addition (i.e., first that b + c = c + b always, and second that b + (c + d) = (b + c) + d always). The same assumption can also you to the commutativity and associativity of multiplication; just think of 







many cubes assembled to form a large rectangular solid. Multiplicative commutativity and associativity are just the assumptions that when you rotate the solid in various ways, the number of cubes will not change. Now these assumptions are not verifiable in all possible cases, because the number of such cases is infinite. We take them for granted; we believe them (if we ever think about them) as deeply as we could believe anything. The amount of money in our pocket will not change as we walk down the street, jostling it up and down; the number of books we have will not change if we pack them up in a box, load them into our car, drive one hundred miles, unload the box, unpack it, and place the books in a new shelf. All of this is part of what we mean by number. There are certain types of people who, as soon as some undeniable fact is written down, find it amusing to show why that "fact" is false after all. I am such a person, and as soon as I had written down the examples above involving sticks, money, and books, I invented situations in which they were wrong. You may have done the same. It goes to show that numbers as abstractions are really quite different from the everyday numbers which we use. People enjoy inventing slogans which violate basic arithmetic but which illustrate "deeper" truths, such as "1 and 1 make 1" (for lovers), or "1 plus 1 plus 1 equals 1" (the Trinity). You can easily pick holes in those slogans, showing why, for instance, using the plus-sign is inappropriate in both cases. But such cases proliferate. Two raindrops running down a windowpane merge; does one plus one make one? A cloud breaks up into two clouds-more evidence for the same? It is not at all easy to draw a sharp line between cases where what is happening could be called "addition", and where some other word is wanted. If you think about the question, you will probably come up with some criterion involving separation of the objects in space, and making sure each one is clearly distinguishable from all the others. But then how could one count ideas? Or the number of gases comprising the atmosphere? Somewhere, if you try to look it up, you can probably find a statement such as, "There are 17 languages in India, and 462 dialects." There is something strange about precise statements like that, when the concepts "language" and "dialect" are themselves fuzzy. 

Ideal Numbers Numbers as realities misbehave. However, there is an ancient and innate sense in people that numbers ought not to misbehave. There is something clean and pure in the abstract notion of number, removed from counting beads, dialects, or clouds; and there ought to be a way of talking about numbers without always having the silliness of reality come in and intrude. The hard-edged rules that govern "ideal" numbers constitute arithmetic, and their more advanced consequences constitute number theory. There is only one relevant question to be asked, in making the transition from numbers as practical things to numbers as formal things. Once you have 







FIGURE 13. Liberation, by M.C. Escher (lithograph, 1955). 







decided to try to capsulize all of number theory in an ideal system, is it really possible to do the job completely? Are numbers so clean and crystalline and regular that their nature can be completely captured in the rules of a formal system? The picture Liberation (Fig. 13), one of Escher's most beautiful, is a marvelous contrast between the formal and the informal, with a fascinating transition region. Are numbers really as free as birds? Do they suffer as much from being crystallized into a rule-obeying system? Is there a magical transition region between numbers in reality and numbers on paper? When I speak of the properties of natural numbers, I don't just mean properties such as the sum of a particular pair of integers. That can be found out by counting, and anybody who has grown up in this century cannot doubt the mechanizability of such processes as counting, adding, multiplying, and so on. I mean the kinds of properties which mathematicians are interested in exploring, questions for which no counting- process is sufficient to provide the answer-not even theoretically sufficient. Let us take a classic example of such a property of natural numbers. The statement is: "There are infinitely many prime numbers." First of all, there is no counting process which will ever be able to confirm, or refute, this assertion. The best we could do would be to count primes for a while and concede that there are "a lot". But no amount of counting alone would ever resolve the question of whether the number of primes is finite or infinite. There could always be more. The statement-and it is called "Euclid's Theorem" (notice the capital "T")-is quite unobvious. It may seem reasonable, or appealing, but it is not obvious. However, mathematicians since Euclid have always called it true. What is the reason? 

Euclid's Proof The reason is that reasoning tells them it is so. Let us follow the reasoning involved. We will look at a variant of Euclid's proof. This proof works by showing that whatever number you pick, there is a prime larger than it. Pick a number-N. Multiply all the positive integers starting with 1 and ending with N; in other words, form the factorial of N, written "N!". What you get is divisible by every number up to N. When you add 1 to N!, the result 

can't be a multiple of 2 (because it leaves 1 over, when you divide by 2); can't be a multiple of 3 (because it leaves I over, when you divide by 3); can't be a multiple of 4 (because it leaves 1 over, when you divide by 4); . . . 







can't be a multiple of N (because it leaves 1 over, when you divide by N); 

In other words, N! + 1, if it is divisible at all (other than by 1 and itself only is divisible by numbers greater than N. So either it is itself prime, or prime divisors are greater than N. But in either case we've shown the must exist a prime above N. The process holds no matter what number is. Whatever N is, there is a prime greater than N. And thus ends the demonstration of the infinitude of the primes. This last step, incidentally, is called generalization, and we will meet again later in a more formal context. It is where we phrase an argument terms of a single number (N), and then point out that N was unspecified and therefore the argument is a general one. Euclid's proof is typical of what constitutes "real mathematics". It simple, compelling, and beautiful. It illustrates that by taking several rash short steps one can get a long way from one's starting point. In our case, t starting points are basic ideas about multiplication and division and forth. The short steps are the steps of reasoning. And though eve individual step of the reasoning seems obvious, the end result is not obvious. We can never check directly whether the statement is true or not; } we believe it, because we believe in reasoning. If you accept reasoning there seems to be no escape route; once you agree to hear Euclid out, you’ll have to agree with his conclusion. That's most fortunate-because it mea that mathematicians will always agree on what statements to label "true and what statements to label "false". This proof exemplifies an orderly thought process. Each statement related to previous ones in an irresistible way. This is why it is called "proof'' rather than just "good evidence". In mathematics the goal always to give an ironclad proof for some unobvious statement. The very fact of the steps being linked together in an ironclad way suggests ti there may be a patterned structure binding these statements together. TI structure can best be exposed by finding a new vocabulary-a stylized vocabulary, consisting of symbols-suitable only for expressing statements about numbers. Then we can look at the proof as it exists in its translated version. It will be a set of statements which are related, line by line, in some detectable way. But the statements, since they're represented by means a small and stylized set of symbols, take on the aspect of patterns. In other words, though when read aloud, they seem to be statements about numb and their properties, still when looked at on paper, they seem to be abstract patterns-and the line-by-line structure of the proof may start to look like slow transformation of patterns according to some few typographical rules. 

Getting Around Infinity Although Euclid's proof is a proof that all numbers have a certain property it avoids treating each of the infinitely many cases separately. It gets around 







it by using phrases like "whatever N is", or "no matter what number N is". We could also phrase-the proof over again, so that it uses the phrase "all N". By knowing the appropriate context and correct ways of using such phrases, we never have to deal with infinitely many statements. We deal with just two or three concepts, such as the word "all"-which, though themselves finite, embody an infinitude; and by using them, we sidestep the apparent problem that there are an infinite number of facts we want to prove. We use the word "all" in a few ways which are defined by the thought processes of reasoning. That is, there are rules which our usage of "all" obeys. We may be unconscious of them, and tend to claim we operate on the basis of the meaning of the word; but that, after all, is only a circumlocution for saying that we are guided by rules which we never make explicit. We have used words all our lives in certain patterns, and instead of calling the patterns "rules", we attribute the courses of our thought processes to the "meanings" of words. That discovery was a crucial recognition in the long path towards the formalization of number theory. If we were to delve into Euclid's proof more and more carefully, we would see that it is composed of many, many small-almost infinitesimal steps. If all those steps were written out line after line, the proof would appear incredibly complicated. To our minds it is clearest when several steps are telescoped together, to form one single sentence. If we tried to look at the proof in slow motion, we would begin to discern individual frames. In other words, the dissection can go only so far, and then we hit the "atomic" nature of reasoning processes. A proof can be broken down into a series of tiny but discontinuous jumps which seem to flow smoothly when perceived from a higher vantage point. In 
\chapter {VIII}
see how incredibly many steps are involved. Perhaps it should not surprise you, though. The operations in Euclid's brain when he invented the proof must have involved millions of neurons (nerve cells), many of which fired several hundred times in a single second. The mere utterance of a sentence involves hundreds of thousands of neurons. If Euclid's thoughts were that complicated, it makes sense for his proof to contain a huge number of steps! (There may be little direct connection between the neural actions in his brain, and a proof in our formal system, but the complexities of the two are comparable. It is as if nature wants the complexity of the proof of the infinitude of primes to be conserved, even when the systems involved are very different from each other.) In Chapters to come, we will lay out a formal system that (1) includes a stylized vocabulary in which all statements about natural numbers can be expressed, and (2) has rules corresponding to all the types of reasoning which seem necessary. A very important question will be whether the rules for symbol manipulation which we have then formulated are really of equal power (as far as number theory is concerned) to our usual mental reasoning abilities-or, more generally, whether it is theoretically possible to attain the level of our thinking abilities, by using some formal system. 







Sonata for Unaccompanied Achilles 

The telephone rings; Achilles picks it up. 


\achilles Hello, this is Achilles.\par

\achilles Oh, hello, Mr. T. How are you?\par

\achilles A torticollis? Oh, I'm sorry to hear it. Do you have any idea what caused it?\par

\achilles How long did you hold it in that position?\par

\achilles Well, no wonder it's stiff, then. What on earth induced you keep your neck\par
twisted that way for so long? 
\achilles Wondrous many of them, eh? What kinds, for example? Achilles: What do you\par
mean, "phantasmagorical beasts"? 

FIGURE 14. Mosaic II, by M. C. Escher (lithograph, 1957). 








\achilles Wasn't it terrifying to see so many of them at the same time? Achilles: A\par
guitar!? Of all things to be in the midst of all those weird creatures. Say, don't you play the guitar? 
\achilles Oh, well, it's all the same to me.\par

\achilles You're right; I wonder why I never noticed that difference between fiddles and\par
guitars before. Speaking of fiddling, how would you like to come over and listen to one of the sonatas for unaccompanied violin by your favorite composer, J. S. Bach? I just bought a marvelous recording of them. I still can't get over the way Bach uses a single violin to create a piece with such interest. 
\achilles A headache too? That's a shame. Perhaps you should just go to bed.\par

\achilles I see. Have you tried counting sheep?\par

\achilles Oh, oh, I see. Yes, I fully know what you mean. Well, if it's THAT distracting,\par
perhaps you'd better tell it to me, and let me try to work on it, too. 
\achilles A word with the letters `A', `D', `A', `C' consecutively inside it ... Hmm ...\par
What about "abracadabra"? 
\achilles True, "ADAC" occurs backwards, not forwards, in that word. Achilles: Hours\par
and hours? It sounds like I'm in for a long puzzle, then. Where did you hear this infernal riddle? 
\achilles You mean he looked like he was meditating on esoteric Buddhist matters, but in\par
reality he was just trying to think up complex word puzzles? 
\achilles Aha!-the snail knew what this fellow was up to. But how did you come to talk\par
to the snail? 
\achilles Say, I once heard a word puzzle a little bit like this one. Do you want to hear it?\par
Or would it just drive you further into distraction? Achilles: I agree-can't do any harm. Here it is: What's a word that begins with the letters "HE" and also ends with "HE"? 
\achilles Very ingenious-but that's almost cheating. It's certainly not what I meant!\par

\achilles Of course you're right-it fulfills the conditions, but it's a sort of "degenerate"\par
solution. There's another solution which I had in mind. Achilles: That's exactly it! How did you come up with it so fast? Achilles: So here's a case where having a headache actually might have helped you, rather than hindering you. Excellent! But I'm still in the dark on your "ADAC" puzzle. 
\achilles Congratulations! Now maybe you'll be able to get to sleep! So tell me, what is\par
the solution? 
\achilles Well, normally I don't like hints, but all right. What's your hint? Achilles: I\par
don't know what you mean by "figure" and "ground" in this case. 
\achilles Certainly I know Mosaic II! I know ALL of Escher's works. After all, he's my\par
favorite artist. In any case, I've got a print of Mosaic II hanging on my wall, in plain view from here. 








\achilles Yes, t see all the black animals.\par

\achilles Yes, I also see how their "negative space" -- what's left out-- defines the white\par
animals. 
\achilles So THAT'S what you mean by "figure" and "ground". But what does that have\par
to do with the "ADAC" puzzle? 
\achilles Oh, this is too tricky for me. I think I'M starting to get a headache\par

\achilles You want to come over now? But I thought--\par

\achilles Very well. Perhaps by then I'll have thought of the right answer to YOUR\par
puzzle, using your figure-ground hint, relating it to MY puzzle 
\achilles I'd love to play them for you.\par

\achilles You've invented a theory about them?\par

\achilles Accompanied by what instrument?\par

\achilles Well, if that's the case, it seems a little strange that he would have written out\par
the harpsichord part, then, and had it published a s well. 
\achilles I see -- sort of an optional feature. One could listen to them either way -- with\par
or without accompaniment. But how would one know what the accompaniment is supposed to sound like? 
\achilles Ah, yes, I guess that it is best, after all, to leave it to the listener’s imagination.\par
And perhaps, as you said, Bach never even had accompaniment in mind at all. Those sonatas seem to work very indeed as they are. 
\achilles Right. Well, I'll see you shortly.\par

\achilles Good-bye, Mr. T.\par









\chapter {III}


Figure and Ground Primes vs. Composites THERE IS A strangeness to the idea that concepts can be captured by simple typographical manipulations. The one concept so far captured is that of addition, and it may not have appeared very strange. But suppose the goal were to create a formal system with theorems of the form Px, the letter `x' standing for a hyphen-string, and where the only such theorems would be ones in which the hyphen-string contained exactly a prime number of hyphens. Thus, P--- would be a theorem, but P---- would not. How could this be done typographically? First, it is important to specify clearly what is meant by typographical operations. The complete repertoire has been presented in the MIU-system and the pq-system, so we really only need to make a list of the kinds of things we have permitted: 

(1) reading and recognizing any of a finite set of symbols; (2) writing down any symbol belonging to that set; (3) copying any of those symbols from one place to another; (4) erasing any of those symbols; (5) checking to see whether one symbol is the same as another; (6) keeping and using a list of previously generated theorems. 

The list is a little redundant, but no matter. What is important is that it clearly involves only trivial abilities, each of them far less than the ability to distinguish primes from nonprimes. How, then, could we compound some of these operations to make a formal system in which primes are distinguished from composite numbers? 

The tq-System A first step might be to try to solve a simpler, but related, problem. We could try to make a system similar to the pq-system, except that it represents multiplication, instead of addition. Let's call it the tq-system, `t' for times'. More specifically, suppose X, Y, and Z are, respectively, the numbers of hyphens in the hyphen-strings x, y, and z. (Notice I am taking special pains to distinguish between a string and the number of hyphens it contains.) Then we wish the string x ty q z to be a theorem if and only if X times Y equals Z. For instance, --t---q----- should be a theorem because 2 times 3 equals 6, but -- t--q--- should not be a theorem. The tq-system can be characterized just about as easily as the pq-system namely, by using just one axiom schema and one rule of inference: 







AXIOM SCHEMA: xt-qx is an axiom, whenever x is a hyphen string. . RULE OF INFERENCE: Suppose that x, y, and z are all hyphen-strings. An suppose that x ty qz is an old theorem. Then, xty-qzx is a ne' theorem. 

Below is the derivation of the theorem --t---q----- : 
\intro (1)       --t-q--         (axiom) \endintro

\intro (2)       --t—q----       (by rule of inference, using line (1) as the old theorem) \endintro

\intro (3)       --t---q ------- (by rule of inference, using line (2) as the old theorem) \endintro


Notice how the middle hyphen-string grows by one hyphen each time the rule of inference is applied; so it is predictable that if you want a theorem with ten hyphens in the middle, you apply the rule of inference nine times in a row. 

Capturing Compositeness Multiplication, a slightly trickier concept than addition, has now bee] "captured" typographically, like the birds in Escher's Liberation. What about primeness? Here's a plan that might seem smart: using the tq-system define a new set of theorems of the form Cx, which characterize compost. numbers, as follows: 

RULE: Suppose x, y, and z are hyphen-strings. If x-ty-qz is a theorem then C z is a theorem. 

This works by saying that Z (the number of hyphens in z) is composite a long as it is the product of two numbers greater than 1-namely, X + (the number of hyphens in x-), and Y + 1 (the number of hyphens in y I am defending this new rule by giving you some "Intelligent mode justifications for it. That is because you are a human being, and want t, know why there is such a rule. If you were operating exclusively in the "Mechanical mode", you would not need any justification, since M-mod. workers just follow the rules mechanically and happily, never questioning; them! Because you work in the I-mode, you will tend to blur in your mind the distinction between strings and their interpretations. You see, things Cal become quite confusing as soon as you perceive "meaning" in the symbol which you are manipulating. You have to fight your own self to keep from thinking that the string'---' is the number 3. The Requirement of Formality, which in Chapter I probably seemed puzzling (because it seemed so obvious), here becomes tricky, and crucial. It is the essential thing which keeps you from mixing up the I-mode with the M-mode; or said another way, it keeps you from mixing up arithmetical facts with typographical theorems. 







Illegally Characterizing Primes It is very tempting to jump from the C-type theorems directly to P-type theorems, by proposing a rule of the following kind: 

PROPOSED RULE: Suppose x is a hyphen-string. If Cx is not a theorem, then Px is a theorem. The fatal flaw here is that checking whether Cx is not a theorem is not an explicitly typographical operation. To know for sure that MU is not a theorem of the MIU-system, you have to go outside of the system ... and so it is with this Proposed Rule. It is a rule which violates the whole idea of formal systems, in that it asks you to operate informally- that is, outside the system. Typographical operation (6) allows you to look into the stockpile of previously found theorems, but this Proposed Rule is asking you to look into a hypothetical "Table of Nontheorems". But in order to generate such a table, you would have to do some reasoning outside the system-reasoning which shows why various strings cannot be generated inside the system. Now it may well be that there is another formal system which can generate the "Table of Nontheorems", by purely typographical means. In fact, our aim is to find just such a system. But the Proposed Rule is not a typographical rule, and must be dropped. This is such an important point that we might dwell on it a bit more. In our C- system (which includes the tq-system and the rule which defines C-type theorems), we have theorems of the form Cx, with `x' standing, as usual, for a hyphen-string. There are also nontheorems of the form Cx. (These are what I mean when I refer to "nontheorems", although of course tt-Cqq and other ill-formed messes are also nontheorems.) The difference is that theorems have a composite number of hyphens, nontheorems have a prime number of hyphens. Now the theorems all have a common "form", that is, originate from a common set of typographical rules. Do all nontheorems also have a common "form", in the same sense? Below is a list of C-type theorems, shown without their derivations. The parenthesized numbers following them simply count the hyphens in them. 

C---- (4) C-------- (6) C----------------(8) C-----------------(9) C--------------------(10) C -------------------- (12) C------------------------ (14) C ------------------------ (15) C----------------------------(16) C----------------------------(18) 



. . . 

I he "holes" in this list are the nontheorems. I o repeat the earlier quest Do the holes also have some "form" in common? Would it be reasonable say that merely by virtue of being the holes in this list, they share a common form? Yes and no. That they share some typographical quality is and able, but whether we want to call it "form" is unclear. The reason hesitating is that the holes are only negatively defined-they are the things that are left out of a list which is positively defined. 

Figure and Ground 

This recalls the famous artistic distinction between figure and ground. When a figure or "positive space" (e.g., a human form, or a letter, or a still life is drawn inside a frame, an unavoidable consequence is that its complementary shape-also called the "ground", or "background", or "negative space"-has also been drawn. In most drawings, however, this fig ground relationship plays little role. The artist is much less interested in ground than in the figure. But sometimes, an artist will take interest in ground as well. There are beautiful alphabets which play with this figure-ground distinction. A message written in such an alphabet is shown below. At fir looks like a collection of somewhat random blobs, but if you step back ways and stare at it for a while, all of a sudden, you will see seven letters appear in this .. . 







FIGURE 15. 



For a similar effect, take a look at my drawing Smoke Signal (Fig. 139). Along these lines, you might consider this puzzle: can you somehow create a drawing containing words in both the figure and the ground? Let us now officially distinguish between two kinds of figures: cursively drawable ones, and recursive ones (by the way, these are my own terms are not in common usage). A cursively drawable figure is one whose ground is merely an accidental by-product of the drawing act. A recursive figure is one whose ground can be seen as a figure in its own right. Usually this is quite deliberate on the part of the artist. The "re" in "recursive" represents the fact that both foreground and background are cursively drawable – the figure is "twice-cursive". Each figure-ground boundary in a recursive figure is a double-edged sword. M. C. Escher was a master at drawing recursive figures-see, for instance, his beautiful recursive drawing of birds (Fig. 16). 







FIGURE 16. Tiling of the plane using birds, by M. C. Escher (from a 1942 notebook). 

Our distinction is not as rigorous as one in mathematics, for who can definitively say that a particular ground is not a figure? Once pointed out, almost any ground has interest of its own. In that sense, every figure is recursive. But that is not what I intended by the term. There is a natural and intuitive notion of recognizable forms. Are both the foreground and background recognizable forms? If so, then the drawing is recursive. If you look at the grounds of most line drawings, you will find them rather unrecognizable. This demonstrates that 

There exist recognizable forms whose negative space is not any recognizable form. 

In more "technical" terminology, this becomes: 

There exist cursively drawable figures which are not recursive. 

Scott Kim's solution to the above puzzle, which I call his "FIGURE-FIGURE Figure", is shown in Figure 17. If you read both black and white, 







FIGURE 17. FIGURE-FIGURE Figure, by Scott E. Kim (1975). 







you will see "FIGURE" everywhere, but "GROUND" nowhere! It is a paragon of recursive figures. In this clever drawing, there are two nonequivalent ways of characterizing the black regions: 

(1) as the negative space to the white regions; (2) as altered copies of the white regions (produced by coloring and shifting each white region). 

(In the special case of the FIGURE-FIGURE Figure, the two characterizations are equivalent-but in most black-and-white pictures, they would not be.) Now in Chapter VIII, when we create our Typographical Number Theory (TNT), it will be our hope that the set of all false statements of number theory can be characterized in two analogous ways: 

(1) as the negative space to the set of all TNT-theorems; (2) as altered copies of the set of all TNT-theorems (produced by negating each TNT-theorem). 

But this hope will be dashed, because: 

(1) inside the set of all nontheorems are found some truths (2) outside the set of all negated theorems are found some falsehoods . You will see why and how this happens, in Chapter XIV. Meanwhile, ponder over a pictorial representation of the situation (Fig. 18). 

Figure and Ground in Music One may also look for figures and grounds in music. One analogue is the distinction between melody and accompaniment-for the melody is always in the forefront of our attention, and the accompaniment is subsidiary, in some sense. Therefore it is surprising when we find, in the lower lines of a piece of music, recognizable melodies. This does not happen too often in post-baroque music. Usually the harmonies are not thought of as foreground. But in baroque music-in Bach above all-the distinct lines, whether high or low or in between, all act as "figures". In this sense, pieces by Bach can be called "recursive". Another figure-ground distinction exists in music: that between on-beat and off- beat. If you count notes in a measure "one-and, two-and, three-and, four-and", most melody-notes will come on numbers, not on "and"'s. But sometimes, a melody will be deliberately pushed onto the "and" 's, for the sheer effect of it. This occurs in several etudes for the piano by Chopin, for instance. It also occurs in Bach-particularly in his Sonatas and Partitas for unaccompanied violin, and his Suites for unaccompanied cello. There, Bach manages to get two or more musical lines going simultaneously. Sometimes he does this by having the solo instrument play "double stops"-two notes at once. Other times, however, he 







FIGURE 18. Considerable visual symbolism is featured in this diagram of the relation between various classes of TNT strings. The biggest box represents the set of all TNT strings The next-biggest box represents the set of all well-formed TNT strings. Within it is found~ set of all sentences of TNT. Now things begin to get interesting. The set of theorems pictured as a tree growing out of a trunk (representing the set of axioms). The tree-symbol chosen because of the recursive growth pattern which it exhibits: new branches (theorems constantly sprouting from old ones. The fingerlike branches probe into the corners of constraining region (the set of truths), yet can never fully occupy it. The boundary beta the set of truths and the set of falsities is meant to suggest a randomly meandering coastline which, no matter how closely you examine it, always has finer levels of structure, an consequently impossible to describe exactly in any finite way. (See B. Mandelbrot's book Fractals.) The reflected tree represents the set of negations of theorems: all of them false yet unable collectively to span the space of false statements. [Drawing by the author.] 

puts one voice on the on-beats, and the other voice on the off-beats, so ear separates them and hears two distinct melodies weaving in and out, - harmonizing with each other. Needless to say, Bach didn't stop at this level of complexity... 

Recursively Enumerable Sets vs. Recursive Sets Now let us carry back the notions of figure and ground to the domain formal systems. In our example, the role of positive space is played by C-type theorems, and the role of negative space is played by strings with a 







prime number of hyphens. So far, the only way we have found to represent prime numbers typographically is as a negative space. Is there, however, some way-I don't care how complicated-of representing the primes as a positive space-that is, as a set of theorems of some formal system? Different people's intuitions give different answers here. I remember quite vividly how puzzled and intrigued I was upon realizing the difference between a positive characterization and a negative characterization. I was quite convinced that not only the primes, but any set of numbers which could be represented negatively, could also be represented positively. The intuition underlying my belief is represented by the question: "How could a figure and its ground not carry exactly the same information?" They seemed to me to embody the same information, just coded in two complementary ways. What seems right to you? It turns out I was right about the primes, but wrong in general. This astonished me, and continues to astonish me even today. It is a fact that: 

There exist formal systems whose negative space (set of nontheorems) is not the positive space (set of theorems) of any formal system. 

This result, it turns out, is of depth equal to Gödel’s Theorem-so it is not surprising that my intuition was upset. I, just like the mathematicians of the early twentieth century, expected the world of formal systems and natural numbers to be more predictable than it is. In more technical terminology, this becomes: 

There exist recursively enumerable sets which are not recursive. 

The phrase recursively enumerable (often abbreviated "r.e.") is the mathematical counterpart to our artistic notion of "cursively drawable"-and recursive is the counterpart of "recursive". For a set of strings to be "r.e." means that it can be generated according to typographical rules-for example, the set of C-type theorems, the set of theorems of the MIU-system-indeed, the set of theorems of any formal system. This could be compared with the conception of a "figure" as "a set of lines which can be generated according to artistic rules" (whatever that might mean!). And a "recursive set" is like a figure whose ground is also a figure-not only is it r.e., but its complement is also r.e. It follows from the above result that: 

There exist formal systems for which there is no typographical decision procedure. 

How does this follow? Very simply. A typographical decision procedure is a method which tells theorems from nontheorems. The existence of such a test allows us to generate all nontheorems systematically, simply by going down a list of all strings and performing the test on them one at a time, discarding ill-formed strings and theorems along the way. This amounts to 







a typographical method for generating the set of nontheorems. But according to the earlier statement (which we here accept on faith), for some systems this is not possible. So we must conclude that typographical decision procedures do not exist for all formal systems. Suppose we found a set F of natural numbers (`F' for `Figure') whi4 we could generate in some formal way-like the composite numbers. Suppose its complement is the set G (for 'Ground')-like the primes. Together F and G make up all the natural numbers, and we know a rule for making all the numbers in set F, but we know no such rule for making all tl numbers in set G. It is important to understand that if the members of were always generated in order of increasing size, then we could always characterize G. The problem is that many r.e. sets are generated I methods which throw in elements in an arbitrary order, so you never know if a number which has been skipped over for a long time will get included you just wait a little longer. We answered no to the artistic question, "Are all figures recursive We have now seen that we must likewise answer no to the analogous question in mathematics: "Are all sets recursive?" With this perspective, 1 us now come back to the elusive word "form". Let us take our figure-set and our ground-set G again. We can agree that all the numbers in set have some common "form"-but can the same be said about numbers in s G? It is a strange question. When we are dealing with an infinite set to sta with-the natural numbers-the holes created by removing some subs may be very hard to define in any explicit way. And so it may be that th< are not connected by any common attribute or "form". In the last analysis it is a matter of taste whether you want to use the word "form"-but just thinking about it is provocative. Perhaps it is best not to define "form", bi to leave it with some intuitive fluidity. Here is a puzzle to think about in connection with the above matter Can you characterize the following set of integers (or its negative space) 

1       3       7      12      18      26      35     45      56      69... 

How is this sequence like the FIGURE-FIGURE Figure? 



Primes as Figure Rather than Ground Finally, what about a formal system for generating primes? How is it don< The trick is to skip right over multiplication, and to go directly to nondivisibility as the thing to represent positively. Here are an axiom schema and rule for producing theorems which represent the notion that one number does not divide (D N D) another number exactly: 

AXIOM SCHEMA: xy D N Dx where x and y are hyphen-strings. 

For example ----D N D--, where x has been replaced by'--'and y by ‘---“. 







RULE: If x D N Dy is a theorem, then so is x D N Dx y. 

If you use the rule twice, you can generate this theorem: 

-----D N D -------------- 

which is interpreted as "5 does not divide 12". But ---D N D------------ is not a theorem. What goes wrong if you try to produce it? Now in order to determine that a given number is prime, we have to build up some knowledge about its nondivisibility properties. In particular, we want to know that it is not divisible by 2 or 3 or 4, etc., all the way up to 1 less than the number itself. But we can't be so vague in formal systems as to say "et cetera". We must spell things out. We would like to have a way of saying, in the language of the system, "the number Z is divisor free up to X", meaning that no number between 2 and X divides Z. This can be done, but there is a trick to it. Think about it if you want. Here is the solution: 

RULE: If --D N D z is a theorem, so is z D F--. 

RULE: If z D Fx is a theorem and also x-D N Dz is a theorem, z D Fx- is a theorem. 

These two rules capture the notion of divisor freeness. All we need to do is to say that primes are numbers which are divisor-free up to 1 less than themselves: 

RULE: If z-DFz is a theorem, then Pz- is a theorem. 

Oh-let's not forget that 2 is prime! 

Axiom: P--. 

And there you have it. The principle of representing primality formally is that there is a test for divisibility which can be done without any backtracking. You march steadily upward, testing first for divisibility by 2, then by 3, and so on. It is this "monotonicity" or unidirectionality-this absence of cross-play between lengthening and shortening, increasing and decreasing-that allows primality to be captured. And it is this potential complexity of formal systems to involve arbitrary amounts of backwards-forwards interference that is responsible for such limitative results as Gödel’s Theorem, Turing's Halting Problem, and the fact that not all recursively enumerable sets are recursive. 







Contracrostipunctus Achilles has come to visit his friend and jogging companion, the Tortoise, at his home 


\achilles Heavens, you certainly have an admirable boomerang collection\par

\tortoise Oh, pshaw. No better than that of any other Tortoise. And now would you like\par
to step into the parlor? 
\achilles Fine. (Walks to the corner of the room.) I see you also have a large collection of\par
records. What sort of music do you enjoy? 
\tortoise Sebastian Bach isn't so bad, in my opinion. But these days, I must say, I am\par
developing more and more of an interest in a rather specialized sort of music. 
\achilles Tell me, what kind of music is that?\par

\tortoise A type of music which you are most unlikely to have heard of. call it "music to\par
break phonographs by". 
\achilles Did you say "to break phonographs by"? That is a curious concept. I can just\par
see you, sledgehammer in hand, whacking on phonograph after another to pieces, to the strains of Beethoven's heroic masterpiece Wellington's Victory. 
\tortoise That's not quite what this music is about. However, you might find its true\par
nature just as intriguing. Perhaps I should give you a brief description of it? 
\achilles Exactly what I was thinking.\par

\tortoise Relatively few people are acquainted with it. It all began whet my friend the\par
Crab-have you met him, by the way?-paid m• a visit. 
\achilles ' twould be a pleasure to make his acquaintance, I'm sure Though I've heard so\par
much about him, I've never met him 
\tortoise Sooner or later I'll get the two of you together. You'd hit it of splendidly.\par
Perhaps we could meet at random in the park on day ... 
\achilles Capital suggestion! I'll be looking forward to it. But you were going to tell me\par
about your weird "music to smash phone graphs by", weren't you? 
\tortoise Oh, yes. Well, you see, the Crab came over to visit one day. You must\par
understand that he's always had a weakness for fang gadgets, and at that time he was quite an aficionado for, of al things, record players. He had just bought his first record player, and being somewhat gullible, believed every word the salesman had told him about it-in particular, that it was capable of reproducing any and all sounds. In short, he was convinced that it was a Perfect phonograph. 








\achilles Naturally, I suppose you disagreed.\par

\tortoise True, but he would hear nothing of my arguments. He staunchly maintained that\par
any sound whatever was reproducible on his machine. Since I couldn't convince him of the contrary, I left it at that. But not long after that, I returned the visit, taking with me a record of a song which I had myself composed. The song was called "I Cannot Be Played on Record Player 1". 
\achilles Rather unusual. Was it a present for the Crab?\par

\tortoise Absolutely. I suggested that we listen to it on his new phonograph, and he was\par
very glad to oblige me. So he put it on. But unfortunately, after only a few notes, the record player began vibrating rather severely, and then with a loud "pop", broke into a large number of fairly small pieces, scattered all about the room. The record was utterly destroyed also, needless to say. 
\achilles Calamitous blow for the poor fellow, I'd say. What was the matter with his\par
record player? 
\tortoise Really, there was nothing the matter, nothing at all. It simply couldn't reproduce\par
the sounds on the record which I had brought him, because they were sounds that would make it vibrate and break. 
\achilles Odd, isn't it? I mean, I thought it was a Perfect phonograph. That's what the\par
salesman had told him, after all. 
\tortoise Surely, Achilles, you don't believe everything that salesmen tell you! Are you\par
as naive as the Crab was? 
\achilles The Crab was naiver by far! I know that salesmen are notorious prevaricators. I\par
wasn't born yesterday! 
\tortoise In that case, maybe you can imagine that this particular salesman had somewhat\par
exaggerated the quality of the Crab's piece of equipment ... perhaps it was indeed less than Perfect, and could not reproduce every possible sound. 
\achilles Perhaps that is an explanation. But there's no explanation for the amazing\par
coincidence that your record had those very sounds on it ... 
\tortoise Unless they got put there deliberately. You see, before returning the Crab's\par
visit, I went to the store where the Crab had bought his machine, and inquired as to the make. Having ascertained that, I sent off to the manufacturers for a description of its design. After receiving that by return mail, I analyzed the entire construction of the phonograph and discovered a certain set of sounds which, if they were produced anywhere in the vicinity, would set the device to shaking and eventually to falling apart. 
\achilles Nasty fellow! You needn't spell out for me the last details: that you recorded\par
those sounds yourself, and offered the dastardly item as a gift ... 








\tortoise Clever devil! You jumped ahead of the story! But that wasn't t end of the\par
adventure, by any means, for the Crab did r believe that his record player was at fault. He was quite stubborn. So he went out and bought a new record player, this o even more expensive, and this time the salesman promised give him double his money back in case the Crab found a soul which it could not reproduce exactly. So the Crab told r excitedly about his new model, and I promised to come over and see it. 
\achilles Tell me if I'm wrong-I bet that before you did so, you on again wrote the\par
manufacturer, and composed and recorded new song called "I Cannot Be Played on Record Player based on the construction of the new model. 
\tortoise Utterly brilliant deduction, Achilles. You've quite got the spirit.\par

\achilles So what happened this time?\par

\tortoise As you might expect, precisely the same thing. The phonograph fell into\par
innumerable pieces, and the record was shattered. Achilles: Consequently, the Crab finally became convinced that there could be no such thing as a Perfect record player. 
\tortoise Rather surprisingly, that's not quite what happened. He was sure that the next\par
model up would fill the bill, and having twice the money, h e-- 
\achilles Oho-I have an idea! He could have easily outwitted you, I obtaining a LOW-\par
fidelity phonograph-one that was not capable of reproducing the sounds which would destroy it. In that way, he would avoid your trick. 
\tortoise Surely, but that would defeat the-original purpose-namely, to have a\par
phonograph which could reproduce any sound whatsoever, even its own self- breaking sound, which is of coup impossible. 
\achilles That's true. I see the dilemma now. If any record player-si\par
Record Player X-is sufficiently high-fidelity, then when attempts to play the song "I Cannot Be Played on Record Player X", it will create just those vibrations which will cause to break. .. So it fails to be Perfect. And yet, the only way to g, around that trickery, namely for Record Player X to be c lower fidelity, even more directly ensures that it is not Perfect It seems that every record player is vulnerable to one or the other of these frailties, and hence all record players are defective. 
\tortoise I don't see why you call them "defective". It is simply an inherent fact about\par
record players that they can't do all that you might wish them to be able to do. But if there is a defect anywhere, is not in THEM, but in your expectations of what they should b able to do! And the Crab was just full of such unrealistic expectations. 








\achilles Compassion for the Crab overwhelms me. High fidelity or low fidelity, he loses\par
either way. 
\tortoise And so, our little game went on like *_his for a few more rounds, and\par
eventually our friend tried to become very smart. He got wind of the principle upon which I was basing my own records, and decided to try to outfox me. He wrote to the phonograph makers, and described a device of his own invention, which they built to specification. He called it "Record Player Omega". It was considerably more sophisticated than an ordinary record player. 
\achilles Let me guess how: Did it have no of cotton? Or\par

\tortoise Let me tell you, instead. That will save some time. In the first place, Record\par
Player Omega incorporated a television camera whose purpose it was to scan any record before playing it. This camera was hooked up to a small built-in computer, which would determine exactly the nature of the sounds, by looking at the groove- patterns. 
\achilles Yes, so far so good. But what could Record Player Omega do with this\par
information? 
\tortoise By elaborate calculations, its little computer figured out what effects the sounds\par
would have upon its phonograph. If it deduced that the sounds were such that they would cause the machine in its present configuration to break, then it did something very clever. Old Omega contained a device which could disassemble large parts of its phonograph subunit, and rebuild them in new ways, so that it could, in effect, change its own structure. If the sounds were "dangerous", a new configuration was chosen, one to which the sounds would pose no threat, and this new configuration would then be built by the rebuilding subunit, under direction of the little computer. Only after this rebuilding operation would Record Player Omega attempt to play the record. 
\achilles Aha! That must have spelled the end of your tricks. I bet you were a little\par
disappointed. 
\tortoise Curious that you should think so ... I don't suppose that you know Godel's\par
Incompleteness Theorem backwards and forwards, do you? 
\achilles Know WHOSE Theorem backwards and forwards? I've\par
heard of anything that sounds like that. I'm sure it's fascinating, but I'd rather hear more about "music to break records by". It's an amusing little story. Actually, I guess I can fill in the end. Obviously, there was no point in going on, and so you sheepishly admitted defeat, and that was that. Isn't that exactly it? 
\tortoise What! It's almost midnight! I'm afraid it's my bedtime. I'd love to talk some\par
more, but really I am growing quite sleepy. 








\achilles As am 1. Well, 1 u be on my way. (As he reaches the door, he suddenly stops,\par
and turns around.) Oh, how silly of me! I almost forgo brought you a little present. Here. (Hands the Tortoise a small neatly wrapped package.) 
\tortoise Really, you shouldn't have! Why, thank you very much indeed think I'll open it\par
now. (Eagerly tears open the package, and ins discovers a glass goblet.) Oh, what an exquisite goblet! Did y know that I am quite an aficionado for, of all things, gl goblets? 
\achilles Didn't have the foggiest. What an agreeable coincidence!\par

\tortoise Say, if you can keep a secret, I'll let you in on something: I trying to find a\par
Perfect goblet: one having no defects of a sort in its shape. Wouldn't it be something if this goblet-h call it "G"-were the one? Tell me, where did you come across Goblet G? 
\achilles Sorry, but that's MY little secret. But you might like to know w its maker is.\par

\tortoise Pray tell, who is it?\par

\achilles Ever hear of the famous glassblower Johann Sebastian Bach? Well, he wasn't\par
exactly famous for glassblowing-but he dabbled at the art as a hobby, though hardly a soul knows it-a: this goblet is the last piece he blew. 
\tortoise Literally his last one? My gracious. If it truly was made by Bach its value is\par
inestimable. But how are you sure of its maker 
\achilles Look at the inscription on the inside-do you see where tletters `B', `A', `C', `H'\par
have been etched? 
\tortoise Sure enough! What an extraordinary thing. (Gently sets Goblet G down on a\par
shelf.) By the way, did you know that each of the four letters in\Bach's name is the name of a musical note? 
\achilles ' tisn't possible, is it? After all, musical notes only go from ‘A’ through `G'.\par

\tortoise Just so; in most countries, that's the case. But in Germany, Bach’s own\par
homeland, the convention has always been similar, except that what we call `B', they call `H', and what we call `B-flat', they call `B'. For instance, we talk about Bach's "Mass in B Minor whereas they talk about his "H-moll Messe". Is that clear? 
\achilles ... hmm ... I guess so. It's a little confusing: H is B, and B B-flat. I suppose\par
his name actually constitutes a melody, then 
\tortoise Strange but true. In fact, he worked that melody subtly into or of his most\par
elaborate musical pieces-namely, the final Contrapunctus in his Art of the Fugue. It was the last fugue Bach ever wrote. When I heard it for the first time, I had no idea how would end. Suddenly, without warning, it broke off. And the ... dead silence. I realized immediately that was where Bach died. It is an indescribably sad moment, and the effect it had o me was-shattering. In any case, B-A-C-H is the last theme c that fugue. It is hidden inside the piece. Bach didn't point it out 







FIGURE 19. The last page of Bach's Art of the Fugue. In the original manuscript, in the handwriting of Bach's son Carl Philipp Emanuel, is written: "N.B. In the course of this fugue, at the point where the name B.A.C.H. was brought in as countersubject, the composer died." (B-A-C-H in box.) I have let this final page of Bach's last fugue serve as an epitaph. [Music Printed by Donald Byrd's program "SMUT", developed at Indiana University] 







Explicitly, but if you know about it, you can find it without much trouble. Ah, me-there are so many clever ways of hiding things in music .. . 
\achilles . . or in poems. Poets used to do very similar things, you know (though it's\par
rather out of style these days). For instance, Lewis Carroll often hid words and names in the first letters (or characters) of the successive lines in poems he wrote. Poems which conceal messages that way are called "acrostics". 
\tortoise Bach, too, occasionally wrote acrostics, which isn't surprising. After all,\par
counterpoint and acrostics, with their levels of hidden meaning, have quite a bit in common. Most acrostics, however, have only one hidden level-but there is no reason that one couldn't make a double-decker-an acrostic on top of an acrostic. Or one could make a "contracrostic"-where the initial letters, taken in reverse order, form a message. Heavens! There's no end to the possibilities inherent in the form. Moreover, it's not limited to poets; anyone could write acrostics-even a dialogician. 
\achilles A dial-a-logician? That's a new one on me.\par

\tortoise Correction: I said "dialogician", by which I meant a writer of dialogues. Hmm\par
... something just occurred to me. In the unlikely event that a dialogician should write a contrapuntal acrostic in homage to J. S. Bach, do you suppose it would be more proper for him to acrostically embed his OWN name-or that of Bach? Oh, well, why worry about such frivolous matters? Anybody who wanted to write such a piece could make up his own mind. Now getting back to Bach's melodic name, did you know that the melody B-A-C-H, if played upside down and backwards, is exactly the same as the original? 
\achilles How can anything be played upside down? Backwards, I can see-you get H-C-\par
A-B-but upside down? You must be pulling my leg. 
\tortoise ' pon my word, you're quite a skeptic, aren't you? Well, I guess I'll have to give\par
you a demonstration. Let me just go and fetch my fiddle- (Walks into the next room, and returns in a jiffy with an ancient-looking violin.) -and play it for you forwards and backwards and every which way. Let's see, now ... (Places his copy of the Art of the Fugue on his music stand and opens it to the last page.) ... here's the last Contrapunctus, and here's the last theme ... 

The Tortoise begins to play: B-A-C- - but as he bows the final H, suddenly, without warning, a shattering sound rudely interrupts his performance. Both he and Achilles spin around, just in time to catch a glimpse of myriad fragments of glass tinkling to the floor from the shelf where Goblet G had stood, only moments before. And then ... dead silence. 








\chapter {IV}


Consistency, Completeness, and Geometry Implicit and Explicit Meaning IN CHAPTER II, we saw how meaning-at least in the relatively simple context of formal systems-arises when there is an isomorphism between rule-governed symbols, and things in the real world. The more complex the isomorphism, in general, the more "equipment"- both hardware and software-is required to extract the meaning from the symbols. If an isomorphism is very simple (or very familiar), we are tempted to say that the meaning which it allows us to see is explicit. We see the meaning without seeing the isomorphism. The most blatant example is human language, where people often attribute meaning to words in themselves, without being in the slightest aware of the very complex "isomorphism" that imbues them with meanings. This is an easy enough error to make. It attributes all the meaning to the object (the word), rather than to the link between that object and the real world. You might compare it to the naive belief that noise is a necessary side effect of any collision of two objects. This is a false belief; if two objects collide in a vacuum, there will be no noise at all. Here again, the error stems from attributing the noise exclusively to the collision, and not recognizing the role of the medium, which carries it from the objects to the ear. Above, I used the word "isomorphism" in quotes to indicate that it must be taken with a grain of salt. The symbolic processes which underlie the understanding of human language are so much more complex than the symbolic processes in typical formal systems, that, if we want to continue thinking of meaning as mediated by isomorphisms, we shall have to adopt a far more flexible conception of what isomorphisms can be than we have up till now. In my opinion, in fact, the key element in answering the question "What is consciousness?" will be the unraveling of the nature of the "isomorphism" which underlies meaning. 

Explicit Meaning of the Contracrostipunctus All this is by way of preparation for a discussion of the Contracrostipunctus-a study in levels of meaning. The Dialogue has both explicit and implicit meanings. Its most explicit meaning is simply the story 







Which was related. This “explicit meaning is, strictly speaking extremely implicit, in the sense that the brain processes required to understand the events in the story, given only the black marks on paper, are incredibly complex. Nevertheless, we shall consider the events in the story to be the explicit meaning of the Dialogue, and assume that every reader of English uses more or less the same "isomorphism" in sucking that meaning from the marks on the paper. 

Even so, I'd like to be a little more explicit about the explicit meaning of the story. First I'll talk about the record players and the records. The main point is that there are two levels of meaning for the grooves in the records. Level One is that of music. Now what is "music"-a sequence of vibrations in the air, or a succession of emotional responses in a brain? It is both. But before there can be emotional responses, there have to be vibrations. Now the vibrations get "pulled" out of the grooves by a record player, a relatively straightforward device; in fact you can do it with a pin, just pulling it down the grooves. After this stage, the ear converts the vibrations into firings of auditory neurons in the brain. Then ensue a number of stages in the brain, which gradually transform the linear sequence of vibrations into a complex pattern of interacting emotional responses-far too complex for us to go into here, much though I would like to. Let us therefore content ourselves with thinking of the sounds in the air as the "Level One" meaning of the grooves. What is the Level Two meaning of the grooves? It is the sequence of vibrations induced in the record player. This meaning can only arise after the Level One meaning has been pulled out of the grooves, since the vibrations in the air cause the vibrations in the phonograph. Therefore, the Level Two meaning depends upon a chain of two isomorphisms: 

(1) Isomorphism between arbitrary groove patterns and air vibrations; (2) Isomorphism between graph vibrations. arbitrary air vibrations and phonograph vibrations 

This chain of two isomorphisms is depicted in Figure 20. Notice that isomorphism I is the one which gives rise to the Level One meaning. The Level Two meaning is more implicit than the Level One meaning, because it is mediated by the chain of two isomorphisms. It is the Level Two meaning which "backfires", causing the record player to break apart. What is of interest is that the production of the Level One meaning forces the production of the Level Two meaning simultaneously-there is no way to have Level One without Level Two. So it was the implicit meaning of the record which turned back on it, and destroyed it. Similar comments apply to the goblet. One difference is that the mapping from letters of the alphabet to musical notes is one more level of isomorphism, which we could call "transcription". That is followed by "translation"-conversion of musical notes into musical sounds. Thereafter, the vibrations act back on the goblet just as they did on the escalating series of phonographs. 







FIGURE 20. Visual rendition of the principle underlying Gödel’s Theorem: two back-to- back mappings which have an unexpected boomeranging effect. The first is from groove patterns to sounds, carried out by a phonograph. The second-familiar, but usually ignored -- is from sounds to vibrations of the phonograph. Note that the second mapping exists independently of the first one, for any sound in the vicinity, not just ones produced by the phonograph itself, will cause such vibrations. The paraphrase of Gödel’s Theorem says that for any record player, there are records which it cannot play because they will cause its indirect self-destruction. [Drawing by the author. 





Implicit Meanings of the Contracrostipunctus What about implicit meanings of the Dialogue? (Yes, it has more than one of these.) The simplest of these has already been pointed out in the paragraphs above-namely, that the events in the two halves of the dialogue are roughly isomorphic to each other: the phonograph becomes a violin, the Tortoise becomes Achilles, the Crab becomes the Tortoise, the grooves become the etched autograph, etc. Once you notice this simple isomorphism, you can go a little further. Observe that in the first half of the story, the Tortoise is the perpetrator of all the mischief, while in the second half, he is the victim. What do you know, but his own method has turned around and backfired on him! Reminiscent of the backfiring of the records' muusic-or the goblet's inscription-or perhaps of the Tortoise's boomerang collection? Yes, indeed. The story is about backfiring on two levels, as follows ... 

Level One: Goblets and records which backfire; Level Two: The Tortoise's devilish method of exploiting implicit meaning to cause backfires-which backfires. 

Therefore we can even make an isomorphism between the two levels of the story, in which we equate the way in which the records and goblet boomerang back to destroy themselves, with the way in which the Tortoise's own fiendish method boomerangs back to get him in the end. Seen this 







way, the story itself is an example of the backfirings which it discusses. So we can think of the Contracrostipunctus as referring to itself indirectly that its own structure is isomorphic to the events it portrays. (Exactly goblet and records refer implicitly to themselves via the back-to-back morphisms of playing and vibration-causing.) One may read the Dialogue without perceiving this fact, of course-but it is there all the time. 

Mapping Between the Contracrostipunctus and Gödel’s Theorem Now you may feel a little dizzy-but the best is yet to come. (Actually, levels of implicit meaning will not even be discussed here-they will 1 for you to ferret out.) The deepest reason for writing this Dialogue illustrate Gödel’s Theorem, which, as I said in the Introduction, heavily on two different levels of meaning of statements of number t1 Each of the two halves of the Contracrostipunctus is an "isomorphic co Gödel’s Theorem. Because this mapping is the central idea of the Dialogue and is rather elaborate, I have carefully charted it out below. 

Phonograph <= =>axiomatic system for number theory 

low-fidelity phonograph <= =>"weak" axiomatic system 

high-fidelity phonograph <= =>"strong" axiomatic system 

"Perfect" phonograph" <= => complete system for number theory' 

Blueprint" of phonograph <= => axioms and rules of formal system 

record <= => string of the formal system 

playable record<= => theorem of the axiomatic system 

unplayable record <= =>nontheorem of the axiomatic system 

sound <= =>true statement of number theory 

reproducible sound <= => 'interpreted theorem of the system 

unreproducible sound <= => true statement which isn't a theorem: 

song title <= =>implicit meaning of Gödel’s string: 

"I Cannot Be Played            "I Cannot Be Derived on Record Player X"            in Formal System X" 

This is not the full extent of the isomorphism between Gödel’s theorem and the Contracrostipunctus, but it is the core of it. You need not if you don't fully grasp Gödel’s Theorem by now-there are still Chapters to go before we reach it! Nevertheless, having read this Dialogue you have already tasted some of the flavor of Gödel’s Theorem without necessarily being aware of it. I now leave you to look for any other types of implicit meaning in the Contracrostipunctus. "Quaerendo invenietis!" 



The Art of the Fugue A few words on the Art of the Fugue ... Composed in the last year of Bach's life, it is a collection of eighteen fugues all based on one theme. Apparently, writing the Musical Offering was an inspiration to Bach. He decided to compose another set of fugues on a much simpler theme, to demonstrate the full range of possibilities inherent in the form. In the Art of the Fugue, Bach uses a very simple theme in the most complex possible ways. The whole work is in a single key. Most of the fugues have four voices, and they gradually increase in complexity and depth of expression. Toward the end, they soar to such heights of intricacy that one suspects he can no longer maintain them. Yet he does . . . until the last Contrapunctus. The circumstances which caused the break-off of the Art of the Fugue (which is to say, of Bach's life) are these: his eyesight having troubled him for years, Bach wished to have an operation. It was done; however, it came out quite poorly, and as a consequence, he lost his sight for the better part of the last year of his life. This did not keep him from vigorous work on his monumental project, however. His aim was to construct a complete exposition of fugal writing, and usage of multiple themes was one important facet of it. In what he planned as the next-to-last fugue, he inserted his own name coded into notes as the third theme. However, upon this very act, his health became so precarious that he was forced to abandon work on his cherished project. In his illness, he managed to dictate to his son-in-law a final chorale prelude, of which Bach's biographer Forkel wrote, "The expression of pious resignation and devotion in it has always affected me whenever I have played it; so that I can hardly say which I would rather miss-this Chorale, or the end of the last fugue." One day, without warning, Bach regained his vision. But a few hours later, he suffered a stroke; and ten days later, he died, leaving it for others to speculate on the incompleteness of the Art of the Fugue. Could it have been caused by Bach's attainment of self-reference? 

Problems Caused by Gödel’s Result The Tortoise says that no sufficiently powerful record player can be perfect, in the sense of being able to reproduce every possible sound from a record. Godel says that no sufficiently powerful formal system can be perfect, in the sense of reproducing every single true statement as a theorem. But as the Tortoise pointed out with respect to phonographs, this fact only seems like a defect if you have unrealistic expectations of what formal systems should be able to do. Nevertheless, mathematicians began this century with just such unrealistic expectations, thinking that axiomatic reasoning was the cure to all ills. They found out otherwise in 1931. The fact that truth transcends theoremhood, in any given formal system, is called "incompleteness" of that system. A most puzzling fact about Gödel’s method of proof is that he uses 







reasoning methods which seemingly cannot be "encapsulated"-they re being incorporated into any formal system. Thus, at first sight, it seems that Gödel has unearthed a hitherto unknown, but deeply significant, difference between human reasoning and mechanical reasoning. This mysterious discrepancy in the power of living and nonliving systems is mirrored in the discrepancy between the notion of truth, and that of theoremhood or at least that is a "romantic" way to view the situation. 

The Modified pq-System and Inconsistency In order to see the situation more realistically, it is necessary to see in, depth why and how meaning is mediated, in formal systems, by isomorphisms. And I believe that this leads to a more romantic way to view i situation. So we now will proceed to investigate some further aspects of 1 relation between meaning and form. Our first step is to make a new formal system by modifying our old friend, the pq-system, very slightly. We a one more axiom schema (retaining the original one, as well as the sin rule of inference): 

Axiom SCHEMA II: If x is a hyphen-string, then xp-qx is an axiom. 

Clearly, then, --p-q-- is a theorem in the new system, and so --p--q---. And yet, their interpretations are, respectively, "2 plus; equals 2", and "2 plus 2 equals 3". It can be seen that our new system contain a lot of false statements (if you consider strings to be statement Thus, our new system is inconsistent with the external world. As if this weren't bad enough, we also have internal problems with < new system, since it contains statements which disagree with one another such as -p-q-- (an old axiom) and -p-q- (a new axiom). So our system is inconsistent in a second sense: internally. Would, therefore, the only reasonable thing to do at this point be drop the new system entirely? Hardly. I have deliberately presented the "inconsistencies" in a wool- pulling manner: that is, I have tried to press fuzzy-headed arguments as strongly as possible, with the purpose of n leading. In fact, you may well have detected the fallacies in what I hi said. The crucial fallacy came when I unquestioningly adopted the very same interpreting words for the new system as I had for the old of Remember that there was only one reason for adopting those words in I last Chapter, and that reason was that the symbols acted isomorphically to concepts which they were matched with, by the interpretation. But when y modify the rules governing the system, you are bound to damage t isomorphism. It just cannot be helped. Thus all the problems which we lamented over in preceding paragraphs were bogus problems; they can made to vanish in no time, by suitably reinterpreting some of the symbols of system. Notice that I said "some"; not necessarily all symbols will have to mapped onto new notions. Some may very well retain their "meaning while others change. 







Suppose, for instance, that we reinterpret just the symbol q, leaving all the others constant; in particular, interpret q by the phrase "is greater than or equal to". Now, our "contradictory" theorems -p-q-and -p-q--come out harmlessly as: "1 plus 1 is greater than or equal to 1", and "1 plus 1 is greater than or equal to 2". We have simultaneously gotten rid of (1) the inconsistency with the external world, and (2) the internal inconsistency. And our new interpretation is a meaningful interpretation; of course the original one is meaningless. That is, it is meaningless for the new system; for the original pq-system, it is fine. But it now seems as pointless and arbitrary to apply it to the new pq-system as it was to apply the "horse-apple-happy" interpretation to the old pq-system. 

The History of Euclidean Geometry Although I have tried to catch you off guard and surprise you a little, this lesson about how to interpret symbols by words may not seem terribly difficult once you have the hang of it. In fact, it is not. And yet it is one of the deepest lessons of all of nineteenth century mathematics! It all begins with Euclid, who, around 300 B.C., compiled and systematized all of what was known about plane and solid geometry in his day. The resulting work, Euclid's Elements, was so solid that it was virtually a bible of geometry for over two thousand years-one of the most enduring works of all time. Why was this so? The principal reason was that Euclid was the founder of rigor in mathematics. The Elements began with very simple concepts, definitions, and so forth, and gradually built up a vast body of results organized in such a way that any given result depended only on foregoing results. Thus, there was a definite plan to the work, an architecture which made it strong and sturdy. Nevertheless, the architecture was of a different type from that of, say, a skyscraper. (See Fig. 21.) In the latter, that it is standing is proof enough that its structural elements are holding it up. But in a book on geometry, when each proposition is claimed to follow logically from earlier propositions, there will be no visible crash if one of the proofs is invalid. The girders and struts are not physical, but abstract. In fact, in Euclid's Elements, the stuff out of which proofs were constructed was human language-that elusive, tricky medium of communication with so many hidden pitfalls. What, then, of the architectural strength of the Elements? Is it certain that it is held up by solid structural elements, or could it have structural weaknesses? Every word which we use has a meaning to us, which guides us in our use of it. The more common the word, the more associations we have with it, and the more deeply rooted is its meaning. Therefore, when someone gives a definition for a common word in the hopes that we will abide by that 







FIGURE 21. Tower of Babel, by M. C. Escher (woodcut, 1928). 







definition, it is a foregone conclusion that we will not do so but will instead be guided, largely unconsciously, by what our minds find in their associative stores. I mention this because it is the sort of problem which Euclid created in his Elements, by attempting to give definitions of ordinary, common words such as "point", "straight line", "circle", and so forth. How can you define something of which everyone already has a clear concept? The only way is if you can make it clear that your word is supposed to be a technical term, and is not to be confused with the everyday word with the same spelling. You have to stress that the connection with the everyday word is only suggestive. Well, Euclid did not do this, because he felt that the points and lines of his Elements were indeed the points and lines of the real world. So by not making sure that all associations were dispelled, Euclid was inviting readers to let their powers of association run free ... This sounds almost anarchic, and is a little unfair to Euclid. He did set down axioms, or postulates, which were supposed to be used in the proofs of propositions. In fact, nothing other than those axioms and postulates was supposed to be used. But this is where he slipped up, for an inevitable consequence of his using ordinary words was that some of the images conjured up by those words crept into the proofs which he created. However, if you read proofs in the Elements, do not by any means expect to find glaring "jumps" in the reasoning. On the contrary, they are very subtle, for Euclid was a penetrating thinker, and would not have made any simpleminded errors. Nonetheless, gaps are there, creating slight imperfections in a classic work. But this is not to be complained about. One should merely gain an appreciation for the difference between absolute rigor and relative rigor. In the long run, Euclid's lack of absolute rigor was the cause of some of the most fertile path-breaking in mathematics, over two thousand years after he wrote his work. Euclid gave five postulates to be used as the "ground story" of the infinite skyscraper of geometry, of which his Elements constituted only the first several hundred stories. The first four postulates are rather terse and elegant: 

(1) A straight line segment can be drawn joining any two points. 

(2) Any straight line segment can be extended indefinitely in a straight line. 

(3) Given any straight line segment, a circle can be drawn having the segment as radius and one end point as center. 

(4) All right angles are congruent. 

The fifth, however, did not share their grace: 

(5) If two lines are drawn which intersect a third in such a way that the sum of the inner angles on one side is less than two right angles, then the two lines inevitably must intersect each other on that side if extended far enough 







Though he never explicitly said so, Euclid considered this postulate to be somehow inferior to the others, since he managed to avoid using it in t proofs of the first twenty- eight propositions. Thus, the first twenty-eight propositions belong to what might be called "four-postulate geometry" that part of geometry which can be derived on the basis of the first to postulates of the Elements, without the help of the fifth postulate. (It is al often called absolute geometry.) Certainly Euclid would have found it 1 preferable to prove this ugly duckling, rather than to have to assume it. B he found no proof, and therefore adopted it. But the disciples of Euclid were no happier about having to assume this fifth postulate. Over the centuries, untold numbers of people ga untold years of their lives in attempting to prove that the fifth postulate s itself part of four-postulate geometry. By 1763, at least twenty-eight deficient proofs had been published-all erroneous! (They were all criticized the dissertation of one G. S. Klugel.) All of these erroneous proofs involve a confusion between everyday intuition and strictly formal properties. It safe to say that today, hardly any of these "proofs" holds any mathematic or historical interest-but there are certain exceptions. 

The Many Faces of Noneuclid Girolamo Saccheri (1667-1733) lived around Bach's time. He had t ambition to free Euclid of every flaw. Based on some earlier work he h; done in logic, he decided to try a novel approach to the proof of the famous fifth: suppose you assume its opposite; then work with that as your fif postulate ... Surely after a while you will create a contradiction. Since i mathematical system can support a contradiction, you will have shown t unsoundness of your own fifth postulate, and therefore the soundness Euclid's fifth postulate. We need not go into details here. Suffice it to s that with great skill, Saccheri worked out proposition after proposition "Saccherian geometry" and eventually became tired of it. At one point, decided he had reached a proposition which was "repugnant to the nature of the straight line". That was what he had been hoping for-to his mind was the long-sought contradiction. At that point, he published his work under the title Euclid Freed of Every Flaw, and then expired. But in so doing, he robbed himself of much posthumous glory, sir he had unwittingly discovered what came later to be known as "hyperbolic geometry". Fifty years after Saccheri, J. H. Lambert repeated the "near miss", this time coming even closer, if possible. Finally, forty years after Lambert, and ninety years after Saccheri, non-Euclidean geometry was recognized for what it was-an authentic new brand of geometry, a bifurcation the hitherto single stream of mathematics. In 1823, non- Euclidean geometry was discovered simultaneously, in one of those inexplicable coincidences, by a Hungarian mathematician, Janos (or Johann) Bolyai, age twenty-one, and a Russian mathematician, Nikolay Lobachevskiy, ag thirty. And, ironically, in that same year, the great French mathematician 







Adrien-Marie Legendre came up with what he was sure was a proof of Euclid's fifth postulate, very much along the lines of Saccheri. Incidentally, Bolyai's father, Farkas (or Wolfgang) Bolyai, a close friend of the great Gauss, invested much effort in trying to prove Euclid's fifth postulate. In a letter to his son Janos, he tried to dissuade him from thinking about such matters: 

You must not attempt this approach to parallels. I know this way to its very end. I have traversed this bottomless night, which extinguished all light and joy of my life. I entreat you, leave the science of parallels alone.... I thought I would sacrifice myself for the sake of the truth. I was ready to become a martyr who would remove the flaw from geometry and return it purified to mankind. I accomplished monstrous, enormous labors; my creations are far better than those of others and yet I have not achieved complete satisfaction. For here it is true that si paullum a summo discessit, vergit ad imum. I turned back when I saw that no man can reach the bottom of this night. I turned back unconsoled, pitying myself and all mankind.... I have traveled past all reefs of this infernal Dead Sea and have always come back with broken mast and torn sail. The ruin of my disposition and my fall date back to this time. I thoughtlessly risked my life and happiness sut Caesar aut nihil.' 

But later, when convinced his son really "had something", he urged him to publish it, anticipating correctly the simultaneity which is so frequent in scientific discovery: 

When the time is ripe for certain things, these things appear in different places in the manner of violets coming to light in early spring. 

How true this was in the case of non-Euclidean geometry! In Germany, Gauss himself and a few others had more or less independently hit upon non-Euclidean ideas. These included a lawyer, F. K. Schweikart, who in 1818 sent a page describing a new "astral" geometry to Gauss; Schweikart's nephew, F. A. Taurinus, who did non-Euclidean trigonometry; and F. L. Wachter, a student of Gauss, who died in 1817, aged twenty-five, having found several deep results in non-Euclidean geometry. The clue to non-Euclidean geometry was "thinking straight" about the propositions which emerge in geometries like Saccheri's and Lambert's. The Saccherian propositions are only "repugnant to the nature of the straight line" if you cannot free yourself of preconceived notions of what "straight line" must mean. If, however, you can divest yourself of those preconceived images, and merely let a "straight line" be something which satisfies the new propositions, then you have achieved a radically new viewpoint. 

Undefined Terms This should begin to sound familiar. In particular, it harks back to the pq-system, and its variant, in which the symbols acquired passive meanings by virtue of their roles in theorems. The symbol q is especially interesting, 







since its "meaning" changed when a new axiom schema was added. In the very same way, one can let the meanings of "point", "line", and so on I determined by the set of theorems (or propositions) in which they occur. This was th great realization of the discoverers of non-Euclidean geometry. The found different sorts of non-Euclidean geometries by denying Euclid's fifth postulate in different ways and following out the consequences. Strict] speaking, they (and Saccheri) did not deny the fifth postulate directly, but rather, they denied an equivalent postulate, called the parallel postulate, which runs as follows: 

Given any straight line, and a point not on it, there exists one, and only one, straight line which passes through that point and never intersects the first line, no matter how far they are extended. 

The second straight line is then said to be parallel to the first. If you assert that no such line exists, then you reach elliptical geometry; if you assert that, at east two such lines exist, you reach hyperbolic geometry. Incidentally, tf reason that such variations are still called "geometries" is that the cot element-absolute, or four-postulate, geometry-is embedded in them. is the presence of this minimal core which makes it sensible to think of the] as describing properties of some sort of geometrical space, even if the spa( is not as intuitive as ordinary space. Actually, elliptical geometry is easily visualized. All "points", "lines and so forth are to be parts of the surface of an ordinary sphere. Let t write "POINT" when the technical term is meant, and "point" when t1 everyday sense is desired. Then, we can say that a POINT consists of a pa of diametrically opposed points of the sphere's surface. A LINE is a great circle on the sphere (a circle which, like the equator, has its center at tI center of the sphere). Under these interpretations, the propositions ( elliptical geometry, though they contain words like "POINT" and "LINE speak of the goings-on on a sphere, not a plane. Notice that two LINT always intersect in exactly two antipodal points of the sphere's surface that is, in exactly one single POINT! And just as two LINES determine POINT, so two POINTS determine a LINE. By treating words such as "POINT" and "LINE" as if they had only tt meaning instilled in them by the propositions in which they occur, we take step towards complete formalization of geometry. This semiformal version still uses a lot of words in English with their usual meanings (words such "the", ` if ", "and", "join", "have"), although the everyday meaning has bee drained out of special words like "POINT" and "LINE", which are consequently called undefined terms. Undefined terms, like the p and q of th pq- system, do get defined in a sense: implicitly-by the totality of all propos dons in which they occur, rather than explicitly, in a definition. One could maintain that a full definition of the undefined tern resides in the postulates alone, since the propositions which follow from them are implicit in the postulates already. This view would say that the postulates are implicit definitions of all the undefined terms, all of the undefined terms being defined in terms of the others. 







The Possibility of Multiple Interpretations A full formalization of geometry would take the drastic step of making every term undefined-that is, turning every term into a "meaningless" symbol of a formal system. I put quotes around "meaningless" because, as you know, the symbols automatically pick up passive meanings in accordance with the theorems they occur in. It is another question, though, whether people discover those meanings, for to do so requires finding a set of concepts which can be linked by an isomorphism to the symbols in the formal system. If one begins with the aim of formalizing geometry, presumably one has an intended interpretation for each symbol, so that the passive meanings are built into the system. That is what I did for p and q when I first created the pq-system. But there may be other passive meanings which are potentially perceptible, which no one has yet noticed. For instance, there were the surprise interpretations of p as "equals" and q as "taken from", in the original pq-system. Although this is rather a trivial example, it contains the essence of the idea that symbols may have many meaningful interpretations-it is up to the observer to look for them. We can summarize our observations so far in terms of the word "consistency". We began our discussion by manufacturing what appeared to be an inconsistent formal system-one which was internally inconsistent, as well as inconsistent with the external world. But a moment later we took it all back, when we realized our error: that we had chosen unfortunate interpretations for the symbols. By changing the interpretations, we regained consistency! It now becomes clear that consistency is not a property of a formal system per se, but depends on the interpretation which is proposed for it. By the same token, inconsistency is not an intrinsic property of any formal system. 

Varieties of Consistency We have been speaking of "consistency" and "inconsistency" all along, without defining them. We have just relied on good old everyday notions. But now let us say exactly what is meant by consistency of a formal system (together with an interpretation): that every theorem, when interpreted, becomes a true statement. And we will say that inconsistency occurs when there is at least one false statement among the interpreted theorems. This definition appears to be talking about inconsistency with the external world- what about internal inconsistencies? Presumably, a system would be internally inconsistent if it contained two or more theorems whose interpretations were incompatible with one another, and internally consistent if all interpreted theorems were compatible with one another. Consider, for example, a formal system which has only the following three theorems: TbZ, ZbE, and EbT. If T is interpreted as "the Tortoise", Z as "Zeno", E as "Egbert", and x by as "x beats y in chess always", then we have the following interpreted theorems: 







The Tortoise always beats Zeno at chess Zeno always beats Egbert at chess. Egbert always beats the Tortoise at chess. 

The statements are not incompatible, although they describe a rather bizarre circle of chess players. Hence, under this interpretation, the form; system in which those three strings are theorems is internally consistent although, in point of fact, none of the three statements is true! Intern< consistency does not require all theorems to come out true, but merely that they come out compatible with one another. Now suppose instead that x by is to be interpreted as "x was invented by y". Then we would have: 

The Tortoise was invented by Zeno. Zeno was invented by Egbert. Egbert was invented by the Tortoise. 

In this case, it doesn't matter whether the individual statements are true c false-and perhaps there is no way to know which ones are true, and which are not. What is nevertheless certain is that not all three can be true at one Thus, the interpretation makes the system internally inconsistent. The internal inconsistency depends not on the interpretations of the three capital letters, but only on that of b, and on the fact that the three capita are cyclically permuted around the occurrences of b. Thus, one can have internal inconsistency without having interpreted all of the symbols of the formal system. (In this case it sufficed to interpret a single symbol.) By tl time sufficiently many symbols have been given interpretations, it may t clear that there is no way that the rest of them can be interpreted so that a theorems will come out true. But it is not just a question of truth-it is question of possibility. All three theorems would come out false if the capitals were interpreted as the names of real people-but that is not why we would call the system internally inconsistent; our grounds for doing s would be the circularity, combined with the interpretation of the letter I (By the way, you'll find more on this "authorship triangle" in Chapter XX.; 

Hypothetical Worlds and Consistency We have given two ways of looking at consistency: the first says that system- plus-interpretation is consistent with the external world if every theorem comes out true when interpreted; the second says that a system-plus: interpretation is internally consistent if all theorems come out mutually compatible when interpreted. Now there is a close relationship between these two types of consistency. In order to determine whether several statements at mutually compatible, you try to imagine a world in which all of them could be simultaneously true. Therefore, internal consistency depends upon consistency with the external world-only now, "the external world" allowed to be any imaginable world, instead of the one we live in. But this is 







an extremely vague, unsatisfactory conclusion. What constitutes an “imaginable" world? After all, it is possible to imagine a world in which three characters invent each other cyclically. Or is it? Is it possible to imagine a world in which there are square circles? Is a world imaginable in which Newton's laws, and not relativity, hold? Is it possible to imagine a world in which something can be simultaneously green and not green? Or a world in which animals exist which are not made of cells? In which Bach improvised an eight-part fugue on a theme of King Frederick the Great? In which mosquitoes are more intelligent than people? In which tortoises can play football-or talk? A tortoise talking football would be an anomaly, of course. Some of these worlds seem more imaginable than others, since some seem to embody logical contradictions-for example, green and not green-while some of them seem, for want of a better word, "plausible" -- such as Bach improvising an eight-part fugue, or animals which are not made of cells. Or even, come to think of it, a world in which the laws of physics are different ... Roughly, then, it should be possible to establish different brands of consistency. For instance, the most lenient would be "logical consistency", putting no restraints on things at all, except those of logic. More specifically, a system-plus-interpretation would be logically consistent just as long as no two of its theorems, when interpreted as statements, directly contradict each other; and mathematically consistent just as long as interpreted theorems do not violate mathematics; and physically consistent just as long as all its interpreted theorems are compatible with physical law; then comes biological consistency, and so on. In a biologically consistent system, there could be a theorem whose interpretation is the statement "Shakespeare wrote an opera", but no theorem whose interpretation is the statement "Cell-less animals exist". Generally speaking, these fancier kinds of inconsistency are not studied, for the reason that they are very hard to disentangle from one another. What kind of inconsistency, for example, should one say is involved in the problem of the three characters who invent each other cyclically? Logical? Physical? Biological? Literary? Usually, the borderline between uninteresting and interesting is drawn between physical consistency and mathematical consistency. (Of course, it is the mathematicians and logicians who do the drawing-hardly an impartial crew . . .) This means that the kinds of inconsistency which "count", for formal systems, are just the logical and mathematical kinds. According to this convention, then, we haven't yet found an interpretation which makes the trio of theorems TbZ, ZbE, EbT inconsistent. We can do so by interpreting b as "is bigger than". What about T and Z and E? They can be interpreted as natural numbers-for example, Z as 0, T as 2, and E as 11. Notice that two theorems come out true this way, one false. If, instead, we had interpreted Z as 3, there would have been two falsehoods and only one truth. But either way, we'd have had inconsistency. In fact, the values assigned to T, Z, and E are irrelevant, as long as it is understood that they are restricted to natural numbers. Once again we see a case where only some of the interpretation is needed, in order to recognize internal inconsistency. 







Embedding of One Formal System In Another The preceding example, in which some symbols could have interpretations while others didn't, is reminiscent of doing geometry in natural languag4 using some words as undefined terms. In such a case, words are divide into two classes: those whose meaning is fixed and immutable, and, those whose meaning is to be adjusted until the system is consistent (these are th undefined terms). Doing geometry in this way requires that meanings have already been established for words in the first class, somewhere outside c geometry. Those words form a rigid skeleton, giving an underlying structure to the system; filling in that skeleton comes other material, which ca vary (Euclidean or non- Euclidean geometry). Formal systems are often built up in just this type of sequential, c hierarchical, manner. For example, Formal System I may be devised, wit rules and axioms that give certain intended passive meanings to its symbol Then Formal System I is incorporated fully into a larger system with more symbols-Formal System II. Since Formal System I's axioms and rules at part of Formal System II, the passive meanings of Formal System I symbols remain valid; they form an immutable skeleton which then plays large role in the determination of the passive meanings of the new symbols of Formal System II. The second system may in turn play the role of skeleton with respect to a third system, and so on. It is also possible-an geometry is a good example of this-to have a system (e.g., absolute geometry) which partly pins down the passive meanings of its undefined terms, and which can be supplemented by extra rules or axioms, which then further restrict the passive meanings of the undefined terms. This the case with Euclidean versus non- Euclidean geometry. 

Layers of Stability in Visual Perception In a similar, hierarchical way, we acquire new knowledge, new vocabulary or perceive unfamiliar objects. It is particularly interesting in the case understanding drawings by Escher, such as Relativity (Fig. 22), in which there occur blatantly impossible images. You might think that we won seek to reinterpret the picture over and over again until we came to interpretation of its parts which was free of contradictions- but we dot do that at all. We sit there amused and puzzled by staircases which go eve which way, and by people going in inconsistent directions on a sing staircase. Those staircases are "islands of certainty" upon which we base of interpretation of the overall picture. Having once identified them, we try extend our understanding, by seeking to establish the relationship which they bear to one another. At that stage, we encounter trouble. But if i attempted to backtrack-that is, to question the "islands of certainty"-s would also encounter trouble, of another sort. There's no way of backtracking and "undeciding" that they are staircases. They are not fishes, or whip or hands-they are just staircases. (There is, actually, one other on t-i leave all the lines of the picture totally uninterpreted, like the "meaningless 







FIGURE 22. Relativity, by M. C. Escher (lithograph, 1953). 

symbols" of a formal system. This ultimate escape route is an example of a "U-mode" response-a Zen attitude towards symbolism.) So we are forced, by the hierarchical nature of our perceptive processes, to see either a crazy world or just a bunch of pointless lines. A similar analysis could be made of dozens of Escher pictures, which rely heavily upon the recognition of certain basic forms, which are then put together in nonstandard ways; and by the time the observer sees the paradox on a high level, it is too late-he can't go back and change his mind about how to interpret the lower-level objects. The difference between an Escher drawing and non-Euclidean geometry is that in the latter, comprehensible interpretations can be found for the undefined terms, resulting in a com 







prehensible total system, whereas for the former, the end result is not reconcilable with one's conception of the world, no matter how long or stares at the pictures. Of course, one can still manufacture hypothetic worlds, in which Escherian events can happen ... but in such worlds, t1 laws of biology, physics, mathematics, or even logic will be violated on or level, while simultaneously being obeyed on another, which makes the: extremely weird worlds. (An example of this is in Waterfall (Fig. 5), whet normal gravitation applies to the moving water, but where the nature space violates the laws of physics.) 

Is Mathematics the Same in Every Conceivable World? We have stressed the fact, above, that internal consistency of a form; system (together with an interpretation) requires that there be some imaginable world-that is, a world whose only restriction is that in it, mathematics and logic should be the same as in our world-in which all the interpreted theorems come out true. External consistency, however consistency with the external world-requires that all theorems come of true in the real world. Now in the special case where one wishes to create consistent formal system whose theorems are to be interpreted as statements of mathematics, it would seem that the difference between the two types of consistency should fade away, since, according to what we sat above, all imaginable worlds have the same mathematics as the real world. Thus, i every conceivable world, 1 plus 1 would have to be 2; likewise, there would have to be infinitely many prime numbers; furthermore, in every conceivable world, all right angles would have to be congruent; and of cours4 through any point not on a given line there would have to be exactly on parallel line ... But wait a minute! That's the parallel postulate-and to assert i universality would be a mistake, in light of what's just been said. If in all conceivable worlds the parallel postulate-is obeyed, then we are asserting that non-Euclidean geometry is inconceivable, which puts us back in the same mental state as Saccheri and Lambert-surely an unwise move. But what, then, if not all of mathematics, must all conceivable worlds share? Could it I as little as logic itself? Or is even logic suspect? Could there be worlds where contradictions are normal parts of existence-worlds where contradictious are not contradictions? Well, in some sense, by merely inventing the concept, we have shoe that such worlds are indeed conceivable; but in a deeper sense, they are al: quite inconceivable. (This in itself is a little contradiction.) Quite serious] however, it seems that if we want to be able to communicate at all, we ha, to adopt some common base, and it pretty well has to include logic. (The are belief systems which reject this point of view-it is too logical. particular, Zen embraces contradictions and non-contradictions with equ eagerness. This may seem inconsistent, but then being inconsistent is pa of Zen, and so ... what can one say?) 







Is Number Theory the Same In All Conceivable Worlds? If we assume that logic is part of every conceivable world (and note that we have not defined logic, but we will in Chapters to come), is that all? Is it really conceivable that, in some worlds, there are not infinitely many primes? Would it not seem necessary that numbers should obey the same laws in all conceivable worlds? Or ... is the concept "natural number" better thought of as an undefined term, like "POINT" or "LINE"? In that case, number theory would be a bifurcated theory, like geometry: there would be standard and nonstandard number theories. But there would have to be some counterpart to absolute geometry: a "core" theory, an invariant ingredient of all number theories which identified them as number theories rather than, say, theories about cocoa or rubber or bananas. It seems to be the consensus of most modern mathematicians and philosophers that there is such a core number theory, which ought to be included, along with logic, in what we consider to be "conceivable worlds". This core of number theory, the counterpart to absolute geometry-is called Peano arithmetic, and we shall formalize it in Chapter VIII. Also, it is now well established-as a matter of fact as a direct consequence of Gödel’s Theorem-that number theory is a bifurcated theory, with standard and nonstandard versions. Unlike the situation in geometry, however, the number of "brands" of number theory is infinite, which makes the situation of number theory considerably more complex. For practical purposes, all number theories are the same. In other words, if bridge building depended on number theory (which in a sense it does), the fact that there are different number theories would not matter, since in the aspects relevant to the real world, all number theories overlap. The same cannot be said of different geometries; for example, the sum of the angles in a triangle is 180 degrees only in Euclidean geometry; it is greater in elliptic geometry, less in hyperbolic. There is a story that Gauss once attempted to measure the sum of the angles in a large triangle defined by three mountain peaks, in order to determine, once and for all, which kind of geometry really rules our universe. It was a hundred years later that Einstein gave a theory (general relativity) which said that the geometry of the universe is determined by its content of matter, so that no one geometry is intrinsic to space itself. Thus to the question, "Which geometry is true?" nature gives an ambiguous answer not only in mathematics, but also in physics. As for the corresponding question, "Which number theory is true?", we shall have more to say on it after going through Gödel’s Theorem in detail. 

Completenes If consistency is the minimal condition under which symbols acquire passive meanings, then its complementary notion, completeness, is the maximal confirmation of those passive meanings. Where consistency is the property 







way round: "Every true statement is produced by the system". Now I refine the notion slightly. We can't mean every true statement in th world-we mean only those which belong to the domain which we at attempting to represent in the system. Therefore, completeness mean! "Every true statement which can be expressed in the notation of the system is a theorem." 

Consistency: when every theorem, upon interpretation, comes out true (in some imaginable world). Completeness: when all statements which are true (in some imaginable world), and which can be expressed as well-formed strings of the system, are theorems. 

An example of a formal system which is complete on its own mode level is the original pq-system, with the original interpretation. All true additions of two positive integers are represented by theorems of th system. We might say this another way: "All true additions of two positive integers are provable within the system." (Warning: When we start using th term "provable statements" instead of "theorems", it shows that we at beginning to blur the distinction between formal systems and their interpretations. This is all right, provided we are very conscious of th blurring that is taking place, and provided that we remember that multiple interpretations are sometimes possible.) The pq-system with the origin interpretation is complete; it is also consistent, since no false statement is-, use our new phrase-provable within the system. Someone might argue that the system is incomplete, on the grounds that additions of three positive integers (such as 2 + 3 + 4 =9) are not represented by theorems of the pq-system, despite being translatable into the notation of the system (e.g., --p---p----q---- --------). However, this string is not well-formed, and hence should be considered to I just as devoid of meaning as is p q p---q p q. Triple additions are simply not expressible in the notation of the system-so the completeness of the system is preserved. Despite the completeness of the pq-system under this interpretation, certainly falls far short of capturing the full notion of truth in numb theory. For example, there is no way that the pq-system tells us how mat prime numbers there are. Gödel’s Incompleteness Theorem says that any system which is "sufficiently powerful" is, by virtue of its power, incomplete, in the sense that there are well-formed strings which express tr statements of number theory, but which are not theorems. (There a truths belonging to number theory which are not provable within the system.) Systems like the pq-system, which are complete but not very powerful, are more like low-fidelity phonographs; they are so poor to beg with that it is obvious that they cannot do what we would wish them do-namely tell us everything about number theory. 







How an Interpretation May Make or Break Completeness What does it mean to say, as I did above, that "completeness is the maximal confirmation of passive meanings"? It means that if a system is consistent but incomplete, there is a mismatch between the symbols and their interpretations. The system does not have the power to justify being interpreted that way. Sometimes, if the interpretations are "trimmed" a little, the system can become complete. To illustrate this idea, let's look at the modified pq-system (including Axiom Schema II) and the interpretation we used for it. After modifying the pq-system, we modified the interpretation for q from "equals" to "is greater than or equal to". We saw that the modified pq-system was consistent under this interpretation; yet something about the new interpretation is not very satisfying. The problem is simple: there are now many expressible truths which are not theorems. For instance, "2 plus 3 is greater than or equal to 1" is expressed by the nontheorem --p---q-. The interpretation is just too sloppy! It doesn't accurately reflect what the theorems in the system do. Under this sloppy interpretation, the pq-system is not complete. We could repair the situation either by (1) adding new rules to the system, making it more powerful, or by (2) tightening up the interpretation. In this case, the sensible alternative seems to be to tighten the interpretation. Instead of interpreting q as "is greater than or equal to", we should say "equals or exceeds by 1". Now the modified pq-system becomes both consistent and complete. And the completeness confirms the appropriateness of the interpretation. 

Incompleteness of Formalized Number Theory In number theory, we will encounter incompleteness again; but there, to remedy the situation, we will be pulled in the other direction-towards adding new rules, to make the system more powerful. The irony is that we think, each time we add a new rule, that we surely have made the system complete now! The nature of the dilemma can be illustrated' by the following allegory ... We have a record player, and we also have a record tentatively labeled "Canon on B-A-C-H". However, when we play the record on the record player, the feedback- induced vibrations (as caused by the Tortoise's records) interfere so much that we do not even recognize the tune. We conclude that something is defective-either our record, or our record player. In order to test our record, we would have to play it on friends' record players, and listen to its quality. In order to test our phonograph, we would have to play friends' records on it, and see if the music we hear agrees with the labels. If our record player passes its test, then we will say the record was defective; contrariwise, if the record passes its test, then we will say our record player was defective. What, however, can we conclude when we find out that both pass their respective tests? That is the moment to remember the chain of two isomorphisms (Fig. 20), and think carefully! 







Little Harmonic Labyrinth 

The Tortoise and Achilles are spending a day at Coney Island After buying a couple of cotton candies, they decide to take a ride on the Ferris wheel. 


\tortoise This is my favorite ride. One seems to move so far, and\par
reality one gets nowhere. 
\achilles I can see why it would appeal to you. Are you all strapped in?\par

\tortoise Yes, I think I've got this buckle done. Well, here we go.\par

\achilles You certainly are exuberant today.\par

\tortoise I have good reason to be. My aunt, who is a fortune-teller me that a stroke of\par
Good Fortune would befall me today. So I am tingling with anticipation. 
\achilles Don't tell me you believe in fortune-telling!\par

\tortoise No ... but they say it works even if you don't believe ii\par

\achilles Well, that's fortunate indeed.\par

\tortoise Ah, what a view of the beach, the crowd, the ocean, the city. . .\par

\achilles Yes, it certainly is splendid. Say, look at that helicopter there. It seems to be\par
flying our way. In fact it's almost directly above us now. 
\tortoise Strange-there's a cable dangling down from it, which is very close to us. It's\par
coming so close we could practically grab it 
\achilles Look! At the end of the line there's a giant hook, with a note\par



\intro (He reaches out and snatches the note. They pass by and are on their z down.) \endintro



\tortoise Can you make out what the note says?\par

\achilles Yes-it reads, "Howdy, friends. Grab a hold of the hook time around, for an\par
Unexpected Surprise." 
\tortoise The note's a little corny but who knows where it might lead, Perhaps it's got\par
something to do with that bit of Good Fortune due me. By all means, let's try it! 
\achilles Let's!\par


(On the trip up they unbuckle their buckles, and at the crest of the ride, grab for the giant hook. All of a sudden they are whooshed up by the ca which quickly reels them skyward into the hovering helicopter. A It strong hand helps them in.) 

Voice: Welcome aboard-Suckers. 
\achilles Wh-who are you?\par








Voce: Allow me to introduce myself. I am Hexachlorophene J. Goodforttune, Kidnapper At-Large, and Devourer of Tortoises par Excellence, at your service. 
\tortoise Gulp!\par
Achilles (whispering to his friend): Uh-oh-I think that this "Goodfortune" is not exactly what we'd anticipated. (To Goodfortune) Ah-if I may be so bold-where are you spiriting us off to? Goodfortune: Ho ho! To my all-electric kitchen-in-the-sky, where I will prepare THIS tasty morsel-(leering at the Tortoise as he says this)-in a delicious pie-in-the-sky! And make no mistake-it's all just for my gobbling pleasure! Ho ho ho! 
\achilles All I can say is you've got a pretty fiendish laugh.\par
Goodfortune (laughing fiendishly): Ho ho ho! For that remark, my friend, you will pay dearly. Ho ho! 
\achilles Good grief-I wonder what he means by that!\par
Goodfortune: Very simple-I've got a Sinister Fate in store for both of you! Just you wait! Ho ho ho! Ho ho ho! 
\achilles Yikes!\par
Goodfortune: Well, we have arrived. Disembark, my friends, into my fabulous all-electric kitchen-in-the-sky. 


\intro (They walk inside.) \endintro


Let me show you around, before I prepare your fates. Here is my bedroom. Here is my study. Please wait here for me for a moment. I've got to go sharpen my knives. While you're waiting, help yourselves to some popcorn. Ho ho ho! Tortoise pie! Tortoise pie! My favorite kind of pie! (Exit.) 


\achilles Oh, boy-popcorn! I'm going to munch my head off!\par

\tortoise Achilles! You just stuffed yourself with cotton candy! Besides, how can you\par
think about food at a time like this? 
\achilles Good gravy-oh, pardon me-I shouldn't use that turn of phrase, should I? I mean\par
in these dire circumstances ... Tortoise: I'm afraid our goose is cooked. 
\achilles Say-take a gander at all these books old Goodfortune has in his study. Quite a\par
collection of esoterica: Birdbrains I Have Known; Chess and Umbrella-Twirling Made Easy; Concerto for Tapdancer and Orchestra ... Hmmm. 
\tortoise What's that small volume lying open over there on the desk, next to the\par
dodecahedron and the open drawing pad? 
\achilles This one? Why, its title is Provocative Adventures of Achilles and the Tortoise\par
Taking Place in Sundry Spots of the Globe. Tortoise: A moderately provocative title. 
\achilles Indeed-and the adventure it's opened to looks provocative. It's called "Djinn and\par
Tonic". 
\tortoise Hmm ... I wonder why. Shall we try reading it? I could take the Tortoise's part,\par
and you could take that of Achilles. 








\achilles I’m game. Here goes nothing . . .\par



\intro (They begin reading "Djinn and Tonic".) \endintro


(Achilles has invited the Tortoise over to see his collection of prints by his favorite artist, M. C. Escher.) 


\tortoise These are wonderful prints, Achilles.\par

\achilles I knew you would enjoy seeing them. Do you have any particular\par
favorite? 
\tortoise One of my favorites is Convex and Concave, where two internally\par
consistent worlds, when juxtaposed, make a completely inconsistent composite world. Inconsistent worlds are always fun places to visit, but I wouldn't want to live there. 
\achilles What do you mean, "fun to visit"? Inconsistent worlds don't EXIST,\par
so how can you visit one? 
\tortoise I beg your pardon, but weren't we just agreeing that in\par


this Escher picture, an inconsistent world is portrayed? 
\achilles Yes, but that's just a two-dimensional world-a fictitious world-a\par
picture. You can't visit that world. 
\tortoise I have my ways ...\par

\achilles How could you propel yourself into a flat picture-universe?\par

\tortoise By drinking a little glass of PUSHING-POTION. That does the\par
trick. 
\achilles What on earth is pushing-potion?\par

\tortoise It's a liquid that comes in small ceramic phials, and which, when\par
drunk by someone looking at a picture, "pushes'' him right into the world of that picture. People who aren't aware of the powers of pushing-potion often are pretty surprised by the situations they wind up in. 
\achilles Is there no antidote? Once pushed, is one irretrievably lost?\par

\tortoise In certain cases, that's not so bad a fate. But there is, in fact, another\par
potion-well, not a potion, actually, but an elixir-no, not an elixir, but a-a 
\tortoise He probably means "tonic".\par

\achilles Tonic?\par
\bye
